\documentclass[11pt]{article}
\usepackage[margin=1in]{geometry}
\usepackage{amsmath,amssymb,amsthm,mathtools}
\usepackage{array,booktabs,longtable}
\usepackage[hidelinks]{hyperref}
\usepackage{cleveref}

\usepackage{xcolor}

\newtheorem{theorem}{Theorem}
\newtheorem{lemma}{Lemma}[section]
\newtheorem{definition}{Definition}[section]
\newtheorem{importedresult}{Imported result}[section]
\newtheorem{corollary}[lemma]{Corollary}

\usepackage{enumitem}
\usepackage{etoolbox}
\newcommand{\thestmt}{}
\AtBeginEnvironment{theorem}{\renewcommand{\thestmt}{\thetheorem}}
\AtBeginEnvironment{lemma}{\renewcommand{\thestmt}{\thelemma}}
\AtBeginEnvironment{definition}{\renewcommand{\thestmt}{\thedefinition}}
\AtBeginEnvironment{importedresult}{\renewcommand{\thestmt}{\theimportedresult}}
\AtBeginEnvironment{corollary}{\renewcommand{\thestmt}{\thecorollary}}
\setlist[enumerate,1]{ref=\thestmt(\arabic*)}
\crefname{enumi}{Item}{Items}
\Crefname{enumi}{Item}{Items}

\newcommand{\R}{\mathbb R}
\newcommand{\D}{\mathbb D}
\newcommand{\Trunc}{\operatorname{Trunc}}
\newcommand{\diag}{\operatorname{diag}}
\newcommand{\nnz}{\operatorname{nnz}}
\newcommand{\clip}{\operatorname{clip}}
\newcommand{\eps}{\varepsilon}
\newcommand{\MM}{\textsc{MatMult}}

\renewcommand{\O}{\widetilde O}

\newcommand{\xini}{x^{(0)}}
\newcommand{\what}{\widehat{w}}
\newcommand{\xhat}{\widehat{x}}
\newcommand{\shat}{\widehat{s}}
\newcommand{\zhat}{\widehat{z}}
\newcommand{\muhat}{\widehat{\mu}}
\newcommand{\tautil}{\widetilde{\tau}}
\newcommand{\ghat}{\widehat{g}}
\newcommand{\sgn}{\operatorname{sgn}}

\newcommand{\defeq}{\stackrel{\Delta}{=}}
\newcommand{\Cnorm}{C_{\mathrm{norm}}}
\newcommand{\Cvalid}{C_{\mathrm{valid}}}
\newcommand{\E}{\mathbb E}

\usepackage{float}
\usepackage{fancybox}
\newenvironment{fminipage}%
  {\begin{Sbox}\begin{minipage}}%
  {\end{minipage}\end{Sbox}\fbox{\TheSbox}}

\newenvironment{algbox}[0]{\vskip 0.2in
\noindent 
\begin{fminipage}{6.3in}
}{
\end{fminipage}
\vskip 0.2in
}

\title{Bit Complexity of LS`15 under Fixed Point}
\author{}
\date{}
\begin{document}
\maketitle

\begin{abstract}
We show that the linear programming runtime of $\O(d^{1/2} ( \nnz(A) + d^2 ))$
arithmetic operations by [Lee-Sidford FOCS`15] holds under bit complexity: on
inputs with $B$ bits the method takes
$\O(d^{1/2} ( \nnz(A) + d^2 ) B^{3})$ bit operations, where
$\O$ hides factors polylogarithmic in
$n$.
\end{abstract}

This document was generated by Aidan Bai, Chenxin Dai, and Richard Peng,
in order to rigorously support the claims made in a another paper.
It took three weeks of systhesis, outline proof reading, and script
modifications using a mixture of Kimi K3, GPT 6 Codex, and the Devin Harness.
We thank Moonshot AI and Cognition AI for providing the compute resources
used throughout this process.

\section{Introduction}
\label{sec:intro}

Let $1\leq d \leq n \leq d^{10}$, let
$A\in\R^{n \times d}$ have full column rank and no zero row, and consider the
linear program
\begin{align}
\min_{x\in\R^n} \quad & c^\top x \tag{primal} \label{eq:primal}\\
\text{subject to} \quad & A^\top x = b, \nonumber\\
& l\le x\le u, \nonumber
\end{align}
whose dual is
\begin{align}
\max_{y\in\R^d,\ sl,su\in\R^n} \quad & b^\top y + l^\top sl - u^\top su \tag{dual} \label{eq:dual}\\
\text{subject to} \quad & Ay + sl - su = c, \nonumber\\
& sl, su \geq 0 . \nonumber
\end{align}
We refer to these two programs as the primal/dual pair.

We work in the fixed point model of \cite{GPV23}: the input
is given, and every number the algorithm forms is kept, as a signed binary
fraction with $B$ bits before and after the decimal point, each arithmetic
operation on such numbers costs $\O(B)$ bit operations, and every intermediate
result is truncated back to this precision.
The model is made precise in Section~\ref{subsec:roundoff}, and all costs below
are bit complexities in it.
Throughout, $\O(\cdot)$ hides factors polynomial in
$\log n$ only; every power of
$B$ is written out.
Since $B \le d^{1/2} \le n$, factors polynomial in
$\log B$, such as the
$\log B$ of fast integer multiplication, are hidden as well.
We assume throughout that $B \le d^{1/2}$.

\begin{theorem}
\label{thm:main}
Given a primal/dual pair $A,b,c,l,u$ and a point
$\xini$ with
$B \le d^{1/2}$ bits before/after decimal place in fixed point, such that
$\|A^{\top} A\|, \|(A^{\top}A)^{-1}\| \leq 2^{B}$, and
$l + 2^{-B} < \xini < u - 2^{-B}$ and
$A^\top \xini = b$, we can compute
$x\in[l,u]$ satisfying
$\|A^\top x-b\|_2\leq 2^{-B}$ and
$c^\top x \le \mathrm{opt} + 2^{-B}$, where
$\mathrm{opt}$ is the optimum of the primal, in bit complexity
$\O(d^{1/2} ( \nnz(A) + d^2 ) B^{3})$ with high probability.
\end{theorem}

The objective guarantee is one sided, as in
\cite{LS15,BLLSSSW21}: the
returned $x$ is only
$2^{-B}$-feasible, and a lower bound on
$c^\top x$ in terms of
$\mathrm{opt}$ would need the sensitivity of the optimum to
$b$, which is not controlled here.
The point $\xini$ is used only to know that the primal is feasible.

\paragraph{Why the bound does not follow from LS15.}
Theorem~\ref{thm:main} is the fixed point counterpart of the running time
$\O(d^{1/2}(\nnz(A) + d^2) \log(1/\epsilon))$ of \cite{LS15}, and
it is not obtained from that paper by inspection.
The count there is of arithmetic operations on real numbers: the inverse
maintenance keeps a $d \times d$ inverse of a sampled matrix, corrects it by low
rank updates, and uses it as a preconditioner for an $\epsilon$-approximate
solve, and none of the growth of the entries of the maintained matrices, the
compounding of truncation errors over the $\O(d^{1/2})$ updates, or the loss of
the spectral guarantee of a row sample chosen from truncated weights, is
accounted for.
Three further points have to be settled before the arithmetic count converts.
The path following analysis is stated for exact Newton directions computed from
approximate weights, so a direction that is itself the truncated output of an
approximate solver needs its own perturbation argument
(Section~\ref{sec:potential_and_progress}); the sampled inverse of
\cite{LS15} is valid only against weights that do not depend on
its random choices, and that paper hides the sample from the path following
method by adding Gaussian noise to a real-valued solve, an argument that has to
be redone when the noise is discrete, the solve is truncated, and everything the
caller derives from the output is itself truncated
(Section~\ref{subsec:randomness}); and the magnitudes of the weights, the
inverses and the iterates along the path have to be bounded by $2^{O(B)}$ in
terms of the input so that the fixed point numbers do not overflow
(Section~\ref{sec:overall}).
The path following analysis that \cite{LS15} invokes is that of
\cite{LS14}, and this analysis has itself been revised substantially
since: its arXiv version \cite{LS13:arxiv} went from 48 pages in
December 2013 to 60 pages over three versions ending in March 2015, and was then
merged with its companion paper into \cite{LS19:arxiv}, 65 to 69 pages
over two versions in 2019 and 2020, which by its own description contains new
results beyond the earlier submissions, among them a nearly optimal
self-concordant barrier and its relation to Lewis weights.
The path following statements imported below are therefore taken from the robust
formulation of \cite{BLLSSSW21} rather than
from \cite{LS14}, and the constants of those statements that we could
not rederive from their proofs are listed in the issues above.

\paragraph{The same bound via the robust interior point method.}
A second route to the bound of Theorem~\ref{thm:main} goes through the robust
interior point method of \cite{BLLSSSW21},
which is stated against three abstract data structures, a heavy hitter, a heavy
sampler and an inverse maintenance, each with an initialization cost $P$, a cost
$c_i \ge \nnz(a_i)$ for changing the scaling of row
$i$, and a query cost
$Q$.
Theorem 6.1 and Lemma 6.7 there bound the total time of path following, in our
dimensions, by
\[
\O\left( (P \|c\|_1)^{1/2}
  + d^{1/2} \left( Q + d \max_i \nnz(a_i) \right) \right)
\]
per factor of $\log(\mu_0 / \mu_f)$.
For a general linear program that paper uses the sketching based heavy hitter of
\cite{BLSS20} with $c_i = \O(d)$ for every row and a heavy
sampler with $Q = \O(d^{2} + n d^{1/2})$ (Lemmas B.1 and B.2 and Corollary B.6
there), and it is these two, $c_i = \O(d)$ in place of
$\nnz(a_i)$ and
$n d^{1/2}$ in place of
$\nnz(A)$, that turn the bound into
$\O(nd + d^{2.5})$ in its Lemma 8.2.
The brute force alternative computes the vector $A h$ explicitly at every query
and reads the heavy entries and the sampling probabilities off it: this is a
heavy hitter and a heavy sampler with $P = \nnz(A)$,
$c_i = \nnz(a_i)$ and
$Q = \nnz(A)$, while the inverse maintenance of Lemma B.2 there has
$P = \nnz(A) + \MM(d, d, d)$,
$c_i = O(1)$ and
$Q = \O(d^{2})$.
Substituting, with $\max_i \nnz(a_i) \le d$ and
$(\nnz(A) \MM(d, d, d))^{1/2} \le \nnz(A) d^{1/2} + \MM(d, d, d) d^{-1/2}$,
gives $\O(d^{1/2}(\nnz(A) + d^{2}))$, formally.
We have not carried this out.
The substitution itself needs only that the interfaces (Definitions 3.1, 6.2 and
6.3 there) accept a query cost $Q$ as large as
$\nnz(A)$, which they do since
$Q$ enters the accounting additively per step; the work is that the whole of
\cite{BLLSSSW21} is in exact arithmetic, so
the magnitude, truncation and randomness accounting of this paper would have to
be repeated for its primal, gradient, dual and Lewis weight maintenance
structures, instead of for the single \textsc{Solve} of
Section~\ref{sec:dynamic}.

\paragraph{Difficulties in reducing the $d^{2.5}$ term.}
The term $d^{1/2} \cdot d^{2}$ is
$d^{1/2}$ path steps times the cost of applying a dense
$d \times d$ maintained inverse to a vector inside \textsc{Solve}, and row
sampling does not lower it: the sample only brings the number of rows down to
$\O(d)$.
In the square dense case the corresponding $n^{2.5}$ has been beaten by
combining a robust central path, which accepts directions accurate only to a
constant in an $\ell_\infty$-type norm, with implicit inverse maintenance, in
which both the inverse and the vector it is applied to are updated lazily in
batches
\cite{CLS21:journal,B20,JSWZ21}.
Those running times are stated in terms of the exponent of square matrix
multiplication and its rectangular and dual variants, and the exponent has moved
over the past five years, from $2.37286$ \cite{AVW21} to
$2.371866$ \cite{DWZ23},
$2.371552$ \cite{VXXZ24},
$2.371339$ \cite{ADVXXZ25} and
$2.371177$
\cite{DEKMRSZAVB26:arxiv},
the last three also improving the rectangular exponents.
For tall sparse $A$, however, the exponent is not what decides whether
$d^{2.5}$ can be lowered.
The inverse side, Woodbury layers with periodic restarts and batched folds,
already costs less than $d^{2}$ per path step under the assumption
$\MM(d, d, d) = \O(d^{1 + \sqrt{2}})$ of this paper, and any of the bounds above
only lowers that; the robust path itself accepts a fixed constant-accurate
inverse with a perturbation argument and no change of exponent.
What blocks the improvement is that the inverse has to be applied to a fresh
Newton right hand side at every step, a vector that changes in all $d$
coordinates.
The square case avoids this by maintaining the right hand side as a slowly
changing vector, and none of the known ways of doing so extends from the square,
or dense, case to tall sparse $A$ without further assumptions on the sparsity of
the rows.

\paragraph{Organization.}
\Cref{sec:prelims} (\nameref{sec:prelims}) fixes the fixed point model, the
notation, and the imported facts on matrix inversion with truncation and on
leverage scores.
\Cref{sec:dynamic} (\nameref{sec:dynamic}) gives the data structure behind
\textsc{Initialize}, \textsc{Update} and \textsc{Solve}: exact inverse
maintenance and its dependence on matrix multiplication, the row sampled
approximate inverse, the hiding of its randomness from an adaptive caller, and
the accounting of the solver against leverage stable weight sequences.
\Cref{sec:lewis} (\nameref{sec:lewis}) collects everything about regularized
Lewis weights, their estimation from a solver, the barrier and its parameters,
and the centrality potential.
\Cref{sec:potential_and_progress} (\nameref{sec:potential_and_progress})
analyzes one step of the method in fixed point: the step itself, direction
finding without a bisection, the effect of solve and truncation errors on
centrality and feasibility, the progress per step, and the stability of the
weights handed to the data structure.
\Cref{sec:overall} (\nameref{sec:overall}) bounds the magnitudes and bit lengths
along the path, verifies the stability hypothesis of the data structure,
constructs the initial point, and assembles the proof of Theorem~\ref{thm:main}.

\section{Prelims}
\label{sec:prelims}

For positive vectors, all powers, quotients, logarithms, and comparisons are
componentwise.
Logarithms are natural, so that $e^{-c \log n} = n^{-c}$.

For vectors $x$ and
$y$ and some
$\epsilon > 0$, we use
$x \approx_\epsilon y$ to denote
$e^{-\epsilon} x_i \le y_i \le e^\epsilon x_i$ for every
$i$.
For symmetric positive definite matrices $X$ and
$Y$ of the same size, we use
$X \approx_\epsilon Y$ to denote
$e^{-\epsilon} Y \preceq X \preceq e^{\epsilon} Y$.
The subscript of $\approx$ is a generic accuracy and takes whatever value is
written there; the unsubscripted $\epsilon$ of Sections~\ref{sec:lewis}
to~\ref{sec:overall} is the fixed centrality parameter of
Equation~\eqref{eq:parameters}.

The matrix $A$ is assumed to have no zero rows, so that
$n \le \nnz(A)$: a zero row constrains nothing and can be dropped in a single
pass.

\begin{definition}
We say invertible matrix $X$ has magnitude
$\rho$ if
$\|X\|, \|X^{-1}\| < \rho$, and
$\rho$ is more than both dimension of
$X$.
\end{definition}

$\MM(n_1, n_2, n_3)$ denotes the arithmetic cost of multiplying a
$n_1 \times n_2$ matrix by an
$n_2 \times n_3$ matrix.
We state all dense costs through $\MM$.
Dense products arise only in the exact inverse maintenance of
Section~\ref{subsec:exact_maintain}, in the static estimate that initializes it,
and in the restarts built on it, so the assumption under which they are
simplified is stated there rather than here.

\subsection{Arithmetic under Roundoffs}
\label{subsec:roundoff}

We work with fixed point numbers.
A \emph{number with $\ell$ bits} is a signed binary number with at most
$\ell$ bits before and
$\ell$ bits after the decimal point, that is, an integer multiple of
$2^{-\ell}$ of absolute value less than
$2^{\ell}$.
For a real $x$ with at most
$\ell$ bits before the decimal point,
$|x| < 2^{\ell}$, we write
$\Trunc(x, \ell)$ for
$x$ truncated toward zero to
$\ell$ bits after the decimal point; it is a number with
$\ell$ bits, has the sign of
$x$, and satisfies
\[
|\Trunc(x, \ell)| \le |x|,
\qquad
|x - \Trunc(x, \ell)| < 2^{-\ell} .
\]
$\Trunc(x, \ell)$ is only ever applied to
$x$ with at most
$\ell$ bits before the decimal point; wherever it appears below, the argument is
bounded in magnitude by a power of $\rho = 2^{O(B)}$ that is far below
$2^{\ell}$, and we do not repeat this check.
Applied to a vector or a matrix, $\Trunc$ acts entrywise.
We say a matrix has entries with $\ell$ bits if all its entries are numbers with
$\ell$ bits.

This is the fixed point model of \cite{GPV23}: an
arithmetic operation on numbers with $\ell$ bits, one of
$+$,
$-$,
$\times$,
$\div$ and the square root, returns the exact result truncated by
$\Trunc(\cdot, \ell)$, and comparisons are exact.

\begin{lemma}
\label{lem:arith}
Arithmetic in numbers with $\ell$ bits can be done at cost
$\O(\ell)$.
\end{lemma}

Two operations lose nothing when given room: the sum or difference of two
numbers with $\ell$ bits is a number with
$\ell + 1$ bits, and their product is a number with
$2 \ell$ bits, so both are exact when carried out on numbers with
$2 \ell$ bits.
This is how the comparisons of ratios in Figure~\ref{fig:direction} are made
exact.

The error of a longer computation is bounded by counting operations and the
number of multiplicative operations on a path.
We say a quantity is computed by a formula with $m$ operations if it is obtained
from numbers with $\ell$ bits by
$m$ arithmetic operations in which every intermediate result is used once, a
result used $k$ times being counted
$k$ times; the multiplicative depth of the formula is the largest number of
operations among $\times$,
$\div$ and square root on a path from an input to the output.

\begin{lemma}
\label{lem:roundoff}
Let a quantity be computed by a formula with $m$ operations and multiplicative
depth $q$ on numbers with
$\ell$ bits, whose inputs are within
$e_1, \ldots, e_k$ of the exact input values, each use of an input counted
separately, and let $M \ge 2$ be such that, in exact arithmetic, every input and
every intermediate result has absolute value at most $M$, every divisor and
every radicand has absolute value at least $1 / M$, and
$M \le 2^{\ell - 1}$.
Write
\[
E
\defeq
\left( m 2^{-\ell} + \sum_{j = 1}^{k} e_j \right)
\left( 2 M^{2} \right)^{q} .
\]
If $E \le 1 / (2 M)$, then the computed value differs from the exact one by at
most $E$.
\end{lemma}

The proof is in Appendix~\ref{app:roundoff}.
Every formula of Sections~\ref{sec:potential_and_progress} and~\ref{sec:overall}
is evaluated on numbers with $40 \log \rho$ bits at
$M = \rho$, and Equation~\eqref{eq:step_roundoff} there is the resulting
standing bound of $\rho^{-20}$ on its error; each proof that uses it records the
operation count and the depth of the formula in question.

\subsection{Matrix Inversion with Fixed Point Roundoffs}
\label{subsec:inverse}

The first imported result, Lemma 3.1 of \cite{GPV23}, is
that a small perturbation of a matrix perturbs its inverse in proportion to the
square of its magnitude.

\begin{importedresult}
\label{imp:perturb}
For invertible matrices $X$ and
$Y$ such that
$Y$ has magnitude
$\rho_Y$, and
$\|X-Y\|_F\le \epsilon$ for some
$0 < \epsilon < 1/(2 \rho_Y)$, we have:
\begin{enumerate}
\item \label{itm:imp_perturb:1} $X$ has magnitude at most
  $2 \rho_Y$.
\item \label{itm:imp_perturb:2} $\|X^{-1}-Y^{-1}\|_F \le 2 \rho_Y^2 \epsilon$.
\end{enumerate}
\end{importedresult}


The second, Theorem 3.2 of \cite{GPV23}, is an approximate
dense inverse in fixed point.

\begin{importedresult}
\label{imp:denseinverse}
For $d \times d$ matrix
$X$ with magnitude
$\rho$, and error
$0 < \epsilon < 0.1$, we can compute explicit matrix
$Z$ with such that
\[
\|Z-X^{-1}\|_F \le \epsilon.
\]
in bit complexity $\O(\MM(d,d,d) \log(\rho/\epsilon))$.
\end{importedresult}

\subsection{Leverage scores}
\label{subsec:leverage}

The one weight vector the data structure of Section~\ref{sec:dynamic} works with
is the classical one.
The leverage score of a row is its diagonal entry in the orthogonal projection
onto the column space, so it lies in $[0, 1]$, the scores sum to the rank
$d$, and a row with a small score is one whose removal barely changes
$A^\top A$.
Sampling rows with probabilities proportional to leverage scores gives a
spectral approximation of $A^\top A$ with
$\O(d)$ rows, which is how Section~\ref{sec:dynamic} keeps its sample small, and
the same scores measure how much a change of weights moves the sample, which is
the stability notion of Definition~\ref{def:tau_stable}.
The regularized Lewis weights of Section~\ref{sec:lewis} are defined through the
leverage scores of a reweighting of $A$ by a power of the weights themselves,
and Section~\ref{sec:potential_and_progress} passes between the two through
Lemma~\ref{lem:sigma_vs_tau}.

\begin{definition}
For a full-column-rank $A\in\R^{n \times d}$, write
$\tau(A)$ for the diagonal of the orthogonal projection
$A(A^\top A)^{-1}A^\top$, and
$\tau_i(A)$ for its
$i$th entry.
\end{definition}

The concentration inequality behind leverage score sampling is the matrix
Chernoff inequality, Theorem 1.1 of \cite{T12:journal}, stated with its
constants.

\begin{importedresult}
\label{imp:matrix_chernoff}
Let $X_1, \ldots, X_N$ be independent random symmetric positive semidefinite
$d \times d$ matrices with
$\lambda_{\max}(X_j) \le R$ almost surely, and let
$\mu_{\min} \defeq \lambda_{\min}(\sum_j \E X_j)$ and
$\mu_{\max} \defeq \lambda_{\max}(\sum_j \E X_j)$.
Then for every $\delta \in [0, 1)$,
\[
\Pr\left[ \lambda_{\min}\left( \sum_{j = 1}^{N} X_j \right)
\le (1 - \delta) \mu_{\min} \right]
\le
d \left( \frac{e^{-\delta}}{(1 - \delta)^{1 - \delta}}
\right)^{\mu_{\min} / R} ,
\]
and for every $\delta \ge 0$,
\[
\Pr\left[ \lambda_{\max}\left( \sum_{j = 1}^{N} X_j \right)
\ge (1 + \delta) \mu_{\max} \right]
\le
d \left( \frac{e^{\delta}}{(1 + \delta)^{1 + \delta}}
\right)^{\mu_{\max} / R} .
\]
\end{importedresult}

\subsection{Notation}
\label{subsec:notation}

Figure~\ref{fig:notation_global} lists the symbols used throughout.
Symbols local to one part are introduced at the start of that part, in a table
of the same form with a pointer to where each symbol is defined and used:
Figure~\ref{fig:notation_dynamic} for the data structure of
Section~\ref{sec:dynamic}, Figure~\ref{fig:notation_lewis} for the weights and
potential of Section~\ref{sec:lewis}, and Figure~\ref{fig:notation_ipm} for the
step and initialization of Sections~\ref{sec:potential_and_progress}
and~\ref{sec:overall}.

\begin{figure}[H]
\begin{center}
\begin{tabular}{@{}ll@{}}
\toprule
Symbol & Meaning \\
\midrule
$B$ & bits before and after the decimal point, $B \le d^{1/2}$ \\
$\O(\cdot)$ & hides factors polynomial in $\log n$; powers of $B$ are explicit \\
$x\approx_\epsilon y$ & $e^{-\epsilon}x_i\le y_i\le e^{\epsilon}x_i$ for every $i$ \\
$\rho$ & magnitude of $X$: $\|X\|,\|X^{-1}\|<\rho$; along the path $\rho=2^{O(B)}$ \\
$\MM(n_1,n_2,n_3)$ & arithmetic cost of multiplying $n_1\times n_2$ by $n_2\times n_3$ \\
$A$ & full-column-rank $n\times d$ matrix, $n>d$, no zero rows, so $n\le\nnz(A)$ \\
$\nnz(A)$ & number of nonzero entries of $A$ \\
$\tau(A)$ & leverage scores of $A$: diagonal of $A(A^\top A)^{-1}A^\top$ \\
$\diag(x)$ & diagonal matrix with $x$ on the diagonal \\
$w$ & positive weight vector; along the path $w_i=(\tau_i\phi_i''(x_i))^{-1}$ \\
$W$ & $W=\diag(w)$; the inverse of $A^\top WA$ is maintained, never exposed \\
$\epsilon$ & target accuracy \\
$b,c,l,u$ & LP data: $\min c^\top x$ over $A^\top x = b$, $l \le x \le u$ \\
$\Trunc(x, \ell)$ & $x$ truncated toward zero to $\ell$ bits after the decimal point, Section~\ref{subsec:roundoff} \\
$r$ & number of rounds, $r = \O(d^{1/2}\log(\mu_0/\mu_f))$ \\
$\kappa$ & per-step contraction of the potential $\Psi$, $\kappa = \alpha^{2}\lambda\gamma/(32 C d^{1/2})$ \\
$\E$ & expectation \\
$\mathcal{S}$, $T_{\mathcal{S}}$ & a linear solver handed to a routine, and its bit complexity per call \\
\textsc{Initialize} & with \textsc{Update}, \textsc{Solve}: the calls of Definition~\ref{def:operations} \\
\bottomrule
\end{tabular}
\end{center}
\caption{Global variables.}
\label{fig:notation_global}
\end{figure}

\section{Dynamic Solver via Row Sampling}
\label{sec:dynamic}

\begin{figure}[H]
\begin{center}\small
\begin{tabular}{@{}lp{0.78\textwidth}@{}}
\toprule
Symbol & Meaning \\
\midrule
$w^{(k)}$ & weights handed to the data structure at round $k$, $W^{(k)} = \diag(w^{(k)})$; Definition~\ref{def:tau_stable} \\
$\sigma^{(k)}$ & leverage scores $\tau((W^{(k)})^{1/2}A)$ of the reweighted matrix at round $k$; Definition~\ref{def:tau_stable}, Corollary~\ref{cor:update_counts} \\
$\what$ & the new weight vector handed to \textsc{Update}; Definition~\ref{def:operations}, Lemma~\ref{lem:inv_maintain_explicit} \\
$\zeta$ & stability parameter of Definition~\ref{def:tau_stable}; Theorem~\ref{thm:inv_maintain_stable}, Lemma~\ref{lem:sampled_maintain} \\
$S^{(k)}$ & the private sampling diagonal of round $k$; $c_k$ of its rows change; Lemma~\ref{lem:sampled_maintain} and its proof, proof of Corollary~\ref{cor:update_counts} \\
$c_s$ & sampling constant, $c_s \defeq 10^{6}$; proof of Lemma~\ref{lem:sampled_maintain} \\
$r_0$ & length of a window between restarts, $r_0 = d^{1/2}$; Lemma~\ref{lem:sampled_maintain}, proof of Corollary~\ref{cor:update_counts} \\
$t_{i, j}$, $(\eta, \eta')$ & refresh rounds of coordinate $i$, and the trigger thresholds of a refresh schedule; Corollary~\ref{cor:update_counts} \\
$Q$ & total number of \textsc{Solve} calls, by the caller and inside the structure; Theorem~\ref{thm:inv_maintain_stable}, Lemma~\ref{lem:sampled_maintain}, Corollary~\ref{cor:adaptive} \\
\bottomrule
\end{tabular}
\end{center}
\caption{Variables local to the dynamic solver of Section~\ref{sec:dynamic}.}
\label{fig:notation_dynamic}
\end{figure}

This section gives the data structure that the interior point method of
Section~\ref{sec:potential_and_progress} runs against.
The interior point method sees the data structure only through the three calls
below.
In particular the maintained inverse, and the row sample it is built from, never
leave the data structure: the only thing that leaves it is the output of
\textsc{Solve}, a vector in $\R^{d}$.
Whatever else the caller needs, it computes itself from that output.

The operations of the data structure are the following.

\begin{definition}
\label{def:operations}
An \emph{inverse maintenance data structure} for an $n \times d$ matrix
$A$ of column rank
$d \le n$ holds a positive vector
$w \in \R^{n}_{>0}$, with diagonal
$W = \diag(w)$, and supports:
\begin{enumerate}
\item \label{itm:operations:1} \textsc{Initialize}$(A, w)$, which stores
  $A$ and the initial
  $w$, together with whatever internal approximation of
  $(A^\top W A)^{-1}$ the implementation maintains.
  Nothing about how the caller arrived at $w$ is passed: the only weights the
  structure ever forms are leverage scores of $W^{1/2} A$.
\item \label{itm:operations:2} \textsc{Update}$(\what)$, which given the full
  vector $\what \in \R^{n}_{>0}$ of new weights replaces
  $w$ by
  $\what$ and restores the internal approximation for the new
  $W$.
  The caller is not asked to say which coordinates moved: reading $\what$ costs
  $O(n)$, below the
  $\nnz(A)$ that every round pays anyway, and the data structure decides for
  itself which of its rows to refresh.
\item \label{itm:operations:3} \textsc{Solve}$(b)$, which given a fixed point
  $b \in \R^{d}$ with
  $20 \log \rho$ bits returns a fixed point
  $y \in \R^{d}$ with
  $40 \log \rho$ bits and
\begin{equation}
\label{eq:solve_guarantee}
\left\| y - (A^\top W A)^{-1} b \right\|_{A^\top W A}
\le
\rho^{-10} \left\| (A^\top W A)^{-1} b \right\|_{A^\top W A} .
\end{equation}
There is no accuracy parameter: every solve is at the high accuracy
$\rho^{-10}$.
That this accuracy, and the quadratic form it is an accuracy for, are what the
two callers of \textsc{Solve} need is checked in
Section~\ref{sec:potential_and_progress}, where the calls are made.
\end{enumerate}
\end{definition}

This interface is deliberately smaller than the ones
of~\cite{BLLSSSW21}, which in addition to a
solve expose a call returning the $\ell_2$ norm of the projected vector.
Such a call is there because those data structures maintain the step implicitly,
in a sketched or coordinate wise lazy representation, and so cannot afford to
write down the vector whose norm the potential argument needs.
We are in the earlier regime, where every quantity the caller handles is written
out explicitly in $\R^{d}$ or
$\R^{n}$: \textsc{Solve} returns all
$d$ entries of
$y$, the caller forms
$A y$ in
$\nnz(A)$ time, and any norm of the projected vector is then computed directly,
within the $\nnz(A) + d^2$ that the round pays anyway.
Only the sample and the inverse built from it stay inside; the vectors do not.

The weights handed to \textsc{Update} are not arbitrary either.
They are the sequence produced by the interior point method of
Section~\ref{sec:potential_and_progress}, and the only property of that sequence
the data structure uses is the following, which is the $\sigma$ stability
assumption of~\cite[Definition 4]{LS15}: the multiplicative
movement of the weights is measured in the norm induced by the leverage scores
of the reweighted matrix.
How the caller arrives at its weights is its own business;
Section~\ref{sec:overall} converts its control of them into the condition below.

\begin{definition}
\label{def:tau_stable}
A weight sequence $w^{(0)}, \ldots, w^{(r)} \in \R^{n}_{>0}$ with diagonals
$W^{(k)}$ and induced leverage scores
\[
\sigma^{(k)}
\defeq
\tau\left( (W^{(k)})^{1/2} A \right)
=
\diag\left( (W^{(k)})^{1/2} A \left( A^\top W^{(k)} A \right)^{-1}
A^\top (W^{(k)})^{1/2} \right)
\]
is $\zeta$-leverage stable, for
$0 < \zeta < 0.1$ and magnitude bound
$\rho$, if for every
$k \in [r]$:
\begin{enumerate}
\item \label{itm:tau_stable:1} the operators stay well conditioned:
  $w^{(k)} \in [\rho^{-1}, \rho]^{n}$ and
  $A^\top W^{(k)} A$ has magnitude at most
  $\rho$;
\item \label{itm:tau_stable:2} coordinate wise things are stable:
\[
w^{(k)} \approx_{\zeta} w^{(k - 1)} ;
\]
\item \label{itm:tau_stable:3} the leverage score norm of the relative change is
  small:
\[
\sum_{i} \sigma^{(k)}_{i}
  \log^2\left( \frac{w^{(k)}_i}{w^{(k - 1)}_i}\right)
\le \zeta^2.
\]
\end{enumerate}
\end{definition}

Condition 1 is the analogue of the $\beta$ containment of~\cite[Definitions 3
and 4]{LS15}, and is what lets $W$, the maintained inverse, and
the intermediate quantities be stored with $O(\log \rho)$ bits.
Condition 3 is weaker than the $\ell_2$ condition
$\|\log w^{(k)} - \log w^{(k - 1)}\|_2 \le \zeta$ that a short step method gives
for free, because $\sigma^{(k)}_i \le 1$; the point of measuring in the leverage
score norm is that rows of small leverage are allowed to move often, and those
are exactly the rows that a leverage score sampled sparsifier rarely keeps.
Two facts about $\sigma^{(k)}$ are used repeatedly:
$\sum_i \sigma^{(k)}_i = d$, and for positive diagonals
$W \approx_{\delta} W'$ the scores satisfy
$\tau(W^{1/2} A) \approx_{2 \delta} \tau(W'^{1/2} A)$, since
$\tau_i(W^{1/2} A) = w_i a_i^\top (A^\top W A)^{-1} a_i$ and both factors move
by at most $e^{\pm \delta}$.

The caller of Section~\ref{sec:potential_and_progress} does not hand over an
exactly stable sequence but an estimate of one, so the definition is used with
the following slack.
A sequence $\what^{(0)}, \ldots, \what^{(r)}$ is
\emph{$\zeta$-leverage stable up to noise} if there is a
$\zeta$-leverage stable
$w^{(0)}, \ldots, w^{(r)}$ with
\[
\what^{(k)} \approx_{\zeta / 10} w^{(k)}
\qquad\text{for every } k .
\]
The noisy sequence then satisfies condition 1 with $e^{0.02} \rho$ in place of
$\rho$ and condition 2 with
$1.2 \zeta$ in place of
$\zeta$, its leverage scores agree with those of
$w^{(k)}$ within
$e^{\pm \zeta / 5}$, and condition 3, which need not hold for
$\what^{(k)}$ itself, enters the analysis only through the counting of
Corollary~\ref{cor:update_counts}, which is stated to absorb the noise.
Everything below is written for an exactly stable sequence and holds for a
sequence stable up to noise with these constants.

The main guarantee of this section is that the operations of
Definition~\ref{def:operations} can be supported along any such sequence, even
one generated adaptively from the outputs of the data structure itself, at an
amortized cost per round of one pass over $A$ plus one dense
$d \times d$ operation.

\begin{theorem}
\label{thm:inv_maintain_stable}
Given $n \times d$ full rank matrix
$A$ where
$A^\top A$ of magnitude at most
$\rho$, and
$w^{(0)}, \ldots, w^{(r)}$,
$r \le n^{2}$ an adaptive
$\zeta$-leverage stable sequence with magnitude bound
$\rho$, for some
$0 < \zeta < 0.1$, possibly up to noise, and let
$Q$ be the total number of \textsc{Solve} calls made during these rounds, by the
caller and by the structure itself.
There is a data structure supporting the operations of
Definition~\ref{def:operations} along this sequence, with probability at least
$1 - n^{-5} - 2 Q n^{-3}$, at bit complexity of:
\begin{enumerate}
\item \label{itm:inv_maintain_stable:1}
  $\O( ( \nnz(A) + \MM(d, d, d) ) \log \rho)$ per \textsc{Initialize},
\item \label{itm:inv_maintain_stable:2} Amortized cost of
  $\O( (\nnz(A) + d^2) \log^2 \rho )$ per
  \textsc{Update}$(\what)$,
\item \label{itm:inv_maintain_stable:3} $\O( (d^2+ \nnz(A)) \log^2 \rho )$ per
  $\textsc{Solve}(b)$.
\end{enumerate}
\end{theorem}

The proof occupies the rest of the section.
We implement the operations twice: once without sampling, by maintaining the
inverse of the full $d \times d$ matrix
$A^\top W A$ explicitly (Section~\ref{subsec:exact_maintain}), and once on a
leverage score sampled row subset, which is what brings the amortized cost per
round down to one pass over $A$ plus one dense
$d \times d$ operation (Section~\ref{subsec:sample_maintain}).
The sampled version is only sound against an adversary that does not see the
samples, so we then inject noise into the solver's output, state the conditions
under which the sampling stays hidden, and reduce adaptively generated updates
to oblivious ones (Section~\ref{subsec:randomness}).
Section~\ref{subsec:inv_maintain_stable_pf} bounds the leverage score mass of
the coordinates that the lazy refresh rule touches over a window, which is the
only way the stability of the sequence enters the cost, and combines the three
with a restart schedule to prove the overall guarantee.

\subsection{Exact inverse maintenance and role of matrix
multiplication exponent}
\label{subsec:exact_maintain}

We start with the guarantees of the square unsampled version, which holds the
inverse of $A^\top W A$ itself.
These are mostly a restatement of~\cite{GPV23}.

This is the only part of the algorithm that multiplies dense matrices, so it is
where the cost of fast matrix multiplication enters, and we record here the
assumption under which those costs are simplified,
\begin{equation}
\label{eq:mm_assumption}
\MM(d, d, d) = \O\left( d^{1 + \sqrt{2}} \right) ,
\end{equation}
which is the assumption made by~\cite{LS15}, whose inverse
maintenance we follow: the restart balance closing Section 3.1 there, on page 8,
uses exactly this exponent to conclude that the preprocessing term dominates the
cost of the periodic recomputations.
It is weaker than what the algorithm of~\cite{V89} asks for: the remark
closing Section 1 there, on page 333, notes that the algorithm ``does not use
the fastest known matrix multiplication algorithm; an algorithm that can
multiply two $n \times n$ matrices in
$O(n^{2.4})$ operations suffices'', and Section 3 there invokes that bound, for
both the product of two $n \times n$ matrices and the inverse of a non-singular
$n \times n$ matrix, when recomputing
$(A^\top D A)^{-1}$ at the start of each period.
Both bounds are met by the fast matrix multiplication algorithms available then
and now.
We use two standard consequences of Equation~\eqref{eq:mm_assumption}: costs are
superadditive in the middle dimension,
\[
\MM(d, c_1, d) + \MM(d, c_2, d)
\le
\MM(d, c_1 + c_2, d) ,
\]
and a thin product is a $c / d$ fraction of a square one, so that for
$1 \le c \le d$,
\[
\MM(d, c, d)
=
\O\left( \frac{c}{d} \MM(d, d, d) \right)
=
\O\left( c \, d^{\sqrt{2}} \right) .
\]
In particular $\MM(d, d, d) = \O(d^{2.5})$ and
$\MM(d, d^{1/2}, d) = \O(d^{1/2 + \sqrt{2}}) = \O(d^{1.92})$, both inside the
budget of Theorem~\ref{thm:main}.
Where the two are balanced against each other, and why $d^{1 + \sqrt{2}}$ is
exactly what that balance can afford, is
Section~\ref{subsec:inv_maintain_stable_pf}.

The unsampled inverse maintenance is the following.

\begin{lemma}
\label{lem:inv_maintain_explicit}
Let $A$ be an
$n \times d$ matrix with column rank
$d \leq n$, such that
$A^\top A$ has magnitude at most
$\rho > 10n$, and let every weight vector handed to the structure lie in
$[\rho^{-1}, \rho]^{n}$ coordinatewise.
There is a deterministic data structure that maintains $(A^{\top} W A)^{-1}$
explicitly to error $\rho^{-O(1)}$ whose cost is:
\begin{enumerate}
\item \label{itm:inv_maintain_explicit:1} Initialize:
  $\O( ( \nnz(A) + \MM(n, d, d)) \log\rho)$,
\item \label{itm:inv_maintain_explicit:2} Update from $w$ to
  $\what$:
  $\O( ( n + \MM(d, \|w - \what\|_{0}, d) ) \log\rho )$, where
  $\|w - \what\|_{0}$ is the number of coordinates on which
  $\what$ differs from
  $w$.
\end{enumerate}
\end{lemma}

The cost of \textsc{Solve} is not an item here.
It is the same for this structure as for the sampled one, since both serve a
solve by Richardson iteration against whatever inverse they hold, and it is
accounted once, in Section~\ref{subsec:inv_maintain_stable_pf}.

\begin{proof}
Throughout, every number is kept to $B = \Theta(\log \rho)$ bits, so by
Lemma~\ref{lem:arith} each arithmetic operation costs $\O(\log \rho)$ and the
arithmetic counts below become bit counts after multiplying by that factor.
The weights are positive with $w_i \in [\rho^{-1}, \rho]$, so
$A^\top W A$ has magnitude
$\rho^{O(1)}$, and the same holds of every matrix formed below.

\emph{\textsc{Initialize}.}
Scaling the rows of $A$ by
$w_i^{1/2}$ takes
$\nnz(A)$ operations, and forming the Gram matrix of the scaled matrix takes
$\MM(n, d, d)$.
Inverting it with Imported result~\ref{imp:denseinverse} at accuracy
$\rho^{-C}$, for a constant
$C$ fixed below, takes
$\O(\MM(d, d, d) \log \rho) \le \O(\MM(n, d, d) \log \rho)$.
This is Item~\ref{itm:inv_maintain_explicit:1}.

\emph{\textsc{Update}.}
Let $K \subseteq [n]$ be the coordinates on which
$\what$ differs from
$w$ and
$k = |K|$, let
$A_K$ be the corresponding
$k$ rows of
$A$, and let
$D \defeq \diag(\what_K - w_K)$, so that
\[
A^\top \widehat{W} A
=
A^\top W A + A_K^\top D A_K .
\]
Writing $M \defeq A^\top W A$, the Woodbury identity gives the new inverse as
\begin{equation}
\label{eq:woodbury}
\left( M + A_K^\top D A_K \right)^{-1}
=
M^{-1}
-
M^{-1} A_K^\top
\left( D^{-1} + A_K M^{-1} A_K^\top \right)^{-1}
A_K M^{-1} .
\end{equation}
Reading $\what$ and locating
$K$ takes
$O(n)$ operations.
Forming $A_K M^{-1}$, a
$k \times d$ matrix, takes
$\MM(k, d, d)$; forming and inverting the
$k \times k$ matrix in the middle takes
$\MM(k, d, k) + \MM(k, k, k)$; and the final product of a
$d \times k$, a
$k \times k$, and a
$k \times d$ matrix takes
$\MM(d, k, d)$.
All of these are $\O(\MM(d, k, d))$, and when
$k \ge d$ one may instead recompute the Gram matrix of the changed rows and
invert from scratch, at $\MM(d, k, d) + \MM(d, d, d) = \O(\MM(d, k, d))$.
This is Item~\ref{itm:inv_maintain_explicit:2}.
The middle matrix is invertible with magnitude $\rho^{O(1)}$:
$D^{-1}$ has entries bounded by
$\rho$ in absolute value since
$w, \what \in [\rho^{-1}, \rho]$ coordinatewise, and
$A_K M^{-1} A_K^\top$ is positive semidefinite with norm
$\rho^{O(1)}$.

\emph{Accumulation of error.}
The one point that is not a count of operations is that
Equation~\eqref{eq:woodbury} is applied to the maintained inverse rather than to
the exact one, so rounding errors are fed back into the structure.
What makes $B = \Theta(\log \rho)$ bits enough is that the recurrence is
backward stable, in the following sense.
Let $Z$ be the maintained matrix before an update and
$N \defeq A^\top W A$ the matrix it approximates, and measure the error
backward, as $\|Z^{-1} - N\|_F$.
Equation~\eqref{eq:woodbury} is an algebraic identity valid for every invertible
$Z$ in place of
$M^{-1}$, as long as the middle matrix is invertible, so the exact Woodbury
update of $Z$ is
\[
Z'
=
\left( Z^{-1} + A_K^\top D A_K \right)^{-1},
\qquad\text{hence}\qquad
(Z')^{-1} - N'
=
Z^{-1} - N
\]
for the new matrix $N' = N + A_K^\top D A_K$: exact Woodbury carries the
backward error over unchanged.
What the $B$-bit computation returns is
$\widetilde{Z}'$ with
$\|\widetilde{Z}' - Z'\|_F \le \rho^{O(1)} 2^{-B}$, since every intermediate has
magnitude $\rho^{O(1)}$, there are
$O(1)$ rounding stages, and the inverse of the middle matrix is computed by
Imported result~\ref{imp:denseinverse} to accuracy $2^{-B}$.
Imported result~\ref{imp:perturb}, applied to $\widetilde{Z}'$ and
$Z'$, converts this into
$\|(\widetilde{Z}')^{-1} - (Z')^{-1}\|_F \le \rho^{O(1)} 2^{-B}$, so each update
adds $\rho^{O(1)} 2^{-B}$ to the backward error rather than multiplying it.
Starting from the backward error $\rho^{O(1)} \rho^{-C}$ of \textsc{Initialize},
after at most $\mathrm{poly}(n) \le \rho$ updates the backward error is at most
$\rho^{O(1)} (\rho^{-C} + 2^{-B})$, which taking the constants in
$C$ and in
$B = \Theta(\log \rho)$ large enough is below
$\rho^{-C' - 30}$ for any target
$C'$, and Imported result~\ref{imp:perturb} turns this into a forward error
$\rho^{-C'}$ on the maintained inverse, which is the accuracy claimed.
The magnitude of the maintained matrix stays $\rho^{O(1)}$ throughout by
Item~\ref{itm:imp_perturb:1}.
\end{proof}

Charging \textsc{Update} a rank one update for every changed coordinate is what
the sampled version of the next subsection avoids: there, a change to $w_i$ only
costs anything when row $i$ is in the sample.

\subsection{Maintaining a sparse approximate inverse}
\label{subsec:sample_maintain}

Running Lemma~\ref{lem:inv_maintain_explicit} on $A$ itself is too expensive in
two places: \textsc{Initialize} pays $\MM(n, d, d)$, and \textsc{Update} pays a
$d \times d$ multiplication for every changed coordinate no matter how little
that coordinate matters.
Both are fixed by running it on a leverage score sampled row subset of $A$
instead.
What the sampled structure holds is never returned; it is used only as the
preconditioner inside \textsc{Solve}, and inside the refresh test of its own
\textsc{Update}, so the following lemma is stated in terms of what it maintains
and what the three calls cost, against an adversary that fixes the weight
sequence in advance.

\begin{lemma}
\label{lem:sampled_maintain}
Let $A$,
$\rho$ be as in Lemma~\ref{lem:inv_maintain_explicit} and let
$w^{(0)}, \ldots, w^{(r_0)}$ be a
$\zeta$-leverage stable sequence chosen by an oblivious adversary, that is, one
fixed independently of the data structure's random coins, and let $Q$ be the
total number of \textsc{Solve} calls made during these rounds, by the caller and
by the structure itself.
There is a data structure which, with probability at least
$1 - n^{-7} - 2 Q n^{-3}$, maintains at every round
$k$ a diagonal
$S^{(k)} \succeq 0$ supported on
$\O(d \log n)$ rows, together with the exact inverse of
$A^\top S^{(k)} W^{(k)} A$, such that
\begin{equation}
\label{eq:sample_spectral}
A^\top S^{(k)} W^{(k)} A
\approx_{0.1}
A^\top W^{(k)} A ,
\end{equation}
with initialization bit complexity $\O( ( \nnz(A) + \MM(d, d, d) ) \log \rho )$
and amortized update bit complexity
\[
\O\left( \left( \nnz(A) + d^{2} \right) \log^{2} \rho \right) .
\]
\end{lemma}

\begin{proof}
Keeping row $i$ independently with probability
$p_i = \min\{1, c_s \sigma^{(k)}_i \log n\}$,
$c_s \defeq 10^{6}$, and setting
$S^{(k)}_{ii} = 1 / p_i$ on the kept rows gives
Equation~\eqref{eq:sample_spectral} with probability $1 - n^{-9}$, with
$\O(d \log n)$ rows retained, since
$\|\sigma^{(k)}\|_1 = d$; this is leverage score sampling,
\cite{SS11:journal}, and the computation at the end of this
proof, from the matrix Chernoff inequality of Imported
result~\ref{imp:matrix_chernoff}, shows that the concentration bound needs only
that the rows are kept independently with $p_i$ between
$\min\{1, (c_s / 100) \sigma^{(k)}_i \log n\}$ and
$\min\{1, 100 c_s \sigma^{(k)}_i \log n\}$.
It is here that obliviousness is used, since the $\sigma^{(k)}$ defining
$p_i$ must not depend on the coins that produced the earlier samples.

Leverage scores are estimated, here and below, by the standard
Johnson--Lindenstrauss argument: since
\[
\sigma^{(k)}_i
=
\left\| (W^{(k)})^{1/2} A \left( A^\top W^{(k)} A \right)^{-1}
a_i \, (w^{(k)}_i)^{1/2} \right\|_2^2 ,
\]
a Gaussian $G$ with
$O(\log n)$ rows gives
$\|G (W^{(k)})^{1/2} A (A^\top W^{(k)} A)^{-1} a_i (w^{(k)}_i)^{1/2}\|_2^2 \approx_{0.1} \sigma^{(k)}_i$
for all $i$ with probability
$1 - n^{-9}$, and the
$n$ numbers on the left are obtained from
$O(\log n)$ solves in
$A^\top W^{(k)} A$, one per row of
$G (W^{(k)})^{1/2} A$, followed by one multiplication of
$A$ against the
$O(\log n)$ solutions.
With a solver at accuracy $\rho^{-10}$ costing
$\O((\nnz(A) + d^{2}) \log^{2} \rho)$ per call, this is
$\O((\nnz(A) + d^{2}) \log^{2} \rho)$ bit operations, all numbers having
$\O(\log \rho)$ bits by Item~\ref{itm:tau_stable:1}.

The probabilities needed to build the first sample are computed this way with a
one-off solver: a sparse embedding \cite{CW17:journal} of
$(W^{(0)})^{1/2} A$ with
$\O(d)$ rows, applied at
$\O(\nnz(A))$ bit operations and a constant factor spectral approximation of
$A^\top W^{(0)} A$ by the embedding guarantee, preconditions
$\O(\log \rho)$ Richardson steps with the dense inverse of its Gram matrix from
Imported result~\ref{imp:denseinverse}, at cost
$\O((\nnz(A) + \MM(d, d, d)) \log \rho)$; this is the only place the full
$n$-row matrix is touched by anything other than a matrix vector product.
The inverse of the sampled $d \times d$ matrix is then formed once, inside the
\textsc{Initialize} of Lemma~\ref{lem:inv_maintain_explicit} applied to the
$\O(d \log n)$ sampled rows, which is
\[
\O(\MM(d, d, d) \log \rho)
\quad\text{rather than}\quad
\O(\MM(n, d, d) \log \rho) .
\]
Item~\ref{itm:tau_stable:1} keeps the sampled matrix of magnitude $\rho^{O(1)}$
throughout by Equation~\eqref{eq:sample_spectral}, so $\O(\log \rho)$ bits
suffice everywhere, and Item~\ref{itm:tau_stable:2} keeps the sampling
probabilities themselves accurate between refreshes, by the
$\tau(W^{1/2} A) \approx_{2 \delta} \tau(W'^{1/2} A)$ fact recorded after
Definition~\ref{def:tau_stable}.

An \textsc{Update} is served in two parts.
First the new sampling probabilities are recomputed: the test below needs the
current $\sigma^{(k)}$, at constant accuracy only, so it is one
Johnson--Lindenstrauss estimate as above with the structure's own \textsc{Solve}
as the solver, at
\[
\O\left( \left( \nnz(A) + d^{2} \right) \log^{2} \rho \right) ,
\]
which also subsumes the $O(n)$ spent reading
$\what$.
Second, rows are resampled lazily, as in~\cite[Algorithm 3]{LS15}:
coordinate $i$ is refreshed only when
$w_i$, or the estimate of
$\sigma_i$, has moved by more than a factor of
$2$ since it was last refreshed, and a refresh changes the maintained matrix
only when it changes the kept status or the scaling of a sampled row.
Each such change is a rank one update to $A^\top S W A$: the structure of
Lemma~\ref{lem:inv_maintain_explicit} is run on the kept rows with weights
$s_i w_i \in [\rho^{-1}, n \rho]$, and a row entering or leaving the sample is a
weight change between $0$ and
$s_i w_i$, for which Equation~\eqref{eq:woodbury} holds verbatim, the middle
matrix being invertible because the updated matrix is positive definite.
So $c_k$ such changes cost
$\O(\MM(d, c_k, d) \log \rho)$ through the \textsc{Update} of
Lemma~\ref{lem:inv_maintain_explicit}.
Bounding $\sum_k c_k$, and so turning the second part into an amortized
$\O(d^{2} \log \rho)$, is what Section~\ref{subsec:inv_maintain_stable_pf} does;
adding the two parts gives the claimed amortized bound.

The sample of round $k$ is not drawn afresh: row
$i$ is kept or not according to coins flipped at the round
$t_i \le k$ at which it was last refreshed, with the probability
$p_i^{(t_i)}$ of that round, and the estimate that decided when to refresh was
itself computed against the maintained form, hence from the earlier coins.
It is worth saying why the sampling argument above still covers this.
Every \textsc{Solve} of the structure, including the $O(\log n)$ internal ones
behind each refresh test, is implemented by the \textsc{NoisySolver} of
Section~\ref{subsec:randomness}, whose output is within total variation distance
$n^{-3}$ of the output of an ideal solver that is a function of
$(A, W^{(k)}, b)$ and fresh Gaussian noise only, by Imported
result~\ref{imp:noisysolver_tv}, and its fixed point implementation, described
after Corollary~\ref{cor:adaptive}, fails to meet the hypotheses of that result
with probability at most $2 n^{-4}$ per call.
So, at the cost of $Q (n^{-3} + 2 n^{-4}) \le 2 Q n^{-3}$ in probability, we may
analyze the process in which every estimate is computed by the ideal solver.
In that process the refresh decisions are functions of the weights, which the
oblivious adversary fixed in advance, and of fresh noise, so they are
independent of the sampling coins; the sampling coins of distinct rows are
independent; and the refresh rule, together with the $0.1$-accuracy of the
estimates, keeps $p_i^{(t_i)}$ within a constant factor of
$c_s \sigma^{(k)}_i \log n$, or equal to one, for every
$i$ and every round
$k \ge t_i$ at which row
$i$ has not been refreshed again: writing
$\widehat{\sigma}^{(k)}_i \approx_{0.1} \sigma^{(k)}_i$ for the estimates, the
rule not having fired between $t_i$ and
$k$ means
$\widehat{\sigma}^{(k)}_i \in [1/2, 2] \, \widehat{\sigma}^{(t_i)}_i$, so
\[
\sigma^{(k)}_i
\in
\left[ \frac{e^{-0.2}}{2}, 2 e^{0.2} \right] \sigma^{(t_i)}_i
\subseteq
\left[ \frac{1}{10}, 10 \right] \sigma^{(t_i)}_i ,
\]
and therefore
\[
p_i^{(t_i)}
=
\min\left\{ 1, c_s \widehat{\sigma}^{(t_i)}_i \log n \right\}
\ge
\min\left\{ 1, \frac{c_s e^{-0.1}}{10} \sigma^{(k)}_i \log n \right\}
\ge
\min\left\{ 1, \frac{c_s}{100} \sigma^{(k)}_i \log n \right\} ,
\]
and symmetrically $p_i^{(t_i)} \le \min\{1, 100 c_s \sigma^{(k)}_i \log n\}$.
So at round $k$ the rows are kept independently, with probabilities fixed
independently of the coins, each at least the probability of leverage score
sampling at the constant $c_s / 100$ and at most that at
$100 c_s$.
We now carry out the concentration computation.
Write $N \defeq A^\top W^{(k)} A$ and
$u_i \defeq N^{-1/2} (w^{(k)}_i)^{1/2} a_i$, so that
$\sum_i u_i u_i^\top = I$ and
$\|u_i\|_2^{2} = (w^{(k)}_i)^{1/2} a_i^\top N^{-1} a_i (w^{(k)}_i)^{1/2} = \sigma^{(k)}_i$,
and Equation~\eqref{eq:sample_spectral} is the statement
\[
e^{-0.1} I
\preceq
\sum_i S^{(k)}_{ii} u_i u_i^\top
\preceq
e^{0.1} I .
\]
Set
\[
R
\defeq
\frac{100}{c_s \log n}
=
\frac{1}{10^{4} \log n} .
\]
We write $\sum_i S^{(k)}_{ii} u_i u_i^\top$ as a sum of independent positive
semidefinite matrices of norm at most $R$.
A row with $p_i^{(t_i)} = 1$ is kept deterministically with
$S^{(k)}_{ii} = 1$; its term
$u_i u_i^\top$ has norm
$\sigma^{(k)}_i \le 1$, and we split it into
$m_i \defeq \lceil \sigma^{(k)}_i / R \rceil$ equal summands
$u_i u_i^\top / m_i$, each of norm at most
$R$, which are constant and hence independent of everything.
For a row with $p_i^{(t_i)} < 1$, let
$\xi_i \in \{0, 1\}$ be its coin and
\[
X_i
\defeq
\frac{\xi_i}{p_i^{(t_i)}} \, u_i u_i^\top ,
\]
so that $\E X_i = u_i u_i^\top$ as
$\E \xi_i = p_i^{(t_i)}$, and, since
$p_i^{(t_i)} < 1$ forces
$p_i^{(t_i)} \ge (c_s / 100) \sigma^{(k)}_i \log n$,
\[
\lambda_{\max}(X_i)
\le
\frac{\sigma^{(k)}_i}{p_i^{(t_i)}}
\le
\frac{100}{c_s \log n}
=
R .
\]
The coins of distinct rows are independent, so all summands are independent, and
their expectations sum to $\sum_i u_i u_i^\top = I$, whence
$\mu_{\min} = \mu_{\max} = 1$ in Imported result~\ref{imp:matrix_chernoff}.
Apply it with $\delta_{+} \defeq e^{0.1} - 1$ in the upper tail and
$\delta_{-} \defeq 1 - e^{-0.1}$ in the lower tail.
Since $\ln(1 + \delta_{+}) = 0.1$ and
$\ln(1 - \delta_{-}) = -0.1$, the two bases are
\[
\frac{e^{\delta_{+}}}{(1 + \delta_{+})^{1 + \delta_{+}}}
=
\exp\left( \delta_{+} - 0.1 (1 + \delta_{+}) \right)
=
\exp\left( 0.9 e^{0.1} - 1 \right)
\le
e^{-0.005}
\]
and
\[
\frac{e^{-\delta_{-}}}{(1 - \delta_{-})^{1 - \delta_{-}}}
=
\exp\left( - \delta_{-} + 0.1 (1 - \delta_{-}) \right)
=
\exp\left( 1.1 e^{-0.1} - 1 \right)
\le
e^{-0.004} ,
\]
using $0.9 e^{0.1} \le 0.9947$ and
$1.1 e^{-0.1} \le 0.9954$.
Raised to the power $1 / R = 10^{4} \log n$, each is at most
$e^{-40 \log n} = n^{-40}$, so
\[
\Pr\left[ \sum_i S^{(k)}_{ii} u_i u_i^\top
\not\approx_{0.1} I \right]
\le
d \, n^{-40} + d \, n^{-40}
\le
n^{-9} ,
\]
as $2 d \le n^{2}$.
The complementary event is Equation~\eqref{eq:sample_spectral} at round $k$.
The number of kept rows has mean at most
$\sum_i p_i^{(t_i)} \le 100 c_s \log n \sum_i \sigma^{(k)}_i = 10^{8} d \log n$,
and the scalar Chernoff bound for a sum of independent indicators,
$\Pr[X \ge 2 \mu'] \le e^{-\mu' / 3}$ for any
$\mu' \ge \E X$, at
$\mu' \defeq 10^{8} d \log n + 30 \log n$ gives at most
$2 \cdot 10^{8} d \log n + 60 \log n = \O(d \log n)$ kept rows except with
probability $n^{-10}$.
Together with the $n^{-9}$ of the leverage score estimate, a round fails with
probability at most $n^{-9} + n^{-9} + n^{-10} \le n^{-8}$, and a union bound
over the $r_0 \le n$ rounds and the
$2 Q n^{-3}$ from the coupling give the probability claimed.
On the event that Equation~\eqref{eq:sample_spectral} holds, the maintained
inverse is a constant factor preconditioner and \textsc{Solve} meets
Equation~\eqref{eq:solve_guarantee}, whatever the coins happened to be.
What obliviousness is used for is the weights and not the probabilities:
$w^{(k)}$ must not be chosen in response to the coins, while the estimate of
$\sigma^{(k)}$ has only to be accurate and, up to the coupling, independent of
the coins.

\textsc{Solve} runs a fixed number $O(\log \rho)$ of steps of preconditioned
Richardson iteration against the maintained inverse, so that in exact arithmetic
it is a linear map of $b$.
The exact operator applied at each step is $A^\top W^{(k)} A$, at the cost of
one multiplication by $A$ and one by
$A^\top$, so the guarantee Equation~\eqref{eq:solve_guarantee} is with respect
to the true matrix, and only the rate of convergence depends on the sample
through Equation~\eqref{eq:sample_spectral} and Imported
result~\ref{imp:perturb}.
The same holds of the leverage score estimates the refresh test runs against it,
and of anything the caller computes from the output of \textsc{Solve}: the
sample only affects how fast the solves converge, not the point they converge
to, so every such quantity is a deterministic function of the inputs up to the
requested accuracy, even though it is computed from a full pass over $A$ while
the preconditioner behind \textsc{Solve} is built from $\O(d \log n)$ rows.

Since the guarantee is relative and the arithmetic is fixed point, the iteration
is run on the scaled right hand side $b' \defeq 2^{-j} b$, with
$j$ the integer for which
$\|b'\|_{\infty} \in [1/2, 1)$, on numbers with
$20 \log \rho$ bits, and the result is multiplied by
$2^{j}$; for
$b = 0$ the answer
$0$ is returned directly.
As $b \ne 0$ has
$20 \log \rho$ bits, its largest entry lies in
$[\rho^{-20}, \rho^{20})$, so
$- 20 \log \rho < j \le 20 \log \rho$.
Write $M \defeq A^\top W^{(k)} A$.
The error $y_t - M^{-1} b'$ of the
$t$-th iterate contracts in the
$M$ norm at rate
$c \le 1/2$ by Equation~\eqref{eq:sample_spectral} and Imported
result~\ref{imp:perturb}, so $\|y_t\|_{M} \le 2 \|M^{-1} b'\|_{M}$ and
$\|y_t\|_2 \le 2 \|M^{-1}\| \, \|b'\|_2 \le 2 \rho d^{1/2}$, up to the
truncation errors bounded below; the products $A y_t$,
$W^{(k)} A y_t$ and
$M y_t$ have magnitude at most
$\rho^{3}$ times this, so every number formed has magnitude below
$\rho^{5} < \rho^{20}$ and fits in
$20 \log \rho$ bits; the output
$2^{j} y'$ has magnitude below
$\rho^{25}$ and at most
$20 \log \rho + 20 \log \rho - 1 < 40 \log \rho$ bits after the point, so it is
a number with $40 \log \rho$ bits as Item~\ref{itm:operations:3} requires, and
costs $O(\log \rho)$ per entry to write.
One Richardson step applies $A$,
$W^{(k)}$,
$A^\top$ and the maintained
$d \times d$ inverse, each product being formed exactly and truncated once to
$20 \log \rho$ bits; each of these
$O(1)$ truncations moves a vector of dimension at most
$n$ by at most
$n^{1/2} \rho^{-20}$ in
$\ell_2$, and is then amplified by the operators applied after it, whose norms
multiply to at most $2 \rho^{3}$, so the step is computed to within
$O(n^{1/2} \rho^{-17})$ in
$\ell_2$ and hence within
$O(n^{1/2} \rho^{-16.5}) \le \rho^{-15}$ in the
$M$ norm, using
$\|M\| \le \rho$ and
$\rho \ge n$.
Against
\[
\left\| M^{-1} b' \right\|_{M}
=
\|b'\|_{M^{-1}}
\ge
\rho^{-1/2} \|b'\|_2
\ge
\frac{\rho^{-1/2}}{2}
\]
this is a relative error of at most $2 \rho^{-14.5}$ per step.
A contraction of rate $c \le 1/2$ run with additive error
$e$ per step converges to within
$e / (1 - c) \le 2 e$ of the exact solution, so
$\O(\log \rho)$ steps reach relative accuracy
$\rho^{-10}$, and the truncation errors never move an iterate by more than
$4 \rho^{-14.5} \|M^{-1} b'\|_{M}$, which justifies the magnitude bound above.
The same scaling is used by every solve the data structure runs; the two inner
solves of \textsc{NoisySolver} below, which are asked for accuracy
$\rho^{-20} / (32 n^{7}) \ge \rho^{-28}$, run on numbers with
$60 \log \rho$ bits instead, which puts the floor of the same computation below
$\rho^{-50}$ and changes the cost by a constant factor.
\end{proof}

\subsection{Hiding the randomness}
\label{subsec:randomness}

Lemma~\ref{lem:sampled_maintain} assumes an oblivious adversary: its sampling is
only valid if the $w^{(k)}$ handed to round
$k$ are independent of the coin flips used in rounds
$1, \ldots, k - 1$.
This fails as stated for our application, where $w^{(k)}$ is computed by the
interior point method from the approximate solutions the data structure itself
returned, and where a single leaked bit about which rows were sampled could be
amplified into a weight sequence that forces a refresh every round.

So we ask for two conditions on the interface of
Definition~\ref{def:operations}, which together make the updates oblivious to
the sampling:
\begin{enumerate}
\item the only outputs of the data structure that depend on the sample are the
  outputs of \textsc{Solve}, which is the only output at all; anything the
  caller computes from them is computed from a full pass over $A$ with the
  sample entering only as a preconditioner, whose accuracy requirement is met
  deterministically;
\item the output of \textsc{Solve} is within total variation distance $n^{-3}$
  of a distribution that is a function of $(A, W, b)$ alone.
\end{enumerate}
The first is exactly the last assertion of Lemma~\ref{lem:sampled_maintain}.
The second is achieved by injecting noise into the solver's output, following
\cite[Section 5]{LS15}: \textsc{Solve} is implemented as a call to
\textsc{NoisySolver} below, and is then the only channel through which anything
about the sampling could leak.

The following is Lemma 16 of \cite{LS15}, the
\textsc{NoisySolver}, stated there for $A^\top A$ and applied here to
$W^{1/2} A$; a linear solver is one whose output is a linear function of its
right hand side.

\begin{importedresult}
\label{imp:noisysolver}
Let $\mathcal{S}$ be a
$T_{\mathcal{S}}$-time linear solver for
$A^\top W A$.
The algorithm that, on input $b$ and accuracy
$\epsilon$, returns
\[
y \defeq y_1 + \alpha_{\mathrm{noise}} y_2,
\qquad
y_1 \defeq \mathcal{S}(b, \epsilon_1),
\qquad
y_2 \defeq \mathcal{S}(A^\top W^{1/2} \xi, \epsilon_2),
\qquad
\alpha_{\mathrm{noise}}
\defeq \frac{\epsilon}{8 n^{1/2}} \|y_1\|_{A^\top W A},
\]
with $\xi \sim N(0, I_n)$,
$\epsilon_1 = \epsilon^2 / (32 n^7)$ and
$\epsilon_2 = \epsilon^2 / (12 d n^6)$, is an
$\epsilon$-accurate
$O(\nnz(A) + T_{\mathcal{S}})$-time solver for
$A^\top W A$.
\end{importedresult}

The following, Theorem 17 of \cite{LS15}, is the
indistinguishability of the noisy solver from an ideal one.

\begin{importedresult}
\label{imp:noisysolver_tv}
Let $\textsc{IdealSolver}(b) \defeq (A^\top W A)^{-1} b + c$ where
$c \sim N(0, \beta_{\mathrm{noise}}^2 (A^\top W A)^{-1})$ and
$\beta_{\mathrm{noise}} = \frac{\epsilon}{8 n^{1/2}} \|(A^\top W A)^{-1} b\|_{A^\top W A}$.
Then the total variation distance between the outputs of \textsc{NoisySolver}
and \textsc{IdealSolver} is at most $n^{-3}$; consequently no algorithm making
at most $n^{2}$ calls can distinguish them.
\end{importedresult}

Both statements are for a linear $\mathcal{S}$: the proof of the second uses
that $y_2$ is a linear image of
$\xi$ and so Gaussian.
The noise is a $(1 \pm \epsilon)$ perturbation in the
$A^\top W A$ norm, so it is absorbed into the solve tolerance of every statement
that consumes the solver, and its distribution depends on the data structure's
randomness not at all.
Since \textsc{IdealSolver} is a deterministic function of $(A, W, b)$ plus fresh
noise, any weight sequence generated from its outputs is independent of the
sampling, and the obliviousness hypothesis of Lemma~\ref{lem:sampled_maintain}
may be dropped.

This gives the guarantee for adaptive weight sequences.

\begin{corollary}
\label{cor:adaptive}
Lemma~\ref{lem:sampled_maintain} holds for weight sequences $w^{(k)}$ produced
by any algorithm that sees the data structure only through the operations of
Definition~\ref{def:operations}, provided every \textsc{Solve} is implemented by
\textsc{NoisySolver} of Imported result~\ref{imp:noisysolver}, at the cost of
$\O(\nnz(A) \log \rho)$ extra bit operations per call; the failure probability
is $n^{-7} + 2 Q n^{-3}$ as in that lemma, with
$Q$ now counting the caller's \textsc{Solve} calls as well.
\end{corollary}

\begin{proof}
By Imported result~\ref{imp:noisysolver_tv} each call is within total variation
distance $n^{-3}$ of \textsc{IdealSolver}, so by a union bound over the
$Q$ calls the whole run may be coupled with a run against \textsc{IdealSolver},
except with probability $2 Q n^{-3}$; the internal calls of the structure are
already coupled this way in the proof of Lemma~\ref{lem:sampled_maintain}, so
the coupling extends to the caller's calls at no further cost than counting them
in $Q$.
Under the coupling every $w^{(k)}$ is a function of
$A$, of
$w^{(0)}, \ldots, w^{(k - 1)}$, and of Gaussian noise independent of the
sampling used inside the data structure, so the sequence is oblivious in the
sense required by Lemma~\ref{lem:sampled_maintain}, which applies verbatim.
The extra cost per call is one additional solve and one multiplication by
$A^\top$, and by Imported result~\ref{imp:noisysolver} the accuracy degrades by
a constant factor, absorbed into $\rho^{-O(1)}$.
\end{proof}

Three points concern fixed point; throughout, $M \defeq A^\top W A$ and
\textsc{NoisySolver} is run at accuracy $\epsilon = \rho^{-10} / 2$, so that
$\epsilon_1 = \epsilon^{2} / (32 n^{7}) \ge \rho^{-28}$ and
$\epsilon_2 = \epsilon^{2} / (12 d n^{6}) \ge \rho^{-28}$ for
$\rho \ge \max(n, 2^{7})$.
First, the output of \textsc{NoisySolver} is formed on the scaled right hand
side $b' = 2^{-j} b$ of the proof of Lemma~\ref{lem:sampled_maintain}, on
numbers with $60 \log \rho$ bits, then truncated to
$20 \log \rho$ bits and multiplied by
$2^{j}$, and the run it is coupled with is against
$2^{j} \Trunc(\textsc{IdealSolver}(b'), 20 \log \rho)$; truncation and scaling
are one map applied to both outputs, and total variation distance does not
increase under a common map, so the bound $n^{-3}$ of Imported
result~\ref{imp:noisysolver_tv} is unaffected, whatever the size of the noise
relative to $\rho^{-20}$, and the truncated ideal output is still a function of
$(A, W, b)$ and fresh noise.
The truncation moves the scaled output by at most $d^{1/2} \rho^{-20}$ in
$\ell_2$, hence by at most
$d^{1/2} \rho^{-19.5} \le \rho^{-18} \|M^{-1} b'\|_{M}$ in the
$M$ norm, by the lower bound
$\|M^{-1} b'\|_{M} \ge \rho^{-1/2} / 2$ displayed in that proof, so the returned
vector is within Equation~\eqref{eq:solve_guarantee}.
Second, $\xi$ is the only place the algorithm needs randomness of more than
$O(\log n)$ bits per row: it is sampled with
$40 \log \rho$ bits per entry, that is, the algorithm holds the truncation
$\xi' = \Trunc(\xi, 40 \log \rho)$ of an exact Gaussian vector
$\xi$, a discrete distribution sampled with
$\mathrm{poly}(\log \rho)$ bit operations per entry, on the event that the
entries of $\xi$ have magnitude below
$\rho^{40}$, which has probability
$1 - n e^{-\rho^{80} / 2}$.
Third, the fixed point computation is not linear in its right hand side, while
Imported result~\ref{imp:noisysolver_tv} is for a linear $\mathcal{S}$.
Let $y'' \defeq y_1 + \alpha_{\mathrm{noise}} Q_2 A^\top W^{1/2} \xi$ be the
output of the exact arithmetic \textsc{NoisySolver} on $b'$ and the exact
Gaussian $\xi$, with the exact arithmetic Richardson iteration, a linear map
$Q_2$ with
$\|Q_2 v - M^{-1} v\|_{M} \le \epsilon_2 \|M^{-1} v\|_{M}$ for every
$v$, as its inner solver; Imported results~\ref{imp:noisysolver}
and~\ref{imp:noisysolver_tv} apply to $y''$.
The computed output differs from $y''$ in two ways.
Replacing $\xi$ by
$\xi'$ moves
$\alpha_{\mathrm{noise}} Q_2 A^\top W^{1/2} \xi$ by at most
$2 \alpha_{\mathrm{noise}} \|A^\top W^{1/2} (\xi - \xi')\|_{M^{-1}} \le 2 \alpha_{\mathrm{noise}} n^{1/2} \rho^{-40} \le \rho^{-48}$
in the $M$ norm, since
$\|A^\top W^{1/2} z\|_{M^{-1}}^{2} = z^\top W^{1/2} A M^{-1} A^\top W^{1/2} z \le \|z\|_2^{2}$
and $\alpha_{\mathrm{noise}} \le \epsilon \|y_1\|_{M} / 8 \le \rho^{-9}$ by
$\|y_1\|_{M} \le 2 \|M^{-1} b'\|_{M} \le 2 \rho^{1/2} d^{1/2}$.
The truncations of the inner solves, of the scale $\alpha_{\mathrm{noise}}$ and
of the final sum are each a relative $\rho^{-50}$ of the quantity truncated, by
the proof of Lemma~\ref{lem:sampled_maintain} run with $60 \log \rho$ bits; by
the $\chi^{2}$ tail bound in the proof of Imported result~\ref{imp:noisysolver},
$\|\xi\|_2^{2} \le 16 n$ except with probability
$n^{-4}$, and then
$\alpha_{\mathrm{noise}} \|y_2\|_{M} \le (\epsilon / 8 n^{1/2}) \|y_1\|_{M} \cdot 4 n^{1/2} (1 + \epsilon_2) \le \|M^{-1} b'\|_{M}$,
using $\|\xi'\|_2 \le \|\xi\|_2$, so these truncations move the output by at
most $\rho^{-48} \|M^{-1} b'\|_{M} \le \rho^{-47}$ in the
$M$ norm.
Altogether the computed output is within $2 \rho^{-47}$ of
$y''$ in the
$M$ norm and hence within
$\rho^{-46}$ of it in every coordinate, using
$\|M^{-1}\| \le \rho$.
The truncations to $20 \log \rho$ bits of the two then agree unless some
coordinate of $y''$ lies within
$\rho^{-46}$ of a multiple of
$h \defeq \rho^{-20}$.
Coordinate $i$ of
$y''$ is Gaussian with mean
$(y_1)_i$ and standard deviation
$s_i = \alpha_{\mathrm{noise}} (Q_2 M Q_2^\top)_{ii}^{1/2} \ge \alpha_{\mathrm{noise}} \rho^{-1/2} / 2 \ge \rho^{-13}$,
where $Q_2 M Q_2^\top \approx_{1} M^{-1}$ by the accuracy of the linear
$Q_2$ and
$\alpha_{\mathrm{noise}} \ge (\epsilon / 8 n^{1/2}) (1 - \epsilon_1) \|M^{-1} b'\|_{M} \ge \rho^{-11} / (64 n^{1/2})$.
For a unimodal density $f$ with
$\sup f \le 1 / s$ and cells of width
$h$, the maxima of
$f$ over the cells sum to at most
$1 / h + 2 / s$, and a
$\delta$-neighbourhood of a lattice point meets two cells, so the probability of
landing within $\delta$ of the lattice
$h \mathbb{Z}$ is at most
$4 \delta (1 / h + 2 / s)$; with
$\delta = \rho^{-46}$ this is at most
$4 \rho^{-46} (\rho^{20} + 2 \rho^{13}) \le \rho^{-25}$ per coordinate and
$\rho^{-24}$ over the
$d \le \rho$ coordinates.
The truncated computed output is therefore within total variation distance
$n^{-3} + n^{-4} + \rho^{-24} \le n^{-3} + 2 n^{-4}$ of
$2^{j} \Trunc(\textsc{IdealSolver}(b'), 20 \log \rho)$, and the excess
$2 n^{-4}$ over the imported
$n^{-3}$ is what is charged per call in the proof of
Lemma~\ref{lem:sampled_maintain}; the coupled ideal output remains a function of
$(A, W, b)$ and fresh noise.

\subsection{Overall accounting and adaptive solver interface}
\label{subsec:inv_maintain_stable_pf}

Condition 3 of Definition~\ref{def:tau_stable} converts into a bound on the
leverage score mass of the coordinates that a lazy refresh rule touches over a
window of rounds.
This bound is the only way the stability of the sequence enters the cost.

The following corollary bounds the refreshed mass.
A refresh schedule is, for each coordinate $i$, an increasing sequence of rounds
$t_{i, 1} < t_{i, 2} < \cdots$ in the window, the rounds at which the coordinate
is refreshed.
It is \emph{$(\eta, \eta')$-triggered}, for
$0 < \eta \le \eta'$, if for every consecutive pair
$t = t_{i, j} < t' = t_{i, j + 1}$ of refresh rounds of a coordinate
$i$ the refresh at
$t'$ was caused by a movement of at least
$\eta$ and not by a smaller one:
\[
w^{(t')}_i \not\approx_{\eta} w^{(t)}_i
\quad\text{or}\quad
\sigma^{(t')}_i \not\approx_{\eta} \sigma^{(t)}_i ,
\qquad\text{while}\qquad
w^{(k)}_i \approx_{\eta'} w^{(t)}_i
\quad\text{and}\quad
\sigma^{(k)}_i \approx_{\eta'} \sigma^{(t)}_i
\quad\text{for } t \le k < t' .
\]
This is what a lazy rule produces when it refreshes a coordinate at the first
round its weight or leverage score is seen to have moved by about a constant
factor, the gap between $\eta$ and
$\eta'$ being the accuracy of the estimate it acts on.

\begin{corollary}
\label{cor:update_counts}
Let $w^{(k_0)}, \ldots, w^{(k_0 + r_0)}$ be
$\zeta$-leverage stable and let a refresh schedule on this window be
$(\eta, \eta')$-triggered.
Then the leverage score mass refreshed over the window, not counting the first
refresh of each coordinate, satisfies
\[
\sum_i \sum_{j \ge 1} \sigma^{(t_{i, j + 1})}_i
\le
\frac{30 e^{3 \eta'} r_0^{2} \zeta^{2}}{\eta^{2}} .
\]
The same bound holds, up to a constant factor, for a sequence
$\what^{(k_0)}, \ldots, \what^{(k_0 + r_0)}$ that is
$\zeta$-leverage stable up to noise, with
$\what$ and its own leverage scores in the trigger conditions, provided
$\eta \ge \zeta$.
\end{corollary}

\begin{proof}
Write $\delta^{(m)} \defeq \log w^{(m + 1)} - \log w^{(m)}$ for the single round
increments, so that condition 3 reads
$\|\delta^{(m)}\|_{\sigma^{(m + 1)}} \le \zeta$, and condition 2 with the fact
recorded after Definition~\ref{def:tau_stable} gives
$\sigma^{(m)} \approx_{2 \zeta} \sigma^{(m + 1)}$.
Fix a consecutive pair $t < t'$ of refresh rounds of a coordinate
$i$.
The second half of the trigger condition, together with
$\sigma^{(t')}_i \approx_{2 \zeta} \sigma^{(t' - 1)}_i$, gives
\begin{equation}
\label{eq:refresh_comparable}
\sigma^{(k)}_i \approx_{2 \eta' + 2 \zeta} \sigma^{(t')}_i
\qquad\text{for every } t \le k \le t' ,
\end{equation}
so the leverage score of a coordinate is comparable to its value at the refresh
throughout the interval that ends at that refresh.
This is the only comparison between different rounds that the proof uses; the
weights themselves may drift by large factors over the window.

\emph{Weight triggers.}
If the pair is triggered by the weight, then
$|\sum_{m = t}^{t' - 1} \delta^{(m)}_i| \ge \eta$, and by Cauchy--Schwarz over
the $t' - t \le r_0$ terms,
\[
\eta^{2} \sigma^{(t')}_i
\le
r_0 \sum_{m = t}^{t' - 1} \sigma^{(t')}_i \left( \delta^{(m)}_i \right)^2
\le
e^{2 \eta' + 2 \zeta} r_0
\sum_{m = t}^{t' - 1} \sigma^{(m + 1)}_i \left( \delta^{(m)}_i \right)^2 ,
\]
by Equation~\eqref{eq:refresh_comparable} at $k = m + 1$.

\emph{Leverage score triggers.}
Interpolate between consecutive weight vectors along
$w(s) \defeq w^{(m)} e^{(s - m) \delta^{(m)}}$ for
$s \in [m, m + 1]$, a continuous path through all the weights of the window,
with $W(s) = \diag(w(s))$,
$M(s) \defeq A^\top W(s) A$,
$\sigma(s) \defeq \tau(W(s)^{1/2} A)$ and
$P(s) \defeq W(s)^{1/2} A M(s)^{-1} A^\top W(s)^{1/2}$ the projection, so that
$\sigma_i(s) = w_i(s) a_i^\top M(s)^{-1} a_i$ and
$P(s)_{ij}^2 = w_i(s) w_j(s) (a_i^\top M(s)^{-1} a_j)^2$.
On $[m, m + 1]$, writing
$u \defeq \delta^{(m)}$, we have
$\frac{d}{ds} M(s)^{-1} = - M(s)^{-1} A^\top W(s) \diag(u) A M(s)^{-1}$, so
\[
\frac{d}{ds} \log \sigma_i(s)
=
u_i
-
\frac{\sum_j w_j(s) u_j (a_i^\top M(s)^{-1} a_j)^2}
{a_i^\top M(s)^{-1} a_i}
=
u_i
-
\frac{1}{\sigma_i(s)} \sum_j P(s)_{ij}^2 u_j .
\]
By Cauchy--Schwarz and $\sum_j P(s)_{ij}^2 = \sigma_i(s)$,
\begin{align*}
\sum_i \sigma_i(s)
\left( \frac{d}{ds} \log \sigma_i(s) \right)^2
&\le
2 \sum_i \sigma_i(s) u_i^2
+
2 \sum_i \frac{1}{\sigma_i(s)}
\left( \sum_j P(s)_{ij}^2 u_j \right)^2
\\
&\le
2 \|u\|_{\sigma(s)}^2
+
2 \sum_j u_j^2 \sum_i P(s)_{ij}^2
=
4 \|u\|_{\sigma(s)}^2 ,
\end{align*}
the last step by symmetry of $P(s)$.
On $[m, m + 1]$ the weights
$w(s)$ are within
$e^{\pm \zeta}$ of
$w^{(m + 1)}$ coordinatewise, so
$\sigma(s) \approx_{2 \zeta} \sigma^{(m + 1)}$ and, with
$D_i(s) \defeq \sigma_i(s) (\frac{d}{ds} \log \sigma_i(s))^2$,
\begin{equation}
\label{eq:lever_path_bound}
\int_m^{m + 1} \sum_i D_i(s) \, ds
\le
4 e^{2 \zeta} \|\delta^{(m)}\|_{\sigma^{(m + 1)}}^2
\le
4 e^{2 \zeta} \zeta^2 .
\end{equation}
If the pair is triggered by the leverage score, then
$|\int_t^{t'} \frac{d}{ds} \log \sigma_i(s) \, ds| \ge \eta$, and by
Cauchy--Schwarz in $s$ over an interval of length at most
$r_0$,
\[
\eta^{2} \sigma^{(t')}_i
\le
r_0 \int_t^{t'} \sigma^{(t')}_i
\left( \frac{d}{ds} \log \sigma_i(s) \right)^2 ds
\le
e^{2 \eta' + 4 \zeta} r_0 \int_t^{t'} D_i(s) \, ds ,
\]
since
$\sigma^{(t')}_i \le e^{2 \eta' + 2 \zeta} \sigma^{(m + 1)}_i \le e^{2 \eta' + 4 \zeta} \sigma_i(s)$
for $s \in [m, m + 1] \subseteq [t, t']$.

\emph{Summing.}
For a fixed $i$ the intervals
$[t_{i, j}, t_{i, j + 1})$ are disjoint, so the sums over
$m$ and the integrals over
$s$ belonging to distinct pairs of the same coordinate do not overlap.
Summing the two trigger bounds over all consecutive pairs of all coordinates,
and using $\zeta \le 0.1$,
\begin{align*}
\eta^{2} \sum_i \sum_{j \ge 1} \sigma_i^{(t_{i, j + 1})}
&\le
e^{2 \eta' + 0.4} r_0
\sum_{m = k_0}^{k_0 + r_0 - 1}
\left(
\|\delta^{(m)}\|_{\sigma^{(m + 1)}}^2
+
\int_m^{m + 1} \sum_i D_i(s) \, ds
\right)
\\
&\le
e^{2 \eta' + 0.4} r_0 \cdot r_0 \cdot 5 e^{0.2} \zeta^{2}
\le
30 e^{3 \eta'} r_0^{2} \zeta^{2} .
\end{align*}

\emph{Noisy version.}
A trigger $\what^{(t')}_i \not\approx_{\eta} \what^{(t)}_i$ together with
$\what^{(m)} \approx_{\zeta / 10} w^{(m)}$ at both ends forces
$w^{(t')}_i \not\approx_{\eta - \zeta / 5} w^{(t)}_i$, and
$\eta - \zeta / 5 \ge 0.8 \eta$ when
$\eta \ge \zeta$; conversely the second half of the trigger condition for
$\what$ gives it for
$w$ with
$\eta' + \zeta / 5 \le 1.2 \eta'$.
The leverage scores of $\what^{(m)}$ are within
$e^{\pm \zeta / 5}$ of
$\sigma^{(m)}$, so the same two implications hold for the leverage score
triggers, and the exact bound for a $(0.8 \eta, 1.2 \eta')$-triggered schedule
gives the noisy one up to a constant factor.
\end{proof}

This bounds the quantity left open in Lemma~\ref{lem:sampled_maintain}, and with
a restart schedule it proves the overall guarantee.

\begin{proof}[Proof of Theorem~\ref{thm:inv_maintain_stable}]
By Corollary~\ref{cor:adaptive} it suffices to prove the statement for an
oblivious sequence, at the cost of the $2 Q n^{-3}$ in the failure probability
and $\O(\nnz(A) \log \rho)$ per \textsc{Solve} already included in the claimed
bounds, so we may apply Lemma~\ref{lem:sampled_maintain}.

Fix a window length $r_0$ and restart the data structure, resampling from
scratch, every $r_0$ rounds.
Inside a window, a coordinate is refreshed when its weight or its estimated
leverage score has moved by more than a factor of $2$ since its last refresh,
which, the estimates being $0.1$-accurate, forces the exact weight or leverage
score to have moved by more than a factor of $2 e^{-0.2} > e^{0.1}$, while until
that round the exact quantities had moved by less than a factor of
$2 e^{0.2} < e^{1}$: the refresh schedule is
$(0.1, 1)$-triggered in the sense of Corollary~\ref{cor:update_counts}, and the
sequence handed to the structure is $\zeta$-leverage stable, or stable up to
noise with $\zeta < 0.1 = \eta$.
So the corollary bounds the leverage score mass of all coordinates refreshed in
the window, beyond the first refresh of each at the restart, by
$\O(r_0^{2} \zeta^{2})$, where the mass of a refreshed coordinate is measured at
the round of its refresh.
A refresh of coordinate $i$ at round
$t'$, its previous refresh being at
$t$, changes the sampled matrix only if the row is kept before or after the
refresh, that is, only if the coin flipped at $t$ or the coin flipped at
$t'$ came up kept; so
$\sum_k c_k$ is at most twice the number of kept coins among the coins flipped
at the refreshes of the window, plus the number of kept coins flipped at the
restart for coordinates refreshed later in the window.
Under the coupling of Lemma~\ref{lem:sampled_maintain} the refresh schedule and
the probabilities $p^{(t)}_i = \min\{1, c_s \widehat{\sigma}^{(t)}_i \log n\}$
are functions of the oblivious weights and of the ideal solver's noise, hence
independent of the coins, and conditioned on them the coins are independent.
The means of the coins counted, each at most twice, sum to at most twice
\begin{align*}
\sum_{\text{refreshes } (i, t')} p^{(t')}_i
+ \sum_{i \in I} p^{(k_0)}_i
&\le
c_s e^{0.1} \log n
\left(
\sum_{\text{refreshes } (i, t')} \sigma^{(t')}_i
+ \sum_{i \in I} \sigma^{(k_0)}_i
\right)
\\
&\le
2 c_s e^{2.3} \log n
\sum_{\text{refreshes } (i, t')} \sigma^{(t')}_i
=
\O\left( r_0^{2} \zeta^{2} \right) ,
\end{align*}
where $I$ is the set of coordinates refreshed at least once after the restart,
the refreshes summed over are those after the restart, and
$\sigma^{(k_0)}_i \le e^{2.2} \sigma^{(t_{i, 2})}_i$ for
$i \in I$ by Equation~\eqref{eq:refresh_comparable} applied to the pair
$(k_0, t_{i, 2})$ with
$2 \eta' + 2 \zeta \le 2.2$, the restart counting as the first refresh
$t_{i, 1} = k_0$ of every coordinate; the last step is the mass bound of the
corollary.
This is the expected count of~\cite[Lemma 15]{LS15}.
A Chernoff bound for independent indicators with this sum of means then gives
\begin{equation}
\label{eq:ck_whp}
\sum_{k \le r_0} c_k
=
\O\left( r_0^2 \zeta^2 + \log n \right)
\end{equation}
with probability at least $1 - n^{-10}$, and the additive
$\log n$ is absorbed below, since
$\MM(d, \log n, d) = \O(d^{\sqrt{2}} \log n)$ is inside the per window budget.
By superadditivity of $\MM(d, \cdot, d)$ in its middle dimension the
\textsc{Update} cost over the window is therefore
\[
\sum_{k \le r_0} \MM(d, c_k, d)
\le
\MM\left( d, \O(r_0^{2} \zeta^{2} + \log n), d \right)
=
\O\left( \frac{r_0^{2}}{d} \MM(d, d, d) \right),
\]
on top of the one $\O(\MM(d, d, d) \log \rho)$ paid at the restart, with
probability at least $1 - n^{-10}$ by Equation~\eqref{eq:ck_whp}, and with
probability at least $1 - n^{-8}$ in every one of the at most
$n^{2}$ windows by a union bound.
The cost bound is therefore a high probability one, as the theorem states, and
not merely a bound in expectation.
Balancing $\MM(d, d, d) / r_0$ against
$r_0 \MM(d, d, d) / d$ gives
$r_0 = d^{1/2}$ and an amortized
\[
\O\left( \frac{\MM(d, d, d)}{d^{1/2}} \log \rho \right)
=
\O\left( d^{1/2 + \sqrt{2}} \log \rho \right)
=
\O\left( d^{2} \log \rho \right)
\]
per round, the last two steps by Equation~\eqref{eq:mm_assumption}, on top of
the $\O((\nnz(A) + d^{2}) \log^{2} \rho)$ per round spent reading the update
vector and re-estimating the leverage scores for the refresh test.
The two together are the amortized
\[
\O\left( \left( \nnz(A) + d^{2} \right) \log^{2} \rho \right)
\]
of Theorem~\ref{thm:inv_maintain_stable}, the second $\log \rho$ being the
iteration count of the solves inside the estimate, the first their bit length.
Resampling at a restart re-estimates the leverage scores against the solver the
structure already has, at $\O((\nnz(A) + d^{2}) \log^{2} \rho)$, the sample of
the previous window being a constant factor preconditioner still, by
Item~\ref{itm:tau_stable:2}, so no later window repeats the dense inverse of the
first round, and the $\nnz(A) + \MM(d, d, d)$ term is paid once.
The failure probability is $n^{-8}$ per round from the proof of
Lemma~\ref{lem:sampled_maintain}, whose union bound was over rounds, so $n^{-6}$
over the $r \le n^{2}$ rounds, plus the
$2 Q n^{-3}$ of its coupling and the
$n^{-8}$ of Equation~\eqref{eq:ck_whp}, at most the
$n^{-5} + 2 Q n^{-3}$ of the theorem.
It remains to account for \textsc{Solve}, which neither
Lemma~\ref{lem:inv_maintain_explicit} nor Lemma~\ref{lem:sampled_maintain}
charges, and which Item~\ref{itm:inv_maintain_stable:3} of the theorem does.

\emph{Correctness of \textsc{Solve}.}
Let $N^{(k)} \defeq A^\top S^{(k)} W^{(k)} A$ be the matrix whose inverse the
structure holds exactly, and $A^\top W^{(k)} A$ the matrix
Equation~\eqref{eq:solve_guarantee} is stated for.
These are within a factor $e^{0.1}$ of each other by
Equation~\eqref{eq:sample_spectral}, so $(N^{(k)})^{-1}$ is a constant factor
preconditioner for $A^\top W^{(k)} A$, and preconditioned Richardson iteration
on $A^\top W^{(k)} A y = b$ contracts the error in the
$A^\top W^{(k)} A$ norm by a constant factor per step, by Imported
result~\ref{imp:perturb}.
The operator applied at each step is the exact $A^\top W^{(k)} A$, one
multiplication by $A$ and one by
$A^\top$, so the sample enters only through the contraction rate: what the
iteration converges to is the true solution, and
Equation~\eqref{eq:solve_guarantee} is a guarantee against $A^\top W^{(k)} A$
and not against the sampled matrix.
Reaching $\rho^{-10}$ takes
$\O(\log \rho)$ steps.

\emph{Cost of \textsc{Solve}.}
A step costs one multiplication by $A$, one by
$A^\top$, and one by the dense
$d \times d$ inverse, so
$\O(\nnz(A) + d^{2})$ arithmetic operations, each on numbers of
$\O(\log \rho)$ bits by Item~\ref{itm:tau_stable:1} and
Equation~\eqref{eq:sample_spectral}.
Over $\O(\log \rho)$ steps this is
\[
\O\left( \left( \nnz(A) + d^{2} \right) \log^{2} \rho \right)
\]
bit operations, which is Item~\ref{itm:inv_maintain_stable:3}; the two
logarithms are the iteration count and the bit length, and it is also what the
refresh test inside \textsc{Update} pays for its own estimate, already counted
above.

\emph{Adaptivity of \textsc{Solve}.}
Every \textsc{Solve} is served by \textsc{NoisySolver} of Imported
result~\ref{imp:noisysolver}, which is one extra solve, one extra multiplication
by $A^\top$, and the
$\O(\nnz(A) \log \rho)$ of sampling and rounding the Gaussian, so a constant
factor on the above, within Item~\ref{itm:inv_maintain_stable:3}.
This is what Corollary~\ref{cor:adaptive} charges for dropping the obliviousness
hypothesis, and it is why the theorem may be stated for an adaptive sequence.
\end{proof}

The restart balance is the only place Equation~\eqref{eq:mm_assumption} is used,
and it is used as \cite{LS15} use it, to make the square
recomputation at a restart fit under the per round budget.
Keeping the Woodbury updates rectangular does not remove it.
Suppose the structure held its inverse in factored form, as the inverse at the
last restart plus the accumulated rank one corrections, applying it to a vector
in $\O(d^{2} + C d)$ after
$C$ corrections and never forming a square product between restarts.
Every window still begins by inverting a $d \times d$ matrix, at
$\MM(d, d, d)$, and
$\O(d^{2} + C d)$ stays inside the per round budget only while
$C = \O(d)$, which by Equation~\eqref{eq:ck_whp} caps the window at
$r_0 = \O(d^{1/2} / \zeta)$.
The amortized restart cost is then $\MM(d, d, d) \zeta / d^{1/2}$, the same
balance up to the constant $\zeta$, so Equation~\eqref{eq:mm_assumption} is
needed either way.

\section{Lewis Weights and Potential}
\label{sec:lewis}

\begin{figure}[H]
\begin{center}\small
\begin{tabular}{@{}lp{0.78\textwidth}@{}}
\toprule
Symbol & Meaning \\
\midrule
$q,\alpha$ & Lewis weight parameter $q=1-\alpha\in[1/2,2)$, $\alpha=1/(10\log(10n/d))$; Lemma~\ref{lem:lewis}, Definition~\ref{def:central_weights} \\
$v$ & regularizer, $d/n\le v_i\le1$ and $\|v\|_1\le4d$; Section~\ref{subsec:weights}, Lemma~\ref{lem:lewis}, Definition~\ref{def:central_weights} \\
$\tau_q(A,v)$ & $v$-regularized $\ell_q$-Lewis weights of $A$; Section~\ref{subsec:weights}, Lemma~\ref{lem:lewis_static} \\
$\tautil$ & an estimate of $\tau_q$, supplied to Lemma~\ref{lem:lewis}; maintained along the path by Definition~\ref{def:invariant} \\
$\varsigma(w),M_w$ & Lewis map $\tau(\diag(w)^{1/2-1/q}A)+v$ and the quadratic form $A^\top\diag(w)^{1-2/q}A$ behind it; Equation~\eqref{eq:lewis_map} \\
$\varepsilon$ & target accuracy of the estimation routines, distinct from the centrality $\epsilon$ of Equation~\eqref{eq:parameters}; Lemma~\ref{lem:lewis_static}, Lemma~\ref{lem:lewis} \\
$\mathcal S,T_{\mathcal S}$ & solver supplied to Lemma~\ref{lem:lewis}, and the cost of one call; Item~\ref{itm:lewis:6} \\
$T(x),\widetilde T$ & $\diag(\tau(x))$ and $\diag(\tautil)$; Definition~\ref{def:invariant} \\
$\phi_i(t)$ & log barrier $-\log(t - l_i) - \log(u_i - t)$; $\Phi''(x) = \diag(\phi_i''(x_i))$; Lemma~\ref{lem:self_concordant} \\
$\tau(x)$ & $\tau_q(\Phi''(x)^{-1/2} A, v)$; Definition~\ref{def:central_weights} \\
$\mu$ & path parameter, from $\mu_0$ down to $\mu_f$; Definition~\ref{def:centered}, Lemma~\ref{lem:pathfollowing} \\
$\lambda$ & potential scale, $\Theta(\log(n / \epsilon) / \epsilon)$; Equation~\eqref{eq:parameters}, Definition~\ref{def:potential} \\
$y(x,s,\mu),\overline y$ & centrality vector of Definition~\ref{def:potential}, and its computed value in a step, Lemma~\ref{lem:ybar} \\
$\Psi(x,s,\mu)$ & centrality potential; Definition~\ref{def:potential}, Lemma~\ref{lem:progress} \\
$C,C_{\mathrm{start}}$ & universal constant of Equation~\eqref{eq:parameters}, and $10C$; Lemma~\ref{lem:lewis_approx} \\
$\gamma,\Cnorm$ & step size and norm scale of Equation~\eqref{eq:parameters}; Definition~\ref{def:centered}, Definition~\ref{def:taunorm} \\
$\|g\|_{\tau+\infty}$ & $\|g\|_\infty + \Cnorm\|g\|_\tau$, where $\|g\|_\tau^2 = \sum_i\tau_ig_i^2$; Definition~\ref{def:taunorm} \\
$g^{\flat(\tau)}$ & the dual unit vector of $g$ in $\|\cdot\|_{\tau+\infty}$; Definition~\ref{def:taunorm}, Definition~\ref{def:admissible} \\
\bottomrule
\end{tabular}
\end{center}
\caption{Variables local to the Lewis weights and potential of
Section~\ref{sec:lewis}.}
\label{fig:notation_lewis}
\end{figure}

This section collects everything about the weights of the central path: the
regularized Lewis weights and the two routines that estimate them, the barrier
and the parameters of the method, and the centrality potential that
Section~\ref{sec:potential_and_progress} drives down one step at a time.
All of it is the fixed point analogue of the regularized Lewis weight central
path of~\cite{BLLSSSW21}, transposed to our
dimensions ($n$ rows,
$d$ columns, so
$m \rightarrow n$ and
$n \rightarrow d$ there), and with the
$\ell_{p}$ exponent there written as
$q$ here.
The data structure of Section~\ref{sec:dynamic} knows nothing of these weights:
it sees only the vector handed to \textsc{Update}.
The estimation routines below read $A$ only through matrix vector products and
through calls of a supplied solver, and the interior point method supplies
\textsc{Solve} of Definition~\ref{def:operations} as that solver.

\subsection{Regularized Lewis weights}
\label{subsec:weights}

\begin{definition}
For $q\in[1/2,2)$,
$n \times d$ matrix
$A$ with
$n > d$ and rank
$d$, and length
$n$ positive vector
$v$ with
$d/n \le v_i \le 1$ and
$\|v\|_1\le4d$, the
$v$-regularized
$\ell_{q}$-Lewis weights of
$A$,
$\tau_{q}(A, v)$ is the unique vector such that
\[
\tau_{q}(A, v)
=
\tau\left(\diag\left(\tau_{q}(A, v)\right)^{1/2-1/q}A\right) + v
\]
\end{definition}

The fixed point that defines the weights is reached by a damped iteration that
touches $A$ only through matrix vector products and through linear systems in a
quadratic form spectrally close to the Lewis one.
We state it twice.
Once from scratch, where the starting point is the all ones vector and the
solver has to be built as well, which is the only place in the paper a dense
$d \times d$ inverse is formed.
And once from a coarse estimate $\tautil$ of the answer together with a solver
for the quadratic form that estimate induces, which is the form the interior
point method calls it in, the estimate being the previous round's and the solver
being the \textsc{Solve} of Definition~\ref{def:operations}.
Throughout this subsection write
\begin{equation}
\label{eq:lewis_map}
\varsigma(w)
\defeq
\tau\left( \diag(w)^{1/2 - 1/q} A \right) + v
\end{equation}
for the map whose fixed point is $\tau_{q}(A, v)$, and
$M_{w} \defeq A^{\top} \diag(w)^{1 - 2/q} A$ for the quadratic form behind it,
so that $\varsigma_i(w) - v_i = w_i^{1 - 2/q} a_i^{\top} M_{w}^{-1} a_i$.

The two facts about Equation~\eqref{eq:lewis_map} that both routines rest on are
Lemma 5.3 and Lemma 5.5 of~\cite{BLLSSSW21},
which transfer the contraction argument of~\cite{CP15} to the
regularized weights; the regularizer $v$ is their
$z$.

\begin{importedresult}
\label{imp:lewis_contraction}
Let $q \in (0, 2)$ and
$w, \overline{\varsigma} \in \R^{n}_{> 0}$ with
$w \approx_{\delta} \varsigma(w)$ and
$\overline{\varsigma} \approx_{\eta} \varsigma(w)$, and let
$w' \defeq (w^{2/q - 1} \overline{\varsigma})^{q/2}$ coordinatewise.
Then $w' \approx_{(\delta + \eta)|1 - q/2| + \eta} \varsigma(w')$ and
$w' \approx_{(\delta + \eta) q/2} w$.
Moreover any $w \approx_{\delta} \varsigma(w)$ satisfies
$w \approx_{\delta} \tau_{q}(A, v)$.
\end{importedresult}

The static Lewis estimation is the following.

\begin{lemma}
\label{lem:lewis_static}
For $n > d$, given:
\begin{enumerate}
\item \label{itm:lewis_static:1} an $n \times d$ matrix
  $A$ with column rank
  $d$ and
  $A^{\top} A$ has magnitude at most
  $\rho$,
\item \label{itm:lewis_static:2} parameter $1/2 \leq q < 2$,
\item \label{itm:lewis_static:3} vector $v$ with
  $d/n \le v_i \le 1$ and
  $\|v\|_1\le 4d$,
\item \label{itm:lewis_static:4} accuracy $0.1 > \varepsilon > \rho^{-1}$,
\end{enumerate}
we can return an $\varepsilon$-approximation to the
$v$-regularized
$\ell_{q}$-Lewis weights of
$A$, with probability at least
$1 - n^{-7}$, in bit complexity
\[
\O\left( \left( \varepsilon^{-2} \left( \nnz(A) + d^{2} \right)
+ \MM(d, d, d) \right) \log \rho \, \log^{2}(1 / \varepsilon) \right) .
\]
\end{lemma}

\begin{proof}
Set $\eta \defeq \varepsilon / 20$ and run
\begin{equation}
\label{eq:lewis_iterate}
w^{(t + 1)}
\defeq
\Trunc\left(
\left( \left( w^{(t)} \right)^{2/q - 1} \overline{\varsigma}^{(t)} \right)^{q/2},
\, 20 \log \rho \right),
\qquad
\overline{\varsigma}^{(t)} \approx_{\eta} \varsigma\left( w^{(t)} \right),
\end{equation}
coordinatewise from $w^{(0)} \defeq \mathbf{1}$ for
$L \defeq \lceil 4 \log(2 \log n / \varepsilon) \rceil$ steps, where
$\overline{\varsigma}^{(t)}$ is the estimate of
$\varsigma(w^{(t)})$ computed below, and return
$w^{(L)}$.

\paragraph{Convergence.}
Let $\delta_t$ be such that
$w^{(t)} \approx_{\delta_t} \varsigma(w^{(t)})$.
Since leverage scores lie in $[0, 1]$ and
$v_i \in [d/n, 1]$,
$\varsigma(\mathbf{1}) \in [d/n, 2]^{n}$, so
$\delta_0 \le \log n$.
As long as $\delta_t \le \log n$ the iterate satisfies
$w^{(t)} \approx_{\log n} \tau_{q}(A, v) \in [d/n, 2]^{n}$ by the last part of
Imported result~\ref{imp:lewis_contraction}, so
$w^{(t)} \in [\rho^{-2}, \rho]^{n}$, the untruncated update has magnitude
$\rho^{O(1)}$, and the truncation in Equation~\eqref{eq:lewis_iterate} changes
each coordinate by a factor $e^{\pm \rho^{-17}}$, which we absorb into
$\eta$ (it is below
$\eta / 100$ as
$\varepsilon > \rho^{-1}$).
The first part of Imported result~\ref{imp:lewis_contraction} with
$|1 - q/2| \le 3/4$ gives
\[
\delta_{t + 1}
\le
\frac{3}{4} \left( \delta_t + \eta \right) + \eta
\le
\frac{3}{4} \delta_t + 2 \eta ,
\qquad\text{so}\qquad
\delta_t
\le
\left( \frac{3}{4} \right)^{t} \log n + 8 \eta
\le
\left( \frac{3}{4} \right)^{t} \log n + \frac{\varepsilon}{2}
\]
for all $t$, which in particular keeps
$\delta_t \le \log n$ throughout.
At $t = L$ the first term is at most
$\varepsilon / 2$, so
$w^{(L)} \approx_{\varepsilon} \varsigma(w^{(L)})$, and the last part of
Imported result~\ref{imp:lewis_contraction} turns this into
$w^{(L)} \approx_{\varepsilon} \tau_{q}(A, v)$.
Since $\varepsilon < 0.1$,
$L = O(\log(1 / \varepsilon) + \log \log n) = \O(\log(1 / \varepsilon))$.

\paragraph{Estimating $\varsigma(w^{(t)})$.}
Write $D \defeq \diag(w^{(t)})^{1/2 - 1/q}$ and
$P \defeq D A M_{w^{(t)}}^{-1} A^{\top} D$ for the projection onto the column
span of $D A$, so
$\varsigma_i(w^{(t)}) - v_i = \|P e_i\|_2^{2}$.
Let $G \in \R^{k \times n}$ have independent
$N(0, 1/k)$ entries with
$k \defeq \lceil 10^{4} \eta^{-2} \log n \rceil$, and set
$\overline{\varsigma}^{(t)}_i \defeq v_i + \|G P e_i\|_2^{2}$.
By the Johnson--Lindenstrauss lemma in its Gaussian form, for each fixed $i$ the
event $\|G P e_i\|_2^{2} \approx_{\eta / 2} \|P e_i\|_2^{2}$ fails with
probability at most $2 e^{-k \eta^{2} / 100} \le n^{-10}$, so all
$n$ coordinates satisfy it with probability at least
$1 - n^{-9}$, and then
$\overline{\varsigma}^{(t)} \approx_{\eta / 2} \varsigma(w^{(t)})$ since adding
$v_i \ge 0$ to both sides can only shrink a ratio.
We also use that $\|G\|_F^{2} \le 2 n$ except with probability
$e^{-k n / 8} \le n^{-9}$:
$k \|G\|_F^{2}$ is a chi-squared variable with
$k n$ degrees of freedom, and
$\Pr[\chi^{2}_{m} \ge 2 m] \le e^{-m / 8}$.
The $k$ rows of
$G P$ are
$(P g_j)^{\top}$ for the rows
$g_j$ of
$G$, and
$P g_j = D A \left( M_{w^{(t)}}^{-1} \left( A^{\top} D g_j \right) \right)$ is
two matrix vector products with $D A$ at
$O(\nnz(A))$ bit operations and one linear solve in
$M_{w^{(t)}}$.
Summing squares of the $k$ rows costs
$O(k n)$, so the estimate costs
$O(k \nnz(A))$ bit operations plus
$k$ solves.

\paragraph{The solves and their accuracy.}
Each solve is run by preconditioned Richardson iteration to relative accuracy
\[
\delta
\defeq
\frac{\eta}{100 n^{2}}
\]
in the $M_{w^{(t)}}$ norm.
The preconditioner is rebuilt at every iterate: apply a sparse embedding $\Pi$
with $\O(d)$ rows and constant factor spectral guarantee to
$D A$ at
$\O(\nnz(A))$ bit operations, form the Gram matrix of
$\Pi D A$ at
$\O(\MM(d, d, d))$ and invert it with Imported result~\ref{imp:denseinverse} at
$\O(\MM(d, d, d) \log \rho)$; the inverse is a constant factor spectral
approximation of $M_{w^{(t)}}^{-1}$, so
$O(\log(1 / \delta)) = O(\log(n / \eta))$ Richardson steps, each one application
of $M_{w^{(t)}}$ at
$O(\nnz(A) \log \rho)$ bit operations and of the dense inverse at
$O(d^{2} \log \rho)$, reach accuracy
$\delta$.
In fixed point the iteration is run exactly as the \textsc{Solve} in the proof
of Lemma~\ref{lem:sampled_maintain}: the right hand side is scaled by a power of
two into $[1/2, 1)$, the numbers have
$20 \log \rho$ bits, and every matrix vector product is formed exactly and
truncated once, so the argument there bounds the relative error of one step by
$2 \rho^{-14.5}$, which is below
$\delta \ge \rho^{-1} / (2000 \rho^{2})$ by a factor of more than
$\rho^{10}$ and is absorbed into the rate.
The residual error is negligible: the computed $j$-th row of
$G P$ is
$(D A \widetilde{y}_j)^{\top}$ with
$\|\widetilde{y}_j - M_{w^{(t)}}^{-1} A^\top D g_j\|_{M_{w^{(t)}}} \le \delta \|P g_j\|_2 \le \delta \|g_j\|_2$,
since $\|y\|_{M_{w^{(t)}}} = \|D A y\|_2$ and
$P$ is a projection; so the row is
$(P g_j)^{\top} + f_j^{\top}$ with
$\|f_j\|_2 \le \delta \|g_j\|_2$.
Writing $a_j \defeq (P g_j)_i$ and using
$\sum_j a_j^{2} \le e^{\eta / 2} \|P e_i\|_2^{2} \le 2$ on the
Johnson--Lindenstrauss event and
$\sum_j (f_j)_i^{2} \le \delta^{2} \|G\|_F^{2} \le 2 n \delta^{2}$,
\[
\begin{aligned}
\left| \sum_j \left( a_j + (f_j)_i \right)^{2} - \sum_j a_j^{2} \right|
&\le
2 \left( \sum_j a_j^{2} \right)^{1/2} \left( \sum_j (f_j)_i^{2} \right)^{1/2}
+ \sum_j (f_j)_i^{2} \\
&\le
2 \cdot 2^{1/2} \cdot (2 n)^{1/2} \delta + 2 n \delta^{2} \\
&\le
6 n^{1/2} \delta
=
\frac{6 \eta}{100 n^{3/2}}
\le
\frac{\eta}{4} \cdot \frac{d}{n}
\le
\frac{\eta}{4} \varsigma_i\left( w^{(t)} \right) ,
\end{aligned}
\]
using $\delta \le 1 / n$ in the third step and
$\varsigma_i(w^{(t)}) \ge v_i \ge d / n$ in the last.
Together with the Johnson--Lindenstrauss error this gives
$\overline{\varsigma}^{(t)} \approx_{\eta} \varsigma(w^{(t)})$, as
Equation~\eqref{eq:lewis_iterate} requires.

\paragraph{Total.}
Per iterate: one preconditioner at $\O(\MM(d, d, d) \log \rho)$, and
$k = \O(\varepsilon^{-2})$ solves at
$O((\nnz(A) + d^{2}) \log \rho \, \log(n / \varepsilon))$ each, plus
$O(k \nnz(A) \log \rho)$ around them.
Over the $L = O(\log(1 / \varepsilon))$ iterates this is
\[
\O\left( \left( \varepsilon^{-2} \left( \nnz(A) + d^{2} \right)
\log(n / \varepsilon) + \MM(d, d, d) \right) \log \rho \, \log(1 / \varepsilon) \right) ,
\]
which is the stated bound since
$\log(n / \varepsilon) \le \log n \cdot \log(1 / \varepsilon)$ for
$\varepsilon < 0.1$ and
$\O$ hides the
$\log n$.
The failure probability is at most $L (n^{-9} + n^{-9}) \le n^{-7}$, one
Johnson--Lindenstrauss event and one Frobenius norm event per iterate.
\end{proof}

Inside the interior point method the estimate of the previous round is already
at hand, and with it a solver, so neither the all ones start nor the dense
inverse is needed.

The fixed-point Lewis estimation from a solver is the following.

\begin{lemma}
\label{lem:lewis}
For $n > d$, given:
\begin{enumerate}
\item \label{itm:lewis:1} an $n \times d$ matrix
  $A$ with column rank
  $d$ and
  $A^{\top} A$ has magnitude at most
  $\rho$,
\item \label{itm:lewis:2} parameter $1/2 \leq q < 2$,
\item \label{itm:lewis:3} vector $v$ with
  $d/n \le v_i \le 1$ and
  $\|v\|_1\le 4d$,
\item \label{itm:lewis:4} accuracy $0.1 > \varepsilon > \rho^{-1}$,
\item \label{itm:lewis:5} a vector $\tautil$ that
  $0.5$-approximates the weights being estimated,
  $\tautil \approx_{0.5} \tau_{q}(A, v)$,
\item \label{itm:lewis:6} a solver $\mathcal{S}$ for a fixed
  $d \times d$ matrix
  $M$ with
  $M \approx_{0.5} A^{\top} \diag(\tautil)^{1 - 2/q} A$, which on input
  $b \in \R^{d}$ returns
  $y$ with
  $\|y - M^{-1} b\|_{M} \le \rho^{-10} \|M^{-1} b\|_{M}$ in
  $T_{\mathcal{S}}$ bit operations,
\end{enumerate}
we can return an $\varepsilon$-approximation to the
$v$-regularized
$\ell_{q}$-Lewis weights of
$A$, with probability at least
$1 - n^{-7}$, in bit complexity
\[
\O\left( \varepsilon^{-2} \left( \nnz(A) \log \rho + T_{\mathcal{S}} \right)
\log^{2}(1 / \varepsilon) \right) .
\]
\end{lemma}

\begin{proof}
Run the iteration of Equation~\eqref{eq:lewis_iterate} from
$w^{(0)} \defeq \tautil$ for
$L \defeq \lceil 4 \log(4 / \varepsilon) \rceil = O(\log(1 / \varepsilon))$
steps, with the solves in $M_{w^{(t)}}$ preconditioned by
$\mathcal{S}$ instead of by a rebuilt dense inverse.

\paragraph{Convergence.}
By Item~\ref{itm:lewis:5},
$\varsigma(\tautil) \approx_{0.5 |1 - 2/q|} \varsigma(\tau_{q}(A, v)) = \tau_{q}(A, v) \approx_{0.5} \tautil$,
where the first step raises the coordinatewise ratio $\tautil / \tau_{q}(A, v)$
to the power $1 - 2/q$ inside the quadratic form and the leverage scores of
$D A$ change by at most the factor
$D$ changes.
As $|1 - 2/q| \le 3$ on
$[1/2, 2)$ this gives
$\delta_0 \le 2$, and the recursion of the proof of Lemma~\ref{lem:lewis_static}
gives $\delta_t \le (3/4)^{t} \cdot 2 + \varepsilon / 2 \le 2$ for all
$t$ and
$\delta_L \le \varepsilon$, so
$w^{(L)} \approx_{\varepsilon} \tau_{q}(A, v)$ by Imported
result~\ref{imp:lewis_contraction}.
The magnitude bookkeeping is as there, now with
$w^{(t)} \approx_{2} \tau_{q}(A, v)$.

\paragraph{The one fixed preconditioner suffices.}
Since every iterate satisfies
$w^{(t)} \approx_{2} \tau_{q}(A, v) \approx_{0.5} \tautil$,
\[
M_{w^{(t)}}
\approx_{2.5 |1 - 2/q|}
A^{\top} \diag(\tautil)^{1 - 2/q} A
\approx_{0.5}
M ,
\]
the second step being Item~\ref{itm:lewis:6}.
So $M$ is an
$e^{8}$ spectral approximation of
$M_{w^{(t)}}$ at every iterate, and a constant factor preconditioner is all
Richardson iteration needs: $O(\log(n / \eta))$ steps, each one application of
$M_{w^{(t)}}$ at
$O(\nnz(A) \log \rho)$ bit operations, its result truncated to the
$20 \log \rho$ bits the solver takes, and one call to
$\mathcal{S}$ at
$T_{\mathcal{S}}$, reach the relative accuracy
$\delta = \eta / (100 n^{2})$ of the proof of Lemma~\ref{lem:lewis_static}, the
fixed point bookkeeping being the one there: the residual is scaled into
$[1/2, 1)$ before it is handed to
$\mathcal{S}$, whose relative error
$\rho^{-10}$ and the truncation of the residual, of relative size at most
$2 \rho^{-14.5}$ as in the proof of Lemma~\ref{lem:sampled_maintain}, are both
far below $\delta \ge \rho^{-3} / 2000$ and are absorbed into the rate.
At the value $q = 1 - \alpha$ of Definition~\ref{def:central_weights} the
exponent is $|1 - 2/q| \le 1.1$ and the constant is below
$e^{3.3}$.
Only the rate at which Richardson contracts depends on $M$, not the point it
converges to, so $M$ enters as a preconditioner only, and both approximations
are asked at the loose accuracy $0.5$; stating Item~\ref{itm:lewis:6} at the
supplied $\tautil$ rather than at the answer is what lets the caller check it.

\paragraph{Total.}
The estimate of $\varsigma(w^{(t)})$, its accuracy, and its failure probability
are exactly as in the proof of Lemma~\ref{lem:lewis_static}:
$k = \O(\varepsilon^{-2})$ solves and
$O(k \nnz(A) \log \rho)$ bit operations per iterate, each solve
$O(\log(n / \varepsilon))$ Richardson steps at
$O(\nnz(A) \log \rho + T_{\mathcal{S}})$, over
$L = O(\log(1 / \varepsilon))$ iterates, which is
$\O(\varepsilon^{-2} (\nnz(A) \log \rho + T_{\mathcal{S}}) \log^{2}(1 / \varepsilon))$
as $\log(n / \varepsilon) \le \log n \cdot \log(1 / \varepsilon)$; the failure
probability is $L (n^{-9} + n^{-9}) \le n^{-7}$.
\end{proof}

\subsection{Barrier, weights, and parameters}

The barrier of the $i$-th box constraint
$l_i \le x_i \le u_i$ is
\[
\phi_i\left( x_i \right)
\defeq
-\log(x_i - l_i) - \log(u_i - x_i),
\]
and $\Phi''(x) \defeq \diag(\phi_i''(x_i))$.

The barrier is highly self-concordant in the sense of Definition 4.1 of
\cite{BLLSSSW21}; the following is Lemma 4.2
there.

\begin{lemma}
\label{lem:self_concordant}
Each $\phi_i$ is highly
$1$-self-concordant on
$(l_i, u_i)$: for all
$t\in(l_i,u_i)$,
\[
|\phi_i'(t)| \le \phi_i''(t)^{1/2},
\quad
|\phi_i'''(t)| \le 2 \phi_i''(t)^{3/2},
\quad
|\phi_i''''(t)| \le 6 \phi_i''(t)^{2}.
\]
Moreover $\phi_i'((l_i + u_i)/2) = 0$ and
$\phi_i''(t) \ge (u_i - l_i)^{-2}$.
\end{lemma}

The fourth derivative bound is what makes the second order expansion of the
weights in Lemma~\ref{lem:step_bounds} valid, and is the reason we cannot work
with an arbitrary $1$-self-concordant barrier.
One consequence is used repeatedly: if
$|\phi_i''(t)^{1/2} (t' - t)| \le \beta \le 1/100$ then, since the two distances
to the bounds change by factors in $[1 - \beta, 1 + \beta]$,
\begin{equation}
\label{eq:barrier_local}
\phi_i''(t') \approx_{3 \beta} \phi_i''(t),
\qquad
\left| \phi_i'(t') - \phi_i'(t) \right| \le 2 \beta \, \phi_i''(t)^{1/2} .
\end{equation}

The central path weights are defined as follows.

\begin{definition}
\label{def:central_weights}
Fix the Lewis exponent and the regularizer
\[
q \defeq 1 - \alpha,
\qquad
\alpha \defeq \frac{1}{10 \log(10 n / d)},
\qquad
v \in \R^{n}_{>0}
\text{ with }
\frac{d}{n} \le v_i \le 1
\text{ and }
\|v\|_1 \le 4 d,
\]
and define the central path weights at a point $x$ with
$l < x < u$ by
\[
\tau(x)
\defeq
\tau_{q}\left( \Phi''(x)^{-1/2} A, v \right),
\qquad
T(x) \defeq \diag(\tau(x)).
\]
\end{definition}

Since $\log(10 n / d) \ge \log 10$ we have
$\alpha \le 1/23$, so
$q \in [1/2, 2)$ as required by Lemma~\ref{lem:lewis}, and the weights are
bounded on both sides: $d / n \le v_i \le \tau_i(x) \le 1 + v_i \le 2$ for every
$i$, and
$\|\tau(x)\|_1 = d + \|v\|_1 \le 5 d$.

The weights are stable under rescaling; the following is Lemma 4.22 of
\cite{BLLSSSW21}.

\begin{lemma}
\label{lem:lewis_approx}
For $q \in (0, 4)$ and positive diagonal
$C \approx_{\epsilon} \widetilde{C}$ we have
$\tau_{q}(C A, v) \approx_{4 \epsilon} \tau_{q}(\widetilde{C} A, v)$.
\end{lemma}

Throughout, $C$ denotes a sufficiently large universal constant and the
algorithm's parameters are
\begin{equation}
\label{eq:parameters}
\begin{aligned}
\epsilon &\defeq \frac{\alpha}{C}, \\
\lambda &\defeq \frac{C \log(C n / \epsilon^2)}{\epsilon}, \\
\gamma &\defeq \frac{\epsilon}{C \lambda}, \\
\Cnorm &\defeq \frac{C}{1 - q} = \frac{C}{\alpha}, \\
\theta &\defeq \frac{\epsilon \gamma}{\Cnorm d^{1/2}}, \\
C_{\mathrm{start}} &\defeq 10 C,
\end{aligned}
\end{equation}
so that $\theta^{-1} = \O(d^{1/2})$, as used in Figure~\ref{fig:step}.
The explicit requirements on $C$ that this section uses are the following.
From $\alpha \le 1/23$,
$\epsilon \le 1/80$ once
$C \ge 4$, as Definition~\ref{def:centered} requires.
Next $e^{\lambda \epsilon} = (C n / \epsilon^2)^{C} \ge n^{10}$ once
$C \ge 10$, which is what converts a polynomial bound on the potential into
$\epsilon$-centrality in Lemma~\ref{lem:pathfollowing}.
And $\lambda \epsilon \le 3 C \log n$ for
$n$ larger than a constant depending on
$C$, which with
$C_{\mathrm{start}} = 10 C$ keeps the initial potential below
$n^{2}$ in Lemma~\ref{lem:pathfollowing}.
Both $\gamma^{-1} = C^{4} \alpha^{-2} \log(C n / \epsilon^2) = O(\log^{3} n)$
and $\Cnorm = O(\log n)$ are
$\O(1)$.
The remaining constraints on $C$ are the hidden constants of the imported
Lemmas~\ref{lem:step_bounds} to~\ref{lem:progress}, which
\cite{BLLSSSW21} asserts for $C$ sufficiently
large without a value.

The algorithm uses, for each of the parameters of Equation~\eqref{eq:parameters}
and for $\alpha$, a number with
$40 \log \rho$ bits within a factor
$1 \pm \rho^{-30}$ of the value displayed; each requirement above is an
inequality with room for such a factor, and we do not distinguish the two in the
notation.

The magnitude bound $\rho$ is, throughout this section, a common bound for the
magnitudes of $A^\top A$ and of every
$A^\top W A$ the data structure holds, and also for
$\phi_i''(x_i)^{\pm 1}$,
$\mu^{\pm 1}$,
$\|A\|$,
$\|c\|_2$,
$\|l\|_{\infty}$ and
$\|u\|_{\infty}$ along the path; we take
$\rho \ge n$ as well, which together with
$\rho = 2^{O(B)}$ is the standing assumption
$\log n = O(B)$ on the bit length of the input.
Lemma~\ref{lem:magnitudes} gives $\rho = 2^{O(B)}$, and enlarging
$\rho$ affects only the
$\log \rho$ factors of the costs.
All arithmetic inside a step is carried out on numbers with
$40 \log \rho = O(B)$ bits, every intermediate result being truncated by
$\Trunc(\cdot, 40 \log \rho)$, except that a right hand side handed to
\textsc{Solve} is truncated to the $20 \log \rho$ bits of
Definition~\ref{def:operations}.
Every quantity a step forms, other than the calls to \textsc{Solve}, is a
formula in the sense of Lemma~\ref{lem:roundoff} whose inputs are stored numbers
or outputs of \textsc{Solve}, whose inputs and intermediate results have
magnitude at most $\rho$, and whose divisors and radicands are at least
$\rho^{-1}$, by Lemma~\ref{lem:magnitudes} with
$\rho$ enlarged by a constant power; so
$M = \rho$, and the requirements
$M \le 2^{\ell - 1}$ and on the bits before the decimal point are met.
The formulas are of two kinds: an entrywise formula with at most $30$ operations
and multiplicative depth at most $6$, and a sum over the rows of at most
$3 n$ products of two stored numbers, of multiplicative depth
$1$ and with at most
$6 n \le 6 \rho$ operations, such as an entry of
$A^\top v$ or of
$A h$; an entrywise formula may take sums of the second kind as inputs.
Lemma~\ref{lem:roundoff} bounds the error of a sum of the second kind by
$6 \rho \cdot 2 \rho^{2} \cdot \rho^{-40} \le \rho^{-36}$, and then that of an
entrywise formula with at most $30$ such inputs by
\begin{equation}
\label{eq:step_roundoff}
\left( 30 \rho^{-40} + 30 \rho^{-36} \right) \left( 2 \rho^{2} \right)^{6}
\le
\rho^{-20}
\end{equation}
for $\rho \ge 4000$, which is below the
$1 / (2 \rho)$ that the lemma requires.
A sum of the second kind whose factors are themselves entrywise formulas, such
as an entry of $A^\top \Phi''(x)^{-1/2} g$, has inputs off by at most
$\rho^{-20}$ each, and Lemma~\ref{lem:roundoff} bounds its error by
\begin{equation}
\label{eq:sum_roundoff}
\left( 6 \rho \cdot \rho^{-40} + 6 n \rho^{-20} \right) 2 \rho^{2}
\le
\rho^{-16} .
\end{equation}
Every truncation error quoted in Sections~\ref{sec:potential_and_progress}
and~\ref{sec:overall} is an instance of Equation~\eqref{eq:step_roundoff}
or~\eqref{eq:sum_roundoff}, or a direct application of Lemma~\ref{lem:roundoff}
with $M = \rho$, and the proof states the operation count and the depth of the
formula concerned.

\subsection{Centrality and the potential}

An $\epsilon$-centered point is defined as follows.

\begin{definition}
\label{def:centered}
A triple $(x, s, \mu)$ with
$l < x < u$ and
$\mu > 0$ is
$\epsilon$-centered for
$\epsilon \in (0, 1/80]$ if
\begin{enumerate}
\item \label{itm:centered:1} (Approximate centrality)
  $\displaystyle \| \frac{s + \mu T(x) \phi'(x)}{\mu T(x) \Phi''(x)^{1/2}} \|_{\infty} \le \epsilon$,
\item \label{itm:centered:2} (Dual feasibility) $s = A z + c$ for some
  $z \in \R^{d}$,
\item \label{itm:centered:3} (Approximate primal feasibility)
  $\displaystyle \| A^\top x - b \| _{( A^\top (T(x) \Phi''(x))^{-1} A )^{-1}} \le \frac{\epsilon \gamma}{\Cnorm}$.
\end{enumerate}
\end{definition}

The algorithm stores the multiplier $z$ of Item~\ref{itm:centered:2} rather than
$s$; wherever
$s$ appears below it stands for
$A z + c$, evaluated in
$O(\nnz(A))$ operations when needed.
This makes dual feasibility hold by construction, truncation included.

The centrality potential is the following.

\begin{definition}
\label{def:potential}
For $\psi(t) \defeq \cosh(\lambda t)$ with
$\lambda$ as in Equation~\eqref{eq:parameters}, and the centrality vector
\[
y_i(x, s, \mu)
\defeq
\frac{s_i + \mu \tau_i(x) \phi_i'(x_i)}
{\mu \tau_i(x) \phi_i''(x_i)^{1/2}},
\qquad
\Psi(x, s, \mu)
\defeq
\sum_{i = 1}^{n}
\psi\left( y_i(x, s, \mu) \right).
\]
\end{definition}

Exact centrality is $s + \mu T(x) \phi'(x) = 0$; the denominators rescale each
coordinate by the local Hessian.
Since $\cosh(\lambda y_i) \le \Psi$, a bound
$\Psi \le n^{10}$ gives
$|y_i| \le \log(2 n^{10}) / \lambda \le \epsilon$ for every
$i$ by the choice
$e^{\lambda \epsilon} \ge n^{10}$ in Equation~\eqref{eq:parameters}, which is
Item~\ref{itm:centered:1}.

The $\tau + \infty$ norm is the following.

\begin{definition}
\label{def:taunorm}
For $g \in \R^{n}$ and a weight vector
$\tau$,
\[
\|g\|_{\tau + \infty}
\defeq
\|g\|_{\infty} + \Cnorm \|g\|_{\tau},
\qquad
\|g\|_{\tau} \defeq \left( \sum_i \tau_i g_i^2 \right)^{1/2},
\]
with dual norm $\|g\|^{*}_{\tau + \infty}$ and
$g^{\flat(\tau)} \defeq \arg\max_{\|h\|_{\tau + \infty} = 1} h^\top g$.
\end{definition}

The step direction is measured in $\|\cdot\|_{\tau + \infty}$ because the
$\ell_\infty$ part keeps every coordinate inside the region where the barrier
Hessian is stable, while the $\tau$ part is what controls the movement of the
Lewis weights, Lemma~\ref{lem:tau_move}, and hence the work of the data
structure.
Two elementary comparisons between the norms are used below: for every $g$,
\begin{equation}
\label{eq:norm_compare}
\|g\|_{\tau} \le \|\tau\|_1^{1/2} \|g\|_{\infty} \le (5 d)^{1/2} \|g\|_{\infty},
\qquad
\|g\|_{\tau + \infty} \le \left( 1 + \Cnorm (5 d)^{1/2} \right) \|g\|_{\infty},
\end{equation}
and if $\tau \approx_{\delta} \tau'$ then
$\|g\|_{\tau + \infty} \le e^{\delta / 2} \|g\|_{\tau' + \infty}$.

The algorithm stores $x$ exactly, as a vector of numbers with
$40 \log \rho$ bits, but never
$\tau(x)$ exactly; it stores the fixed point approximation of
Lemma~\ref{lem:lewis}, run against the \textsc{Solve} call of
Definition~\ref{def:operations}.

The approximations maintained across steps are the following.

\begin{definition}
\label{def:invariant}
At the start of every call to \textsc{Step} we hold a point $\overline{x}$ and a
vector $\tautil$ with
\[
\left\| \Phi''(x)^{1/2} (\overline{x} - x) \right\|_{\infty} \le \epsilon,
\qquad
\tautil \approx_{\epsilon} \tau(\overline{x}),
\]
and write $\widetilde{T} \defeq \diag(\tautil)$.
\end{definition}

This is Invariant 4.10 of~\cite{BLLSSSW21}.
There the proxy $\overline{x}$ is needed because the iterate lives implicitly
and only $\O(d)$ of its coordinates may be read per step; here every coordinate
of $x$ is explicit, so we take
$\overline{x} \defeq x$ throughout, the first condition holds with slack zero,
and $\overline{x}$ is kept in the statements only so that they read as in the
source.
The estimate $\tautil$ is produced by Lemma~\ref{lem:lewis} at accuracy
$\gamma / 100 \le \epsilon$ (step 2 of Figure~\ref{fig:step}); the finer
accuracy is what step 4 needs, see Lemma~\ref{lem:ybar}.

\section{Effect of One Step}
\label{sec:potential_and_progress}

\begin{figure}[H]
\begin{center}\small
\begin{tabular}{@{}lp{0.78\textwidth}@{}}
\toprule
Symbol & Meaning \\
\midrule
$\xini$ & initial point given with the LP data, $l+2^{-B}<\xini<u-2^{-B}$; Theorem~\ref{thm:main}, Lemma~\ref{lem:initialization} \\
$\overline x$ & proxy for $x$ inside a step, Definition~\ref{def:invariant}; equal to $x$ here, Figure~\ref{fig:step} \\
$x,sl,su$ & primal iterate and dual slacks, $sl - su - c$ in the range of $A$; Definition~\ref{def:centered}, Lemma~\ref{lem:initialization} \\
$z,s$ & stored dual multiplier and the slack it defines, $s = Az + c$; Item~\ref{itm:centered:2}, Figure~\ref{fig:step} \\
$\theta$ & path step size, $\theta^{-1} = \O(d^{1/2})$; Figure~\ref{fig:step}, Lemma~\ref{lem:progress}, Lemma~\ref{lem:pathfollowing} \\
$\overline y, g$ & centrality proxy of step 4 of Figure~\ref{fig:step}, Lemma~\ref{lem:ybar}, and the admissible direction, Definition~\ref{def:admissible} \\
$\ghat, \pi, s_k, h^{(k)}, \nu$ & inside \textsc{Direction}, Figure~\ref{fig:direction}: computed gradient, its normalization, scale, candidate, certified norm \\
$\widetilde{q}_k, \widetilde{\nu}_k$ & inside \textsc{Direction}: values computed on numbers with $40 \log \rho$ bits; proof of Lemma~\ref{lem:direction} \\
$\overline A$ & rescaled matrix $\widetilde T^{-1/2}\Phi''(\overline x)^{-1/2}A$, so $\overline A^\top\overline A = A^\top WA$; Figure~\ref{fig:step}, Lemma~\ref{lem:solve_error} \\
$\delta_1,\delta_2,\delta_r$ & the two solve residuals of a step and their sum; Figure~\ref{fig:step}, Lemma~\ref{lem:solve_error} \\
$R,\Cvalid$ & step-sampling diagonal of~\cite{BLLSSSW21} and its parameter; $R = I$ here, Section~\ref{sec:potential_and_progress} \\
$\xhat,\zhat,\shat,\muhat$ & the iterate, multiplier, slack and path parameter a step returns; Figure~\ref{fig:step}, Lemma~\ref{lem:computed_step} \\
$\xhat^\star,\zhat^\star,h_j^\star$ & the exact step of Equation~\eqref{eq:exact_step} and its exact solves; Lemmas~\ref{lem:solve_error}, \ref{lem:step_bounds}, \ref{lem:progress} \\
$\beta$ & size of a perturbation in Lemma~\ref{lem:perturb}; proof of Lemma~\ref{lem:computed_step} \\
$D_c$ & $\diag(c) = \Phi''(x)^{-1/2}$, the barrier scaling of Lemma~\ref{lem:step_bounds}; Lemma~\ref{lem:solver_form}, Lemma~\ref{lem:tau_move} \\
$\tau^{(k)}, T^{(k)}$ & the path Lewis weights $\tau(x^{(k)})$ at round $k$ and their diagonal; Lemma~\ref{lem:stability}, Lemma~\ref{lem:ipm_is_stable} \\
\bottomrule
\end{tabular}
\end{center}
\caption{Variables local to the interior point method, its step
(Section~\ref{sec:potential_and_progress}) and initialization
(Section~\ref{sec:overall}).}
\label{fig:notation_ipm}
\end{figure}

This section states the top level guarantees of one step of the interior point
method, in the setting of Section~\ref{sec:lewis}, and chains them into the path
following bound.
Two simplifications are specific to this version.
A step reads and writes all $n$ coordinates explicitly, which the budget
$\nnz(A) \ge n$ per step allows, so the step is never subsampled: the random
diagonal $R$ of~\cite{BLLSSSW21} is the
identity here, and the imported per-step statements, which are in expectation
over $R$ there, are deterministic.
And the only errors beyond the analysis there are the $\rho^{-10}$ accuracy of
\textsc{Solve} and fixed point truncation; both are handled by one perturbation
lemma, Lemma~\ref{lem:perturb}, which compares the computed step with the exact
step on the same inputs.
The section ends with the movement bounds on the path weights that
Section~\ref{sec:overall} converts into the leverage stability of
Definition~\ref{def:tau_stable}.

A step touches $A$ only through the calls of Definition~\ref{def:operations} and
through $O(1)$ matrix vector products per solver call, so it costs
\[
\O\left( \gamma^{-2} \left( \nnz(A) + d^{2} \right) \log \rho \right)
+ \left( \text{cost of one \textsc{Update}} \right)
\]
bit operations, the $d^{2}$ being the applications of the maintained inverse
inside the \textsc{Solve} calls, including those the Lewis iteration makes, and
$\gamma^{-2} = \O(1)$ the number of Johnson--Lindenstrauss directions that
iteration uses at accuracy $\gamma / 100$; the coordinatewise work of the step,
including the computation of the direction in Lemma~\ref{lem:direction}, is
$\O(n \log \rho) \le \O(\nnz(A) \log \rho)$ and is absorbed.
No matrix produced inside the data structure is visible here.

\paragraph{No subsampling of the step.}
Algorithm 1 of~\cite{BLLSSSW21} multiplies the
residual direction $\delta_r$ of Figure~\ref{fig:step} by a random diagonal
$R$ drawn from a
$\Cvalid$-valid distribution, \cite[Definition
4.13]{BLLSSSW21}, because there a step may
only touch $\O(d)$ coordinates.
Here a step already pays $\nnz(A) \ge n$, so we take
$R \defeq I$.
The identity is $\Cvalid$-valid for every
$\Cvalid \ge \Cnorm$: its expectation is
$I$; the variances of its entries and the deviation
$\|R \delta_r - \delta_r\|_{\infty}$ are zero;
$\E[R_{ii}^{2}] = 1 \le 2 \tau(\overline{A})_i^{-1}$ and
$\E[R_{ii} R_{jj}] = 1 \le 2$ since leverage scores are at most one; and
$\overline{A}^\top R \overline{A} = \overline{A}^\top \overline{A}$ exactly.
Consequently every statement
of~\cite{BLLSSSW21} that is in expectation
over $R$, or holds with probability
$1 - d^{-10}$ over
$R$, holds deterministically here, and the imported lemmas below are stated in
that form.
The only randomness left in a step is inside \textsc{Solve} and inside the Lewis
estimate.

\subsection{The step}

\begin{figure}[H]
\begin{algbox}
\textsc{Step}$(x, z, \mu)$, with
$s \defeq A z + c$:
\begin{enumerate}
\item $\muhat \defeq (1 - \theta) \mu$.
\item Set $\overline{x} \defeq x$ and
\[
\tautil
\defeq
\textsc{LewisEstimate}\left(
\Phi''(x)^{-1/2} A, q, v, \gamma / 100, \textsc{Solve}
\right),
\]
the routine of Lemma~\ref{lem:lewis} run with the data structure's
\textsc{Solve} as its solver and the previous round's $\tautil$ as its starting
estimate, giving $\tautil \approx_{\gamma / 100} \tau(x)$.
\item Call
\[
\textsc{Update}(\what),
\qquad
\what_i
\defeq
\Trunc\left(
\left( \tautil_i \phi_i''(x_i) \right)^{-1}, \, 40 \log \rho
\right)
\text{ for every } i \in [n] ,
\]
so that the data structure now holds
$\widehat W \approx_{\rho^{-18}} W \defeq \widetilde{T}^{-1} \Phi''(x)^{-1}$ by
Equation~\eqref{eq:step_roundoff}, as $w_i \ge \rho^{-1}$ by
Lemma~\ref{lem:magnitudes}.
\item Let
\[
\overline{y}_i
\defeq
\frac{s_i + \mu \tautil_i \phi_i'(x_i)}
{\mu \tautil_i \phi_i''(x_i)^{1/2}},
\qquad
g \defeq \textsc{Direction}(\overline{y}, \tautil),
\]
the routine of Figure~\ref{fig:direction}, which returns a fixed point
approximation of $- \gamma \nabla \Psi(\overline{y})^{\flat(\tautil)}$ that is
admissible in the sense of Definition~\ref{def:admissible}; in particular
$\|g\|_{\tautil + \infty} \le \gamma$.
\item With $\overline{A} \defeq \widetilde{T}^{-1/2} \Phi''(x)^{-1/2} A$, so
  that $\overline{A}^\top \overline{A} = A^\top W A$, solve the two Newton
  systems by
\[
h_1
\defeq
\textsc{Solve}\left(
\Trunc\left( A^\top \Phi''(x)^{-1/2} g, \, 20 \log \rho \right) \right),
\qquad
h_2
\defeq
\textsc{Solve}\left(
\Trunc\left( A^\top x - b, \, 20 \log \rho \right) \right),
\]
and set
\[
\delta_1
\defeq
\widetilde{T}^{-1} \Phi''(x)^{-1/2} A \, h_1,
\qquad
\delta_2
\defeq
\widetilde{T}^{-1} \Phi''(x)^{-1/2} A \, h_2,
\qquad
\delta_r \defeq \delta_1 + \delta_2 .
\]
\item Set
\[
\xhat \defeq \Trunc\left( x + \Phi''(x)^{-1/2} (g - \delta_r), \, 40 \log \rho \right),
\qquad
\zhat \defeq \Trunc\left( z + \mu h_1, \, 40 \log \rho \right),
\]
so that $\shat \defeq A \zhat + c$ equals
$s + \mu \widetilde{T} \Phi''(x)^{1/2} \delta_1$ up to the solve error and the
truncation.
\item \textbf{return} $(\xhat, \zhat)$.
\end{enumerate}
\end{algbox}
\caption{One step of the path-following algorithm for the primal.
Given an $\epsilon$-centered pair
$(x, s = A z + c)$ at path parameter
$\mu$, \textsc{Step} decreases the path parameter by a factor
$1 - \theta$, refreshes the Lewis-weight estimate and the maintained
inverse, and takes one Newton step on the penalized objective;
Lemma~\ref{lem:computed_step} shows the output is $\epsilon$-centered at
the new parameter with contracted potential.}
\label{fig:step}
\end{figure}

Three things about the call order in Figure~\ref{fig:step} are worth recording.
First, the Lewis estimate of step 2 is made before \textsc{Update} in step 3, so
the solver it is handed is the one belonging to the previous round's weights;
that this solver satisfies Item~\ref{itm:lewis:6} is checked in
Lemma~\ref{lem:solver_form} below.
Second, \textsc{Solve} is accurate to $\rho^{-10}$ in the norm of the held
matrix $A^\top \widehat W A \approx_{\rho^{-18}} A^\top W A$, hence to
$3 \rho^{-10}$ in the
$A^\top W A$ norm, and the two places its output enters the step, the direction
$\delta_r$ and the multiplier update
$\mu h_1$, are both linear in it; Lemma~\ref{lem:solve_error} converts the
guarantee into coordinatewise bounds on these two quantities, and
Lemma~\ref{lem:perturb} shows that perturbations of that size do not disturb
centrality or feasibility.
Third, step 3 hands \textsc{Update} the whole weight vector rather than the
coordinates that moved, which costs $O(n) \le O(\nnz(A))$ per round and leaves
the data structure to decide which of its rows to refresh, which is the decision
Corollary~\ref{cor:update_counts} counts.

The following lemma shows that the maintained form serves the Lewis estimate.

\begin{lemma}
\label{lem:solver_form}
Let $\tautil'$ be the estimate of the previous round,
$x'$ the previous iterate, and
$W' = (\widetilde{T}')^{-1} \Phi''(x')^{-1}$, so that the data structure holds
$\widehat W' \approx_{\rho^{-18}} W'$ when step 2 of Figure~\ref{fig:step} runs
at the current iterate $x$.
If $\tautil' \approx_{\gamma / 100} \tau(x')$ and
$\|\Phi''(x')^{1/2} (x - x')\|_{\infty} \le 1 / 100$, then
$\tautil'$ and
$M \defeq A^\top \widehat W' A$ satisfy Items~\ref{itm:lewis:5}
and~\ref{itm:lewis:6} for the matrix $D_c A$,
$D_c = \Phi''(x)^{-1/2}$, whose weights step 2 estimates.
\end{lemma}

\begin{proof}
Item~\ref{itm:lewis:5} asks $\tautil' \approx_{0.5} \tau(x)$.
By Equation~\eqref{eq:barrier_local}
$\Phi''(x)^{-1/2} \approx_{3/200} \Phi''(x')^{-1/2}$, so
Lemma~\ref{lem:lewis_approx} gives
$\tau(x) \approx_{3/50} \tau(x') \approx_{\gamma / 100} \tautil'$.
Item~\ref{itm:lewis:6} asks
$M \approx_{0.5} (D_c A)^\top \diag(\tautil')^{1 - 2/q} D_c A$.
At $q = 1 - \alpha$,
\[
1 - \frac{2}{q}
=
- \frac{1 + \alpha}{1 - \alpha}
=
-1 - \frac{2 \alpha}{1 - \alpha},
\]
so
\begin{equation}
\label{eq:solver_is_lewis_form}
\diag(\tautil')^{1 - 2/q} D_c^{2}
=
\diag(\tautil')^{-2 \alpha / (1 - \alpha)}
\cdot \frac{\Phi''(x')}{\Phi''(x)} \cdot W' ,
\end{equation}
the quotient of diagonal matrices being entrywise.
The second factor is within $e^{\pm 3/100}$ of the identity by
Equation~\eqref{eq:barrier_local}.
For the first, $d / n \le \tau_i(x') \le 5$ gives
$e^{-1/100} d / n \le \tautil'_i \le 6$, so
$|\log \tautil'_i| \le \log(10 n / d)$, and with
$\alpha = 1 / (10 \log(10 n / d)) \le 1/23$,
\[
\frac{2 \alpha}{1 - \alpha} \left| \log \tautil'_i \right|
\le
\frac{2}{10 (1 - 1/23)}
=
\frac{23}{110}
<
\frac{1}{4} .
\]
Hence $A^\top W' A$ and the form of Item~\ref{itm:lewis:6} are within
$e^{\pm (1/4 + 3/100)}$, and
$M \approx_{\rho^{-18}} A^\top W' A$ adds
$\rho^{-18}$ to the exponent, still inside the
$0.5$ of Item~\ref{itm:lewis:6}.
\end{proof}

Along the path the hypotheses of Lemma~\ref{lem:solver_form} are supplied by the
previous round: the accuracy of $\tautil'$ by its step 2, and the movement bound
by Item~\ref{itm:step_bounds:2} together with Lemma~\ref{lem:computed_step}, as
$\gamma \le 1/100$.

The following lemma controls the computed centrality vector.

\begin{lemma}
\label{lem:ybar}
If $(x, s, \mu)$ is
$\epsilon$-centered and
$\tautil \approx_{\gamma / 100} \tau(x)$, then the vector
$\overline{y}$ of step 4 satisfies
$\|\overline{y} - y(x, s, \mu)\|_{\infty} \le \gamma / 20$, truncation included.
\end{lemma}

\begin{proof}
Write $a_i \defeq s_i / (\mu \tau_i(x) \phi_i''(x_i)^{1/2})$ and
$b_i \defeq \phi_i'(x_i) / \phi_i''(x_i)^{1/2}$, so
$y_i = a_i + b_i$,
$|b_i| \le 1$ by Lemma~\ref{lem:self_concordant}, and
$|a_i| \le |y_i| + 1 \le 1 + \epsilon$ by centrality.
Only the first term depends on $\tau$, and
\[
\left| \overline{y}_i - y_i \right|
=
|a_i| \left| \frac{\tau_i(x)}{\tautil_i} - 1 \right|
\le
(1 + \epsilon) \left( e^{\gamma / 100} - 1 \right)
\le
\frac{1.03 \gamma}{100} .
\]
Each entry of $\overline{y}$ is a formula in the stored numbers
$x_i, l_i, u_i, \mu, \tautil_i, c_i$ and the sum
$(A z)_i$, with at most
$20$ operations, namely the two distances to the bounds, their reciprocals and
squared reciprocals for $\phi_i'$ and
$\phi_i''$, one square root, four products and one quotient, of multiplicative
depth $5$; so Equation~\eqref{eq:step_roundoff} bounds the truncation error by
$\rho^{-20} \le \gamma / 100$ per entry.
\end{proof}

\subsection{Direction finding}

The exact direction $-\gamma \nabla \Psi(\overline{y})^{\flat(\tautil)}$
of~\cite{BLLSSSW21} is the maximizer of a
linear functional over the unit ball of $\|\cdot\|_{\tautil + \infty}$; it is
not a rational function of $\overline{y}$ and
$\tautil$, and in fixed point only a near maximizer can be formed.
The progress analysis uses the direction through the three properties below
only, and tolerates a constant fraction of the margin $\alpha / 4$ that the
imported proof has, so a near maximizer suffices.

An admissible direction is defined as follows.

\begin{definition}
\label{def:admissible}
Given $\overline{y} \in \R^{n}$ and a positive
$\tautil \in \R^{n}$, a vector
$g \in \R^{n}$ is an \emph{admissible direction} at
$(\overline{y}, \tautil)$ if
\begin{enumerate}
\item \label{itm:admissible:1} $\|g\|_{\tautil + \infty} \le \gamma$;
\item \label{itm:admissible:2} $g_i \, \nabla \Psi(\overline{y})_i \le 0$ for
  every $i \in [n]$;
\item \label{itm:admissible:3}
  $- \nabla \Psi(\overline{y})^\top g \ge (1 - \alpha / 16) \, \gamma \, \|\nabla \Psi(\overline{y})\|^{*}_{\tautil + \infty}$.
\end{enumerate}
\end{definition}

The exact direction is admissible with $0$ in place of
$\alpha / 16$: Item~\ref{itm:admissible:1} holds with equality by
Definition~\ref{def:taunorm}, Item~\ref{itm:admissible:3} is the definition of
the dual norm, and Item~\ref{itm:admissible:2} holds for the maximizer because
flipping the sign of a coordinate of $h$ that disagrees in sign with
$\nabla \Psi(\overline{y})_i$ leaves
$\|h\|_{\tautil + \infty}$ unchanged and does not decrease
$h^\top \nabla \Psi(\overline{y})$.

\begin{figure}[H]
\begin{algbox}
\textsc{Direction}$(\overline{y}, \tautil)$, all arithmetic on numbers with
$40 \log \rho$ bits, every intermediate result truncated by
$\Trunc(\cdot, 40 \log \rho)$; a tilde marks a quantity computed this way, so
$\widetilde{q}_k$ is the computed value of
$\pi^\top h^{(k)}$ and
$\widetilde{\nu}_k$ that of
$\|h^{(k)}\|_{\tautil + \infty}$.
\begin{enumerate}
\item If $\overline{y} = 0$, \textbf{return}
  $0$.
  Otherwise let $j$ be the integer with
  $2^{-j-1} < \|\overline{y}\|_{\infty} \le 2^{-j}$, and for every
  $i$ compute a number
  $\ghat_i$ with
  $40 \log \rho$ bits such that
  $|\ghat_i - 2^{j} \lambda \sinh(\lambda \overline{y}_i)| \le \rho^{-20}$, as
  $\ghat_i \defeq \Trunc(\Trunc(\lambda z_i) \, \widetilde{S}_i, \, 40 \log \rho)$
  with $z_i \defeq 2^{j} \overline{y}_i$ and
  $\widetilde{S}_i$ the partial sum of the Taylor series of
  $\sinh(t) / t$ at
  $t = \Trunc(\lambda \overline{y}_i, 40 \log \rho)$ with
  $\lceil 6 \lambda \epsilon + 40 \log \rho \rceil$ terms.
\item Normalize and threshold:
  $\pi_i \defeq \Trunc(|\ghat_i| / \|\ghat\|_{\infty}, \, 40 \log \rho)$, then
  $\pi_i \defeq 0$ wherever
  $\pi_i < \rho^{-10}$.
\item Let $\eta$ be a number with
  $40 \log \rho$ bits in
  $[\alpha / 128, \alpha / 64]$ and
  $K \defeq \lceil 512 \alpha^{-1} \log(12 \rho^{12}) \rceil$.
  Set $s_0 \defeq \Trunc(\rho^{-2}, 40 \log \rho)$,
  $s_{k + 1} \defeq \Trunc((1 + \eta) s_k, \, 40 \log \rho)$ for
  $k < K$, and for every
  $k \le K$ form
\[
h^{(k)}_i
\defeq
\min\left( 1, \frac{s_k \, \pi_i}{\tautil_i} \right),
\qquad
\nu_k \defeq \widetilde{\nu}_k + \rho^{-17} ,
\]
the computed norm shifted up so that $\nu_k \ge \|h^{(k)}\|_{\tautil + \infty}$.
\item Let $k^{\star}$ maximize
  $\widetilde{q}_k / \nu_k$ over
  $k$, comparing two fractions by cross multiplication, and write
  $h \defeq h^{(k^{\star})}$,
  $\nu \defeq \nu_{k^{\star}}$.
\item \textbf{return} $g$ with
  $g_i \defeq - \sgn(\ghat_i) \, \Trunc(\gamma h_i / \nu, \, 40 \log \rho)$.
\end{enumerate}
\end{algbox}
\caption{Computing the step direction in fixed point.  Up to scale,
the maximizer of $\pi^\top h$ over the unit ball of
$\|\cdot\|_{\tautil + \infty}$ has the form
$h_i = \min(1, s \pi_i / \tautil_i)$ for a single scalar
$s$, and the ratio
$\pi^\top h / \|h\|_{\tautil + \infty}$ is scale free, so no constraint
has to be solved: the geometric sequence $s_k$ replaces the exact
$s$ by one within a factor
$1 + \eta$, and the final normalization by
$\nu$ makes Item~\ref{itm:admissible:1} hold exactly.}
\label{fig:direction}
\end{figure}

The following lemma shows that \textsc{Direction} computes an admissible
direction.

\begin{lemma}
\label{lem:direction}
Let $(x, s, \mu)$ be
$\epsilon$-centered with
$\Psi(x, s, \mu) \le n^{10}$, let
$\tautil \approx_{\gamma / 100} \tau(x)$, and let
$\overline{y}$ be the vector of step 4 of Figure~\ref{fig:step}.
Then \textsc{Direction}$(\overline{y}, \tautil)$ returns an admissible direction
at $(\overline{y}, \tautil)$, using
$\O(n \log \rho)$ bit operations.
\end{lemma}

\begin{proof}
Write $p \defeq \nabla \Psi(\overline{y})$, so
$p_i = \lambda \sinh(\lambda \overline{y}_i)$, let
$U \defeq \{ h : \|h\|_{\tautil + \infty} \le 1 \}$, and abbreviate
\[
\|\cdot\|^{*}
\defeq
\|\cdot\|^{*}_{\tautil + \infty}
=
\max_{h \in U} \langle \cdot, h \rangle .
\]
As noted after Definition~\ref{def:potential}, $\Psi \le n^{10}$ gives
$|y_i| \le \epsilon$, so
$|\overline{y}_i| \le \epsilon + \gamma / 20 \le 2 \epsilon$ by
Lemma~\ref{lem:ybar} and, by Equation~\eqref{eq:parameters},
\[
|p_i|
\le
\lambda e^{2 \lambda \epsilon}
=
\lambda \left( \frac{C n}{\epsilon^2} \right)^{2 C}
=
n^{O(1)} ;
\]
and $d / (2 n) \le \tautil_i \le 6$ as in the proof of
Lemma~\ref{lem:solver_form}.
So every quantity below has magnitude $\rho^{O(1)}$ and occupies
$O(\log \rho)$ bits.
Two facts about the dual norm are used repeatedly.
Since $U$ is symmetric under sign changes of coordinates,
\begin{equation}
\label{eq:dual_abs}
\|p\|^{*}_{\tautil + \infty}
=
\max_{h \in U} \sum_i |p_i| \, |h_i| ,
\end{equation}
so the dual norm is monotone in $|p|$ entrywise, and
$\|p - p'\|^{*}_{\tautil + \infty} \le \|p - p'\|_1$ because
$\|h\|_{\infty} \le 1$ on
$U$.
And testing $h = e_i / \|e_i\|_{\tautil + \infty}$ for the
$i$ maximizing
$|p_i|$, where
$\|e_i\|_{\tautil + \infty} = 1 + \Cnorm \tautil_i^{1/2} \le 4 \Cnorm$,
\begin{equation}
\label{eq:dual_lower}
\|p\|^{*}_{\tautil + \infty}
\ge
\frac{\|p\|_{\infty}}{4 \Cnorm}
\ge
\frac{\|p\|_{1}}{4 n \Cnorm} .
\end{equation}

The gradient, step 1.
If $\overline{y} = 0$ then
$p = 0$, and
$g = 0$ is admissible.
Otherwise Items~\ref{itm:admissible:2} and~\ref{itm:admissible:3} are unchanged
when $p$ is scaled by a positive constant, so it suffices to prove them for
$\widehat{p} \defeq 2^{j} p$.
Write $z_i \defeq 2^{j} \overline{y}_i$, which is an exact shift of a stored
number since $j \ge 0$, with
$|z_i| \le 1$ and
$\|z\|_{\infty} > 1/2$, and
$S(t) \defeq \sinh(t) / t = \sum_{k \ge 0} t^{2k} / (2k + 1)!$, so that
$\widehat{p}_i = \lambda z_i S(t_i)$ for
$t_i \defeq \lambda \overline{y}_i$, and
$1 \le S(t) \le e^{|t|}$.
Here $|t_i| \le 2 \lambda \epsilon$ and
$e^{2 \lambda \epsilon} = n^{O(1)} \le \rho$ once
$\rho$ is enlarged by a constant power as agreed after
Equation~\eqref{eq:parameters}, so $\|\widehat{p}\|_{\infty} \ge \lambda / 2$
and every quantity below is at most $\rho$ in magnitude.
The computed $\widetilde{t}_i = \Trunc(\lambda \overline{y}_i, 40 \log \rho)$ is
within $\rho^{-40}$ of
$t_i$, and
$c_i \defeq \Trunc(\widetilde{t}_i^{2}, 40 \log \rho)$ within
$3 \rho^{-39}$ of
$t_i^{2}$.
The partial sum $\widetilde{S}_i$ is formed as
$b_0 \defeq 1$,
$b_k \defeq \Trunc(\Trunc(b_{k - 1} c_i) / (2k (2k + 1)), 40 \log \rho)$ and
$\widetilde{S}_i \defeq \Trunc(\sum_{k \le K} b_k)$ with
$K \defeq \lceil 6 \lambda \epsilon + 40 \log \rho \rceil$, the sum accumulated
left to right.
This recurrence has multiplicative depth $K$, for which the bound of
Lemma~\ref{lem:roundoff} is too coarse, so we bound its error directly.
Let $\beta_k \defeq c_i^{k} / (2k + 1)!$ be the exact terms at
$c_i$ and
$e_k \defeq |b_k - \beta_k|$; then
$e_0 = 0$ and
$e_k \le r_k e_{k - 1} + 2 \rho^{-40}$ with
$r_k \defeq c_i / (2k (2k + 1))$, since the product is off by
$c_i e_{k - 1} + \rho^{-40}$ and the division by an integer at least
$2$ does not increase this before its own truncation.
Unrolling,
\[
e_k
\le
2 \rho^{-40} \sum_{m = 0}^{k} \prod_{i' = m + 1}^{k} r_{i'}
=
2 \rho^{-40} \sum_{m = 0}^{k} \frac{\beta_k}{\beta_m}
\le
2 \rho^{-40} (k + 1) \max_{m \le k} \frac{c_i^{k - m}}{(2 (k - m))!}
\le
2 e \rho^{-39} (k + 1),
\]
using $(2k + 1)! / (2m + 1)! \ge (2(k - m))!$ and
$\sum_{j' \ge 0} c^{j'} / (2j')! = \cosh(c^{1/2}) \le e^{|t_i| + 1} \le e \rho$,
as $c_i \le t_i^{2} + 1$.
Hence
$|\widetilde{S}_i - \sum_{k \le K} \beta_k| \le \sum_{k \le K} e_k + K \rho^{-40} \le 2 e (K + 1)^{2} \rho^{-39} \le \rho^{-36}$,
as $K + 1 \le \rho$.
Next, $\sum_{k \le K} \beta_k$ differs from
$\sum_{k \le K} t_i^{2k} / (2k + 1)!$ by at most
$|c_i - t_i^{2}| \sum_{k} k u^{k - 1} / (2k + 1)! \le 3 \rho^{-39} \cosh(u^{1/2}) \le 3 e \rho^{-38}$
for some $u$ between
$c_i$ and
$t_i^{2}$, and the tail of the series of
$S(t_i)$ after
$K$ terms is at most
$\rho^{-42}$: for
$k \ge K / 2 \ge 3 \lambda \epsilon$ the ratio of consecutive terms is
$t_i^{2} / (2k (2k + 1)) \le 1/9$, so the term at
$K$ is at most
$3^{-K} \cosh(t_i) \le \rho^{-43.9} \cdot \rho$ and the tail at most
$9/8$ of it.
Altogether $|\widetilde{S}_i - S(t_i)| \le 2 \rho^{-36}$, and
$\ghat_i = \Trunc(\Trunc(\lambda z_i) \widetilde{S}_i)$ is within
\[
\rho^{-40} \left( \rho + 1 \right) + \rho \cdot 2 \rho^{-36} + \rho^{-40}
\le
\rho^{-20}
\]
of $\lambda z_i S(t_i) = \widehat{p}_i$, as step 1 requires.
In particular $\|\ghat\|_{\infty} \ge \lambda / 2 - \rho^{-20} \ge 1/4$, so
$\ghat \ne 0$.
We work with $\ghat$ in place of
$\widehat{p}$ and charge the difference at the end: wherever
$|\ghat_i| \ge \rho^{-10} \|\ghat\|_{\infty} \ge \rho^{-10} / 4 > \rho^{-20}$,
which is the case wherever step 2 leaves $\pi_i > 0$ since the truncation there
is toward zero, the sign of $\ghat_i$ is that of
$\widehat{p}_i$; and since
$\|\ghat - \widehat{p}\|^{*} \le \|\ghat - \widehat{p}\|_1 \le n \rho^{-20}$
while
$\|\widehat{p}\|^{*} \ge \|\widehat{p}\|_{\infty} / (4 \Cnorm) \ge 1 / (8 \Cnorm)$
by Equation~\eqref{eq:dual_lower},
\[
\left\| \ghat \right\|^{*}
\ge
\left( 1 - 8 n \Cnorm \rho^{-20} \right) \left\| \widehat{p} \right\|^{*}
\ge
\left( 1 - \rho^{-17} \right) \left\| \widehat{p} \right\|^{*},
\]
using $\rho \ge n \ge 8 \Cnorm$.

Normalization and thresholding, step 2.
Let $\overline{\pi} \defeq |\ghat| / \|\ghat\|_{\infty} \in [0, 1]^{n}$, a
vector with a unit entry, so $\|\overline{\pi}\|^{*} \ge 1 / (4 \Cnorm)$ by
Equation~\eqref{eq:dual_lower}.
Each $\pi_i$ differs from
$\overline{\pi}_i$ by at most
$\rho^{-20}$ from the truncation to
$40 \log \rho$ bits, one operation on stored numbers in
Equation~\eqref{eq:step_roundoff}, and by at most $\rho^{-10}$ from the
thresholding, so, using $\rho \ge n \ge \Cnorm$,
\[
\left\| \pi - \overline{\pi} \right\|_1
\le
2 n \rho^{-10}
\le
2 \rho^{-9}
\le
8 \Cnorm \rho^{-9} \left\| \overline{\pi} \right\|^{*}
\le
\rho^{-8} \left\| \overline{\pi} \right\|^{*} .
\]
Hence, for every $h \in U$,
$\overline{\pi}^\top h \ge \pi^\top h - \rho^{-8} \|\overline{\pi}\|^{*}$, and
$\|\pi\|^{*} \ge (1 - \rho^{-8}) \|\overline{\pi}\|^{*}$: a vector
$h \in U$ with
$\pi^\top h \ge (1 - \eta') \|\pi\|^{*}$ has
$\overline{\pi}^\top h \ge (1 - \eta' - 2 \rho^{-8}) \|\overline{\pi}\|^{*}$.
The nonzero entries of $\pi$ lie in
$[\rho^{-10}, 1]$, and
$\pi_i > 0$ implies
$\ghat_i \ne 0$.
Since $-\ghat^\top g = \gamma \|\ghat\|_{\infty} \, \overline{\pi}^\top h / \nu$
up to the final truncation and
$\|\ghat\|^{*} = \|\ghat\|_{\infty} \|\overline{\pi}\|^{*}$, it remains to find
$h \ge 0$ with
$h_i = 0$ wherever
$\pi_i = 0$ and
\begin{equation}
\label{eq:direction_goal}
\frac{\pi^\top h}{\|h\|_{\tautil + \infty}}
\ge
\left( 1 - \frac{\alpha}{32} - \rho^{-9} \right) \|\pi\|^{*} .
\end{equation}

Structure of the maximizer, in exact arithmetic.
By Equation~\eqref{eq:dual_abs} the maximizer defining $\|\pi\|^{*}$ may be
taken nonnegative, and setting to $0$ its coordinates with
$\pi_i = 0$ leaves
$\pi^\top h$ unchanged and does not increase the norm.
As $\pi^\top h$ and
$\|h\|_{\tautil + \infty}$ are both homogeneous of degree one,
\begin{equation}
\label{eq:dual_ratio}
\|\pi\|^{*}
=
\max\left\{
\frac{\pi^\top h}{\|h\|_{\tautil + \infty}}
:
h \ge 0, \ h \ne 0, \ h_i = 0 \text{ wherever } \pi_i = 0
\right\} ,
\end{equation}
and only the direction of $h$ matters.
Let $h^{\star}$ attain the maximum, scaled so that
$\|h^{\star}\|_{\infty} = 1$, and put
$a^{\star} \defeq \Cnorm \|h^{\star}\|_{\tautil}$.
Then $h^{\star}$ maximizes
$\pi^\top h$ over
$\{ h \ge 0 : \|h\|_{\infty} \le 1, \ \Cnorm \|h\|_{\tautil} \le a^{\star} \}$,
since every $h$ in this set has
\[
\pi^\top h
\le
\|\pi\|^{*} \|h\|_{\tautil + \infty}
\le
\|\pi\|^{*} \left( 1 + a^{\star} \right)
=
\pi^\top h^{\star} .
\]
This is a linear objective over the intersection of a box with an ellipsoid, a
convex program with a strictly feasible point, so the Karush--Kuhn--Tucker
conditions hold at $h^{\star}$: for some multiplier
$\xi \ge 0$ of the ellipsoid constraint,
$h^{\star}_i = \min(1, \pi_i / (2 \xi \tautil_i))$ for every
$i$ with
$\pi_i > 0$ when
$\xi > 0$, and
$h^{\star}_i = 1$ for every such
$i$ when
$\xi = 0$.
For $s \in (0, \infty]$ write
\[
h^{(s)}_i
\defeq
\min\left( 1, \frac{s \, \pi_i}{\tautil_i} \right),
\qquad
r(s)
\defeq
\frac{\pi^\top h^{(s)}}{\left\| h^{(s)} \right\|_{\tautil + \infty}} ,
\]
so that $h^{\star} = h^{(s^{\star})}$ with
$s^{\star} = 1 / (2 \xi)$, or
$s^{\star} = \infty$, and
$r(s^{\star}) = \|\pi\|^{*}$.

The range of $s$.
Let $s_{\min} \defeq \min_{i : \pi_i > 0} \tautil_i / \pi_i$ and
$s_{\max} \defeq \max_{i : \pi_i > 0} \tautil_i / \pi_i$.
For $s \le s_{\min}$ no coordinate is capped and
$h^{(s)} = s \, \pi / \tautil$, so
$r(s) = r(s_{\min})$; for
$s \ge s_{\max}$ every coordinate with
$\pi_i > 0$ is capped and
$h^{(s)} = h^{(\infty)}$, so
$r(s) = r(s_{\max})$.
Hence $s^{\star}$ may be taken in
$[s_{\min}, s_{\max}]$.
As the nonzero $\pi_i$ lie in
$[\rho^{-10}, 1]$,
$d / (2 n) \le \tautil_i \le 6$ and
$\rho \ge n \ge 2$,
\[
s_{\min}
\ge
\frac{d}{2 n}
\ge
\rho^{-2},
\qquad
s_{\max}
\le
6 \rho^{10} .
\]

Monotonicity in $s$.
Since $\min(1, c z) \le c \min(1, z)$ for
$c \ge 1$ and
$z \ge 0$, for
$c \ge 1$
\begin{equation}
\label{eq:flat_scaling}
h^{(s)}
\le
h^{(c s)}
\le
c \, h^{(s)}
\qquad
\text{entrywise},
\end{equation}
and as $\pi^\top h$ and
$\|h\|_{\tautil + \infty}$ are monotone and homogeneous on
$h \ge 0$,
\begin{equation}
\label{eq:ratio_loss}
r(c s)
\ge
\frac{\pi^\top h^{(s)}}{c \left\| h^{(s)} \right\|_{\tautil + \infty}}
=
\frac{r(s)}{c},
\qquad
r(s)
\ge
\frac{r(c s)}{c}
\qquad
\text{for } c \ge 1 :
\end{equation}
changing $s$ by a factor
$c$ changes
$r$ by at most the factor
$c$.

The sequence $s_k$.
The computed $s_{k + 1}$ differs from
$(1 + \eta) s_k$ by less than
$\rho^{-20}$, two operations of depth one in Equation~\eqref{eq:step_roundoff},
and by induction
\[
s_k
\ge
s_0
\ge
\rho^{-2} - \rho^{-20}
\ge
\frac{\rho^{-2}}{2},
\]
so with $\alpha = 1 / (10 \log(10 n / d)) \ge \rho^{-1}$,
\[
\left( 1 + \frac{\alpha}{256} \right) s_k
\le
(1 + \eta) \left( 1 - 2 \rho^{-18} \right) s_k
\le
s_{k + 1}
\le
\left( 1 + \frac{\alpha}{64} \right) s_k .
\]
Hence $s_0 \le \rho^{-2} \le s_{\min} \le s^{\star}$, and, by
$\log(1 + t) \ge t / 2$ for
$t \le 1$ and the choice of
$K$,
\[
s_K
\ge
\left( 1 + \frac{\alpha}{256} \right)^{K} \frac{\rho^{-2}}{2}
\ge
e^{\alpha K / 512} \, \frac{\rho^{-2}}{2}
\ge
12 \rho^{12} \cdot \frac{\rho^{-2}}{2}
=
6 \rho^{10}
\ge
s_{\max}
\ge
s^{\star} .
\]
So the first $\overline{k}$ with
$s_{\overline{k}} \ge s^{\star}$ has
\[
s_{\overline{k}}
\le
\left( 1 + \frac{\alpha}{64} \right) s_{\overline{k} - 1}
<
\left( 1 + \frac{\alpha}{64} \right) s^{\star},
\]
and Equation~\eqref{eq:ratio_loss} gives
\begin{equation}
\label{eq:good_k}
r(s_{\overline{k}})
\ge
\frac{\|\pi\|^{*}}{1 + \alpha / 64}
\ge
\left( 1 - \frac{\alpha}{64} \right) \|\pi\|^{*} .
\end{equation}

Truncation of the candidates, step 3.
Fix $k$ and write
$\overline{h} \defeq h^{(s_k)}$.
The computed $h^{(k)}_i$ is
$\min(1, \cdot)$ of the truncated quotient of the truncated product
$s_k \pi_i$ by
$\tautil_i$; both truncations are toward zero on nonnegative quantities, each by
less than $\rho^{-40}$, and the division by
$\tautil_i \ge \rho^{-2}$ amplifies the error of the product by at most
$\rho^{2}$, so
\[
0
\le
\overline{h}_i - h^{(k)}_i
\le
\rho^{-40} \left( 1 + \rho^{2} \right)
\le
\rho^{-17},
\]
and $h^{(k)}_i = 0$ wherever
$\pi_i = 0$.
Changing $h$ by at most
$\rho^{-17}$ per entry changes
$\pi^\top h$ by at most
$n \rho^{-17}$ and
$\|h\|_{\tautil + \infty}$ by at most
$\rho^{-17} (1 + \Cnorm (6 d)^{1/2})$, both at most
$\rho^{-15}$.
On the other side, for the $i_0$ with
$\pi_{i_0} = 1$,
\[
\overline{h}_{i_0}
=
\min\left( 1, \frac{s_k}{\tautil_{i_0}} \right)
\ge
\frac{\rho^{-2}}{12}
\ge
\rho^{-3},
\qquad
\text{so}
\qquad
\pi^\top \overline{h}
\ge
\overline{h}_{i_0}
\ge
\rho^{-3},
\qquad
\left\| \overline{h} \right\|_{\tautil + \infty}
\ge
\left\| \overline{h} \right\|_{\infty}
\ge
\rho^{-3} .
\]
Hence
\[
\left( 1 - \rho^{-12} \right) \pi^\top \overline{h}
\le
\pi^\top h^{(k)}
\le
\pi^\top \overline{h},
\qquad
\left( 1 - \rho^{-12} \right)
\left\| \overline{h} \right\|_{\tautil + \infty}
\le
\left\| h^{(k)} \right\|_{\tautil + \infty}
\le
\left\| \overline{h} \right\|_{\tautil + \infty} .
\]
The computed $\widetilde{q}_k$ is the sum of the
$n$ products
$\pi_i h^{(k)}_i$ of stored numbers, a formula with
$2 n \le 2 \rho$ operations of multiplicative depth
$1$ on numbers of magnitude at most
$\rho$, so Lemma~\ref{lem:roundoff} with
$M = \rho$ and
$\ell = 40 \log \rho$ puts it within
$2 \rho \cdot 2 \rho^{2} \cdot \rho^{-40} = 4 \rho^{-37}$ of
$\pi^\top h^{(k)}$.
For $\widetilde{\nu}_k$, the radicand
$\sum_i \tautil_i (h^{(k)}_i)^{2}$ is a formula with
$3 n \le 3 \rho$ operations of depth
$2$ on stored numbers, so it is off by at most
$3 \rho \cdot (2 \rho^{2})^{2} \cdot \rho^{-40} = 12 \rho^{-35}$ by the same
lemma; it is at least
$\tautil_{i_0} (h^{(k)}_{i_0})^{2} \ge \rho^{-2} \cdot \rho^{-6} / 4 \ge \rho^{-9}$,
so the square root, whose derivative is at most $\rho^{4.5} / 2$ there, is off
by at most $6 \rho^{-30.5} + \rho^{-40}$, and the product with
$\Cnorm \le \rho$ and the sum with the exact maximum
$\|h^{(k)}\|_{\infty}$ leave
$\widetilde{\nu}_k$ within
$\rho^{-40} + \rho \left( 6 \rho^{-30.5} + \rho^{-40} \right) + 2 \rho^{-40} \le \rho^{-17}$
of $\|h^{(k)}\|_{\tautil + \infty}$; hence
\[
\left\| h^{(k)} \right\|_{\tautil + \infty}
\le
\nu_k
\le
\left\| h^{(k)} \right\|_{\tautil + \infty} + 2 \rho^{-17},
\]
and against the lower bounds $\rho^{-3} / 2$ on
$\pi^\top h^{(k)}$ and
$\|h^{(k)}\|_{\tautil + \infty}$ both
\begin{equation}
\label{eq:candidate_ratio}
\frac{\pi^\top h^{(k)}}{\nu_k}
\quad \text{and} \quad
\frac{\widetilde{q}_k}{\nu_k}
\quad \text{lie in} \quad
\left[ \left( 1 - \rho^{-11} \right) r(s_k), \
\left( 1 + \rho^{-11} \right) r(s_k) \right] .
\end{equation}

The selection, steps 4 and 5.
The cross multiplication compares $\widetilde{q}_k \nu_{k'}$ with
$\widetilde{q}_{k'} \nu_k$, products of numbers with
$40 \log \rho$ bits that are exact on
$80 \log \rho$ bits, so
$k^{\star}$ maximizes
$\widetilde{q}_k / \nu_k$ exactly, and by Equation~\eqref{eq:candidate_ratio} at
$k^{\star}$ and at
$\overline{k}$ and Equation~\eqref{eq:good_k},
\[
\frac{\pi^\top h}{\nu}
\ge
\left( 1 - \rho^{-11} \right) r(s_{k^{\star}})
\ge
\left( 1 - \rho^{-11} \right)^{2}
\frac{\widetilde{q}_{k^{\star}}}{\nu_{k^{\star}}}
\ge
\left( 1 - \rho^{-11} \right)^{2}
\frac{\widetilde{q}_{\overline{k}}}{\nu_{\overline{k}}}
\]
and
\[
\left( 1 - \rho^{-11} \right)^{2}
\frac{\widetilde{q}_{\overline{k}}}{\nu_{\overline{k}}}
\ge
\left( 1 - \rho^{-11} \right)^{3} r(s_{\overline{k}})
\ge
\left( 1 - \frac{\alpha}{64} - \rho^{-10} \right) \|\pi\|^{*},
\]
which is Equation~\eqref{eq:direction_goal} for $h / \nu$.
Finally
\[
|g_i|
=
\Trunc\left( \frac{\gamma h_i}{\nu}, \, 40 \log \rho \right)
\le
\frac{\gamma h_i}{\nu},
\qquad
\text{so}
\qquad
\|g\|_{\tautil + \infty}
\le
\frac{\gamma \|h\|_{\tautil + \infty}}{\nu}
\le
\gamma,
\]
which is Item~\ref{itm:admissible:1}, and $h_i = 0$ where
$\pi_i = 0$ while
$\sgn(g_i) = -\sgn(\ghat_i) = -\sgn(\widehat{p}_i) = -\sgn(p_i)$ elsewhere,
which is Item~\ref{itm:admissible:2}.
The truncation toward zero, two operations of depth two in
Equation~\eqref{eq:step_roundoff}, lowers each $|g_i|$ by at most
$\rho^{-20}$ and so
$\sum_i |\ghat_i| \, |g_i|$ by at most
\[
n \rho^{-20} \|\ghat\|_{\infty}
\le
4 \Cnorm \rho^{-19} \|\ghat\|^{*}
\le
\rho^{-17} \gamma \|\ghat\|^{*},
\]
so that
\[
- \ghat^\top g
\ge
\left( 1 - \frac{\alpha}{64} - \rho^{-10} - 2 \rho^{-8} - \rho^{-17} \right)
\gamma \|\ghat\|^{*} .
\]
Passing from $\ghat$ to
$\widehat{p}$: on the support of
$g$ the signs agree and
$|\widehat{p}_i| \ge |\ghat_i| - \rho^{-20}$, so
\[
- \widehat{p}^\top g
\ge
- \ghat^\top g - n \rho^{-20} \gamma
\ge
- \ghat^\top g - 16 n \Cnorm \rho^{-20} \gamma \|\ghat\|^{*}
\ge
- \ghat^\top g - \rho^{-17} \gamma \|\ghat\|^{*},
\]
using $\|\ghat\|^{*} \ge \|\ghat\|_{\infty} / (4 \Cnorm) \ge 1 / (16 \Cnorm)$
and $\rho \ge n \ge 16 \Cnorm$, and
$\|\ghat\|^{*} \ge (1 - \rho^{-17}) \|\widehat{p}\|^{*}$ from step 1.
Collecting the losses of the normalization, the gradient and the truncation,
\[
- \widehat{p}^\top g
\ge
\left( 1 - \rho^{-17} \right)
\left( 1 - \frac{\alpha}{64} - \rho^{-10} - 2 \rho^{-8} - 2 \rho^{-17} \right)
\gamma \|\widehat{p}\|^{*}
\]
and
\[
\left( 1 - \rho^{-17} \right)
\left( 1 - \frac{\alpha}{64} - \rho^{-10} - 2 \rho^{-8} - 2 \rho^{-17} \right)
\ge
1 - \frac{\alpha}{64} - \rho^{-7}
\ge
1 - \frac{\alpha}{16},
\]
so
\[
- \widehat{p}^\top g
\ge
\left( 1 - \frac{\alpha}{16} \right)
\gamma \|\widehat{p}\|^{*},
\]
which is Item~\ref{itm:admissible:3} for $\widehat{p}$, hence for
$p = 2^{-j} \widehat{p}$.

Cost.
Step 1 is $O(n K) = O(n (\lambda \epsilon + \log \rho)) = O(n \log \rho)$
operations, as $\lambda \epsilon \le 3 C \log n$ and
$\rho \ge n$; step 2 is
$O(n)$, and step 3 is
$K + 1 = O(\alpha^{-1} \log \rho) = O(\log n \log \rho)$ candidates at
$O(n + \log \rho)$ operations each, all on numbers with
$40 \log \rho$ bits:
$\O(n \log \rho)$ bit operations in total, within the budget of a step.
\end{proof}

\subsection{Centrality and feasibility}

The following lemma bounds the solve error in the step.

\begin{lemma}
\label{lem:solve_error}
In step 5 of Figure~\ref{fig:step} let
$h_j^{\star} \defeq (A^\top W A)^{-1} b_j$ be the exact solutions for the right
hand sides $b_1 = A^\top \Phi''(x)^{-1/2} g$ and
$b_2 = A^\top x - b$, and let
$\delta_j^{\star}$,
$\delta_r^{\star}$ be formed from them as
$\delta_j, \delta_r$ are from
$h_j$.
If $(x, s, \mu)$ is
$\epsilon$-centered,
$\tautil \approx_{\gamma / 100} \tau(x)$ and
$\|g\|_{\tautil + \infty} \le \gamma$, then
\[
\left\| h_1^{\star} \right\|_{A^\top W A} \le \gamma,
\qquad
\left\| h_2^{\star} \right\|_{A^\top W A} \le \gamma,
\]
and, for the outputs of \textsc{Solve},
\[
\left\| \delta_r - \delta_r^{\star} \right\|_{\infty}
\le
12 \left( \frac{n}{d} \right)^{1/2} \rho^{-10} \gamma,
\qquad
\left\| \left( \mu \widetilde{T} \Phi''(x)^{1/2} \right)^{-1}
A \left( \mu h_1 - \mu h_1^{\star} \right) \right\|_{\infty}
\le
6 \left( \frac{n}{d} \right)^{1/2} \rho^{-10} \gamma .
\]
\end{lemma}

\begin{proof}
For any $e \in \R^{d}$,
$\widetilde{T}^{-1} \Phi''(x)^{-1/2} A e = \widetilde{T}^{-1/2} \overline{A} e$
and
\[
\left\| \overline{A} e \right\|_{2}^{2}
=
e^\top \overline{A}^\top \overline{A} e
=
\left\| e \right\|_{A^\top W A}^{2},
\]
so every coordinate of $\widetilde{T}^{-1} \Phi''(x)^{-1/2} A e$ is at most
$\tautil_i^{-1/2} \|e\|_{A^\top W A} \le 2 (n/d)^{1/2} \|e\|_{A^\top W A}$,
using $\tautil_i \ge e^{-\gamma/100} d / n$.
Both displayed bounds follow by taking $e = h_j - h_j^{\star}$ once
$\|e\|_{A^\top W A} \le 3 \rho^{-10} \|h_j^{\star}\|_{A^\top W A}$ is shown and
the two norms $\|h_j^{\star}\|_{A^\top W A}$ are bounded by
$\gamma$.
For the first, write $H \defeq A^\top W A$ and
$\widehat H \defeq A^\top \widehat W A$ for the matrix the data structure holds
after step 3, so that $\widehat H \approx_{\rho^{-18}} H$ and
Equation~\eqref{eq:solve_guarantee} reads
$\|h_j - \widehat H^{-1} b_j\|_{\widehat H} \le \rho^{-10} \|b_j\|_{\widehat H^{-1}}$.
For symmetric positive definite $\widehat H \approx_{\delta} H$ and any
$b$,
$\|b\|_{\widehat H^{-1}} \le e^{\delta / 2} \|b\|_{H^{-1}}$ and
\[
\left\| \widehat H^{-1} b - H^{-1} b \right\|_{H}
=
\left\| \left( H^{1/2} \widehat H^{-1} H^{1/2} - I \right)
H^{-1/2} b \right\|_{2}
\le
\left( e^{\delta} - 1 \right) \|b\|_{H^{-1}},
\]
so with $\delta = \rho^{-18}$ and
$\|b_j\|_{H^{-1}} = \|h_j^{\star}\|_{H}$,
\[
\left\| h_j - h_j^{\star} \right\|_{H}
\le
e^{\delta / 2}
\left\| h_j - \widehat H^{-1} b_j \right\|_{\widehat H}
+ \left\| \widehat H^{-1} b_j - H^{-1} b_j \right\|_{H}
\le
\left( e^{\delta} \rho^{-10} + e^{\delta} - 1 \right)
\left\| h_j^{\star} \right\|_{H}
\le
3 \rho^{-10} \left\| h_j^{\star} \right\|_{H} .
\]

For $h_1^{\star}$:
$b_1 = \overline{A}^\top \widetilde{T}^{1/2} g$, so
$\overline{A} h_1^{\star} = P \widetilde{T}^{1/2} g$ with
$P$ the orthogonal projection onto the range of
$\overline{A}$, and
\[
\left\| h_1^{\star} \right\|_{A^\top W A}
=
\left\| P \widetilde{T}^{1/2} g \right\|_{2}
\le
\left\| \widetilde{T}^{1/2} g \right\|_{2}
=
\|g\|_{\tautil}
\le
\frac{\gamma}{\Cnorm} .
\]
For $h_2^{\star}$:
$\|h_2^{\star}\|_{A^\top W A}^{2} = (A^\top x - b)^\top (A^\top W A)^{-1} (A^\top x - b)$,
and
\[
W
=
\left( \widetilde{T} \Phi''(x) \right)^{-1}
\approx_{\gamma / 100}
\left( T(x) \Phi''(x) \right)^{-1},
\]
so
\[
\left\| h_2^{\star} \right\|_{A^\top W A}
\le
e^{\gamma / 200}
\left\| A^\top x - b \right\|_{(A^\top (T(x) \Phi''(x))^{-1} A)^{-1}}
\le
e^{\gamma / 200} \frac{\epsilon \gamma}{\Cnorm}
\le
\gamma
\]
by Item~\ref{itm:centered:3}.
\end{proof}

The following lemma bounds the effect of perturbations on centrality and
feasibility.

\begin{lemma}
\label{lem:perturb}
Let $s = A z + c$ with
$\Psi(x, s, \mu) \le n^{10}$, and let
$x', z'$ satisfy
\[
\left\| \Phi''(x)^{1/2} (x' - x) \right\|_{\infty} \le \beta,
\qquad
\left\| \left( \mu T(x) \Phi''(x)^{1/2} \right)^{-1}
A (z' - z) \right\|_{\infty} \le \beta
\]
for some $\beta \le 1/100$, and put
$s' \defeq A z' + c$.
Then
\begin{enumerate}
\item \label{itm:perturb:1}
  $\Psi(x', s', \mu) \le e^{40 \lambda \beta} \Psi(x, s, \mu)$;
\item \label{itm:perturb:2}
  $\displaystyle \left\| A^\top x' - b \right\|_{(A^\top (T(x') \Phi''(x'))^{-1} A)^{-1}} \le e^{5 \beta} \left( \left\| A^\top x - b \right\|_{(A^\top (T(x) \Phi''(x))^{-1} A)^{-1}} + (5 d)^{1/2} \beta \right)$;
\item \label{itm:perturb:3} $s'$ is dual feasible.
\end{enumerate}
\end{lemma}

\begin{proof}
By Equation~\eqref{eq:barrier_local}, $\Phi''(x') \approx_{3 \beta} \Phi''(x)$
and $|\phi_i'(x'_i) - \phi_i'(x_i)| \le 2 \beta \phi_i''(x_i)^{1/2}$; by
Lemma~\ref{lem:lewis_approx} applied to
$\Phi''(x')^{-1/2} \approx_{3 \beta / 2} \Phi''(x)^{-1/2}$,
$\tau(x') \approx_{6 \beta} \tau(x)$.

For Item~\ref{itm:perturb:1}, write $y_i = a_i + b_i$ with
$a_i = s_i / (\mu \tau_i(x) \phi_i''(x_i)^{1/2})$ and
$b_i = \phi_i'(x_i) / \phi_i''(x_i)^{1/2}$, and
$y'_i = a'_i + b'_i$ likewise at
$(x', s')$.
From $\cosh(\lambda y_i) \le n^{10}$,
$|y_i| \le \log(2 n^{10}) / \lambda \le 1$, so
$|a_i| \le 2$ and
$|b_i| \le 1$.
The numerator of $a'_i$ is
$s_i + (A (z' - z))_i$, and the second summand is at most
$\beta$ times the denominator of
$a_i$; the denominator of
$a'_i$ is within
$e^{\pm 7.5 \beta}$ of that of
$a_i$.
Hence
\[
|a'_i - a_i|
\le
(2 + \beta) \left( e^{7.5 \beta} - 1 \right) + \beta
\le
18 \beta,
\qquad
|b'_i - b_i|
\le
\left( e^{1.5 \beta} - 1 \right) + 2 \beta e^{1.5 \beta}
\le
4 \beta,
\]
so $|y'_i - y_i| \le 22 \beta$.
Finally
\[
\cosh(a + b)
=
\cosh a \cosh b + \sinh a \sinh b
\le
\cosh(a) \, e^{|b|},
\]
so $\psi(y'_i) \le \psi(y_i) e^{22 \lambda \beta}$ for every
$i$, and summing gives Item~\ref{itm:perturb:1}.

For Item~\ref{itm:perturb:2} put $D \defeq T(x) \Phi''(x)$ and
$D' \defeq T(x') \Phi''(x')$, so
$D' \approx_{9 \beta} D$, and the two norms in the statement are within
$e^{\pm 4.5 \beta}$ of each other.
By the triangle inequality it remains to bound
$\|A^\top (x' - x)\|_{(A^\top D^{-1} A)^{-1}}$, and with
$r \defeq x' - x$,
\[
\left\| A^\top r \right\|_{(A^\top D^{-1} A)^{-1}}^{2}
=
\left( D^{1/2} r \right)^\top
D^{-1/2} A \left( A^\top D^{-1} A \right)^{-1} A^\top D^{-1/2}
\left( D^{1/2} r \right)
\le
\left\| D^{1/2} r \right\|_{2}^{2},
\]
because the middle matrix is an orthogonal projection, and
\[
\left\| D^{1/2} r \right\|_{2}^{2}
=
\sum_i \tau_i(x) \phi_i''(x_i) r_i^{2}
\le
5 d \beta^{2}
\]
by $\|\tau(x)\|_1 \le 5 d$.
Item~\ref{itm:perturb:3} is the definition of $s'$.
\end{proof}

\subsection{Per-step bounds and progress}

We now import the per-step analysis.
To separate what is imported from what is new, define the \emph{exact step} on
the inputs of \textsc{Step} and the direction $g$ of its step 4 as the pair
\begin{equation}
\label{eq:exact_step}
\xhat^{\star}
\defeq
x + \Phi''(x)^{-1/2} \left( g - \delta_r^{\star} \right),
\qquad
\zhat^{\star}
\defeq
z + \mu h_1^{\star},
\qquad
\shat^{\star}
\defeq
A \zhat^{\star} + c
=
s + \mu \widetilde{T} \Phi''(x)^{1/2} \delta_1^{\star},
\end{equation}
that is, steps 5 and 6 with exact solves and no truncation.
The exact step is Algorithm 1
of~\cite{BLLSSSW21} with $R = I$,
$\overline{x} = x$ and
$\overline{\tau} = \tautil$, and the three lemmas below are its analysis there,
stated for $R = I$.
Their hypotheses are those of Definitions~\ref{def:centered}
and~\ref{def:invariant} together with $\|g\|_{\tautil + \infty} \le \gamma$; the
source states the norm bound on $g$ in
$\|\cdot\|_{\tau(x) + \infty}$, which differs from
$\|\cdot\|_{\tautil + \infty}$ by the factor
$e^{\gamma / 200}$ of Equation~\eqref{eq:norm_compare}, absorbed by the
$O(\gamma)$ slack of the statements.

The bounds on the exact step are Lemmas 4.29, 4.31 and 4.33 of
\cite{BLLSSSW21}.

\begin{lemma}
\label{lem:step_bounds}
Let $(x, s, \mu)$ be
$\epsilon$-centered, let
$\overline{x} = x, \tautil$ satisfy Definition~\ref{def:invariant}, let
$\|g\|_{\tautil + \infty} \le \gamma$, and let
$\delta_x \defeq \xhat^{\star} - x$,
$c \defeq \Phi''(x)^{-1/2}$,
$D_c \defeq \diag(c)$,
$\delta_c \defeq \Phi''(\xhat^{\star})^{-1/2} - \Phi''(x)^{-1/2}$, and
$\tau \defeq \tau(x)$.
Then
\begin{enumerate}
\item \label{itm:step_bounds:1}
  $\|D_c^{-1} \delta_x\|_{\tau + \infty} \le \gamma + O( (\epsilon + \Cnorm^{-1}) \gamma )$;
\item \label{itm:step_bounds:2} $\|D_c^{-1} \delta_x\|_{\infty} \lesssim \gamma$
  and $\|D_c^{-1} \delta_c\|_{\infty} \lesssim \gamma$;
\item \label{itm:step_bounds:3}
  $\|D_c^{-1} \delta_c\|_{\tau + \infty} \le \gamma + O( (\epsilon + \Cnorm^{-1}) \gamma )$;
\item \label{itm:step_bounds:4}
  $\|D_c^{-2} \delta_c^2\|_{\tau + \infty} \lesssim \gamma^2$;
\item \label{itm:step_bounds:5}
  $\|D_c^{-1} |\delta_x| \|_{P^{(2)}} \lesssim \gamma / \Cnorm$, where
  $P = \overline{A}(\overline{A}^\top \overline{A})^{-1}\overline{A}^\top$ and
  $P^{(2)}$ is its entrywise square.
\end{enumerate}
\end{lemma}

Because $\delta_c$ is
$\gamma$-bounded in the sense of~\cite[Definition
4.23]{BLLSSSW21}, the Jacobian bounds for the
regularized Lewis weight map transfer these bounds to the weights themselves.
This is the $\tau_q$ statement referred to in Section~\ref{sec:dynamic}, at the
value $q = 1 - \alpha$ of Definition~\ref{def:central_weights}.

The movement of $\tau_q$ in one exact step is Lemmas 4.20 and 4.35 of
\cite{BLLSSSW21}.

\begin{lemma}
\label{lem:tau_move}
In the setting of Lemma~\ref{lem:step_bounds}, let
$\delta_\tau \defeq \tau(\xhat^{\star}) - \tau(x)$,
$T \defeq T(x)$, and let
$J$ be the Jacobian of
$c \mapsto \tau_{q}(\diag(c) A, v)$ at
$c$.
Then
\begin{enumerate}
\item \label{itm:tau_move:1} $\|T^{-1} \delta_\tau\|_{\infty} \lesssim \gamma$;
\item \label{itm:tau_move:2}
  $\|T^{-1} (\delta_\tau - J \delta_c)\|_{\tau + \infty} \lesssim \gamma^2$;
\item \label{itm:tau_move:3}
  $\|T^{-1} J h\|_{\tau + \infty} \lesssim \|D_c^{-1} h\|_{\tau + \infty}$ for
  every $h \in \R^{n}$.
\end{enumerate}
Consequently $\|T^{-1} \delta_\tau\|_{\tau + \infty} \lesssim \gamma$.
\end{lemma}

Item~\ref{itm:tau_move:3} is the decomposition $T^{-1} J D_c = \Lambda + K$
of~\cite[Lemma 4.20]{BLLSSSW21} into a
diagonal $0 \preceq \Lambda \preceq I$ and a matrix
$K$ with
$\|K h\|_{\infty} \lesssim \|h\|_{\infty}$ and
$\|K h\|_{\tau} \lesssim \||h|\|_{P^{(2)}} \le \|h\|_{\tau}$, the last step by
Cauchy--Schwarz and $\sum_j P_{ij}^{2} = P_{ii} \le \tau_i$.
The consequence is Items~\ref{itm:tau_move:2} and~\ref{itm:tau_move:3} applied
to $h = \delta_c$, with Item~\ref{itm:step_bounds:3}.

The progress of the exact step with an admissible direction follows Lemmas 4.38,
4.39 and Corollary 4.40 of \cite{BLLSSSW21},
with the hypothesis on the direction relaxed.

\begin{lemma}
\label{lem:progress}
In the setting of Lemma~\ref{lem:step_bounds}, suppose moreover that
$\|\overline{y} - y(x, s, \mu)\|_{\infty} \le \gamma / 20$, that
$g$ is an admissible direction at
$(\overline{y}, \tautil)$ in the sense of Definition~\ref{def:admissible}, and
let $\muhat = (1 - \theta) \mu$.
Then $\shat^{\star}$ is dual feasible,
\[
\left\| A^\top \xhat^{\star} - b
\right\|_{(A^\top (T(\xhat^{\star}) \Phi''(\xhat^{\star}))^{-1} A)^{-1}}
\le
\frac{\epsilon \gamma}{2 \Cnorm},
\]
and
\[
\Psi(\xhat^{\star}, \shat^{\star}, \muhat)
\le
\left( 1 - \kappa \right) \Psi(x, s, \mu) + n .
\]
\end{lemma}

\begin{proof}
Write $y \defeq y(x, s, \mu)$,
$\psi'(y)_i \defeq \lambda \sinh(\lambda y_i)$ and
$\psi''(y)_i \defeq \lambda^{2} \cosh(\lambda y_i)$, so that
$\nabla \Psi(y) = \psi'(y)$, and abbreviate
$\|\cdot\|^{*} \defeq \|\cdot\|^{*}_{\tautil + \infty}$ and
$U \defeq \{ h : \|h\|_{\tautil + \infty} \le 1 \}$.
Dual feasibility and the feasibility bound are~\cite[Lemma
4.38]{BLLSSSW21}, and~\cite[Lemma
4.39]{BLLSSSW21} gives, for a universal
constant $c_{2}$,
\begin{equation}
\label{eq:first_drop}
\Psi(\xhat^{\star}, \shat^{\star}, \muhat)
\le
\Psi(x, s, \mu)
+ \psi'(y)^\top g
+ \left( 1 - \frac{\alpha}{4} \right) \gamma \|\psi'(y)\|^{*}
+ c_{2} \gamma^{2} \|\psi''(y)\|^{*} .
\end{equation}
Both statements use $g$ only through
$\|g\|_{\tautil + \infty} \le \gamma$, which is Item~\ref{itm:admissible:1}; the
direction enters Lemma 4.39 there only through the term $\psi'(y)^\top g$, which
is kept explicit in Equation~\eqref{eq:first_drop}.
The source then bounds this term for the exact maximizer by \cite[Lemma
4.36]{BLNPSSSW20} (the argument
is~\cite[Corollary 4.40]{BLLSSSW21}); we rerun
that argument with Items~\ref{itm:admissible:2} and~\ref{itm:admissible:3} in
place of the maximizer.

The elementary inequality.
For $a \ge 0$ and
$|t| \le \delta' \le 1 / 5$, expanding
$\sinh(a + t) = \sinh(a) \cosh(t) + \cosh(a) \sinh(t)$ with
$\cosh(t) \ge 1$,
$\sinh(t) \ge - \sinh(\delta') \ge - 2 \delta'$ and
$\cosh(a) \le \sinh(a) + 1$ gives
\begin{equation}
\label{eq:sinh_shift}
\sinh(a + t)
\ge
\left( 1 - 2 \delta' \right) \sinh(a) - 2 \delta' .
\end{equation}
Let $\delta \defeq \gamma / 20$, so that
$\lambda \delta \le 1 / 5$ by Equation~\eqref{eq:parameters}, and write
$p \defeq \nabla \Psi(\overline{y})$.

The correlation with $\psi'(y)$.
Fix $i$ with
$g_i \ne 0$; by Item~\ref{itm:admissible:2},
$p_i \ne 0$ and
$g_i$ has the sign opposite to
$\overline{y}_i$, so with
$a \defeq \lambda |\overline{y}_i|$ and
$t \defeq \lambda \sgn(\overline{y}_i) (y_i - \overline{y}_i)$,
$|t| \le \lambda \delta$, Equation~\eqref{eq:sinh_shift} gives
\[
- \psi'(y)_i g_i
=
|g_i| \, \lambda \sinh(a + t)
\ge
\left( 1 - 2 \lambda \delta \right) |g_i| \, |p_i| - 2 \lambda^{2} \delta \, |g_i| .
\]
Summing over $i$ and using Item~\ref{itm:admissible:3} together with
$\|g\|_1 \le n \|g\|_{\infty} \le n \gamma$,
\begin{equation}
\label{eq:corr_shift}
- \psi'(y)^\top g
\ge
\left( 1 - 2 \lambda \delta \right)
\left( 1 - \frac{\alpha}{16} \right) \gamma \|p\|^{*}
- 2 \lambda^{2} \delta \gamma n .
\end{equation}
Next, $\|p\|^{*}$ is compared with
$\|\psi'(y)\|^{*}$.
By Equation~\eqref{eq:dual_abs} the dual norm is
$\max_{h \in U} \sum_i |p_i| |h_i|$; for the maximizer
$h$ of
$\psi'(y)$, Equation~\eqref{eq:sinh_shift} with
$a \defeq \lambda |y_i|$ and
$t \defeq \lambda (|\overline{y}_i| - |y_i|)$ gives
$|p_i| \ge (1 - 2 \lambda \delta) |\psi'(y)_i| - 2 \lambda^{2} \delta$, and
since $\|h\|_{\infty} \le 1$ on
$U$,
\[
\|p\|^{*}
\ge
\sum_i |p_i| \, |h_i|
\ge
\left( 1 - 2 \lambda \delta \right) \|\psi'(y)\|^{*} - 2 \lambda^{2} \delta n .
\]
Substituting into Equation~\eqref{eq:corr_shift} and using
$4 \lambda \delta = \lambda \gamma / 5 \le \alpha / (5 C^{2})$,
\begin{equation}
\label{eq:corr_final}
\psi'(y)^\top g
\le
- \left( 1 - \frac{\alpha}{16} - \frac{\lambda \gamma}{5} \right)
\gamma \|\psi'(y)\|^{*}
+ \frac{\lambda^{2} \gamma^{2} n}{5} .
\end{equation}

The second order term.
Entrywise
\[
\psi''(y)_i
=
\lambda^{2} \cosh(\lambda y_i)
\le
\lambda \left| \psi'(y)_i \right| + \lambda^{2},
\]
so by the monotonicity of the dual norm in the absolute values and
$\|1\|^{*} \le n$,
\[
c_{2} \gamma^{2} \|\psi''(y)\|^{*}
\le
c_{2} \lambda \gamma \cdot \gamma \|\psi'(y)\|^{*} + c_{2} \lambda^{2} \gamma^{2} n
\le
\frac{\alpha}{32} \gamma \|\psi'(y)\|^{*} + c_{2} \lambda^{2} \gamma^{2} n,
\]
using $c_{2} \lambda \gamma = c_{2} \alpha / C^{2} \le \alpha / 32$ for
$C \ge 32 c_{2}$.

Collecting.
Equations~\eqref{eq:first_drop}, \eqref{eq:corr_final} and the last display give
\[
\begin{aligned}
\Psi(\xhat^{\star}, \shat^{\star}, \muhat)
&\le
\Psi(x, s, \mu)
- \left( \frac{\alpha}{4} - \frac{\alpha}{16} - \frac{\alpha}{32}
- \frac{\lambda \gamma}{5} \right) \gamma \|\psi'(y)\|^{*}
+ \left( c_{2} + \frac{1}{5} \right) \lambda^{2} \gamma^{2} n
\\
&\le
\Psi(x, s, \mu)
- \frac{\alpha}{8} \gamma \|\psi'(y)\|^{*}
+ \left( c_{2} + \frac{1}{5} \right) \lambda^{2} \gamma^{2} n .
\end{aligned}
\]
Finally $\|\psi'(y)\|^{*}$ is converted into the potential.
Since $\|h\|_{\tautil + \infty} \le \|h\|_{\infty} \|1\|_{\tautil + \infty}$ and
\[
\|1\|_{\tautil + \infty}
=
1 + \Cnorm \|\tautil\|_1^{1/2}
\le
1 + \Cnorm (6 d)^{1/2}
\le
4 \Cnorm d^{1/2}
=
\frac{4 C d^{1/2}}{\alpha},
\]
the vector $h \defeq \alpha \, \sgn(\psi'(y)) / (4 C d^{1/2})$ lies in
$U$ and gives
$\|\psi'(y)\|^{*} \ge \alpha \|\psi'(y)\|_1 / (4 C d^{1/2})$, while
$|\lambda \sinh(\lambda y_i)| \ge \lambda (\cosh(\lambda y_i) - 1)$ gives
$\|\psi'(y)\|_1 \ge \lambda \Psi(x, s, \mu) - \lambda n$.
Hence, with $\kappa = \alpha^{2} \lambda \gamma / (32 C d^{1/2})$,
\[
\Psi(\xhat^{\star}, \shat^{\star}, \muhat)
\le
\left( 1 - \kappa \right) \Psi(x, s, \mu)
+ \kappa n
+ \left( c_{2} + \frac{1}{5} \right) \lambda^{2} \gamma^{2} n
\le
\left( 1 - \kappa \right) \Psi(x, s, \mu) + n,
\]
since $\kappa \le \lambda \gamma / 32 \le 1 / 256$ and
$(c_{2} + 1/5) \lambda^{2} \gamma^{2} \le (1 / 32 + 1 / 40) / 8$ by
$\lambda \gamma \le 1 / 8$ and
$c_{2} \lambda \gamma \le 1 / 32$.
\end{proof}

The following lemma is the guarantee for the computed step.

\begin{lemma}
\label{lem:computed_step}
Let $(x, s = A z + c, \mu)$ be
$\epsilon$-centered with
$\Psi(x, s, \mu) \le n^{2}$, and let the estimate
$\tautil$ of step 2 and the \textsc{Solve} calls of step 5 meet their
guarantees.
Then the output $(\xhat, \zhat)$ of
\textsc{Step}$(x, z, \mu)$ satisfies
\[
\left\| \Phi''(x)^{1/2} \left( \xhat - \xhat^{\star} \right)
\right\|_{\infty}
\le \rho^{-8} \gamma,
\qquad
\left\| \left( \muhat T(\xhat^{\star})
\Phi''(\xhat^{\star})^{1/2} \right)^{-1}
A \left( \zhat - \zhat^{\star} \right) \right\|_{\infty}
\le \rho^{-8} \gamma,
\]
and, with $\shat = A \zhat + c$, the triple
$(\xhat, \shat, \muhat)$ is
$\epsilon$-centered and
\[
\Psi(\xhat, \shat, \muhat)
\le
\left( 1 - \frac{\kappa}{2} \right) \Psi(x, s, \mu) + 2 n .
\]
\end{lemma}

\begin{proof}
Step 2 gives $\tautil \approx_{\gamma / 100} \tau(x)$, so
Definition~\ref{def:invariant} holds, Lemma~\ref{lem:ybar} supplies the
hypothesis on $\overline{y}$ of Lemma~\ref{lem:progress}, and
Lemma~\ref{lem:direction}, applicable as $\Psi \le n^{2} \le n^{10}$, shows that
the $g$ of step 4 is an admissible direction, which is the remaining hypothesis
of Lemma~\ref{lem:progress}; in particular
$\|g\|_{\tautil + \infty} \le \gamma$, as Lemmas~\ref{lem:step_bounds}
and~\ref{lem:solve_error} require.

The two displayed bounds.
By Equation~\eqref{eq:exact_step},
$\xhat - \xhat^{\star} = \Phi''(x)^{-1/2} (\delta_r^{\star} - \delta_r) + r_x$
and $\zhat - \zhat^{\star} = \mu (h_1 - h_1^{\star}) + r_z$, where
$r_x, r_z$ collect the truncation.
The solve error terms are bounded by Lemma~\ref{lem:solve_error}: the first by
$12 (n/d)^{1/2} \rho^{-10} \gamma \le 12 \rho^{-9.5} \gamma$ as
$\rho \ge n$, and the second by
$6 (n/d)^{1/2} \rho^{-10} \gamma$ in the norm with
$\mu \widetilde{T} \Phi''(x)^{1/2}$ in place of
$\muhat T(\xhat^{\star}) \Phi''(\xhat^{\star})^{1/2}$; the two diagonals are
within a factor $3$ since
$\muhat \ge \mu / 2$,
$\tau(\xhat^{\star}) \approx_{O(\gamma)} \tau(x) \approx_{\gamma/100} \tautil$
and $\Phi''(\xhat^{\star}) \approx_{O(\gamma)} \Phi''(x)$ by
Item~\ref{itm:step_bounds:2} and Equation~\eqref{eq:barrier_local}.
For the truncation, every entry of $\xhat$ and
$\zhat$ is a formula with at most
$12$ operations and multiplicative depth at most
$6$ on the stored numbers, the outputs
$h_1, h_2$ of \textsc{Solve} and the sums
$(A h_j)_i$: two distances to the bounds, their squared reciprocals, a sum, a
square root and a reciprocal for $\phi_i''(x_i)^{-1/2}$, then the difference
$g_i - (\delta_r)_i$, one product and one sum for
$\xhat_i$, and one product and one sum for
$\zhat_i$; so Equation~\eqref{eq:step_roundoff} gives
$\|r_x\|_{\infty}, \|r_z\|_{\infty} \le \rho^{-20}$.
The right hand sides passed to \textsc{Solve} are sums of the second kind over
entrywise formulas, cut to $20 \log \rho$ bits, so by
Equation~\eqref{eq:sum_roundoff} each entry of $b_j$ is off by at most
$\rho^{-16} + \rho^{-20} \le 2 \rho^{-16}$, which moves
$h_j^{\star} = (A^\top W A)^{-1} b_j$ by at most
$\|(A^\top W A)^{-1/2}\| \cdot 2 d^{1/2} \rho^{-16} \le 2 \rho^{1/2} d^{1/2} \rho^{-16}$
in the $A^\top W A$ norm, and this is handled like the solve error.
Then $\|\Phi''(x)^{1/2} r_x\|_{\infty} \le \rho^{1/2} \rho^{-20}$, and
\[
\left| (A r_z)_i \right|
\le
\|A\| \, \|r_z\|_2
\le
\rho^{1/2} d^{1/2} \rho^{-20}
\]
against a denominator
$\muhat \tau_i \phi_i''^{1/2} \ge \rho^{-1} (d / 2n) \rho^{-1/2}$.
All of these are below $\rho^{-8} \gamma / 2$ for
$\rho \ge n$ large enough, and the two bounds follow.

The potential and feasibility.
Lemma~\ref{lem:progress} gives
$\Psi(\xhat^{\star}, \shat^{\star}, \muhat) \le n^{2} + n \le n^{10}$, so
Lemma~\ref{lem:perturb} applies at $(\xhat^{\star}, \zhat^{\star}, \muhat)$ with
$\beta = \rho^{-8} \gamma$:
\[
\begin{aligned}
\Psi(\xhat, \shat, \muhat)
&\le
e^{40 \lambda \gamma \rho^{-8}}
\left( (1 - \kappa) \Psi(x, s, \mu) + n \right)
\\
&\le
\left( 1 + \frac{\kappa}{4} \right)
\left( (1 - \kappa) \Psi(x, s, \mu) + n \right)
\\
&\le
\left( 1 - \frac{\kappa}{2} \right) \Psi(x, s, \mu) + 2 n,
\end{aligned}
\]
where the middle step uses
\[
40 \lambda \gamma \rho^{-8}
\le
\frac{\kappa}{8}
=
\frac{\alpha^{2} \lambda \gamma}{256 C d^{1/2}},
\qquad
\text{i.e.}
\qquad
\rho^{8}
\ge
10240 \, C d^{1/2} \alpha^{-2},
\]
true for $\rho \ge n$ and
$n$ larger than a constant.
Item~\ref{itm:perturb:2} with the feasibility bound of Lemma~\ref{lem:progress}
gives
\[
\left\| A^\top \xhat - b
\right\|_{(A^\top (T(\xhat) \Phi''(\xhat))^{-1} A)^{-1}}
\le
e^{5 \beta} \left( \frac{\epsilon \gamma}{2 \Cnorm}
+ (5 d)^{1/2} \rho^{-8} \gamma \right)
\le
\frac{\epsilon \gamma}{\Cnorm},
\]
Item~\ref{itm:perturb:3} gives dual feasibility, and centrality follows from
$\Psi(\xhat, \shat, \muhat) \le n^{2} + 2 n \le n^{10}$ as noted after
Definition~\ref{def:potential}.
\end{proof}

The path following lemma is Lemma 4.12 of
\cite{BLLSSSW21}.

\begin{lemma}
\label{lem:pathfollowing}
Start from an $\epsilon / C_{\mathrm{start}}$-centered triple
$(x, s = A z + c, \mu_0)$ with an estimate
$\tautil \approx_{\gamma / 100} \tau(x)$, and iterate \textsc{Step} until
$\mu \le \mu_f$.
This makes
\[
r
=
\left\lceil \theta^{-1} \log(\mu_0 / \mu_f) \right\rceil
=
\O\left( d^{1/2} \log(\mu_0 / \mu_f) \right)
\]
calls, and on the event that every \textsc{Solve} call along the way meets the
guarantee of Definition~\ref{def:operations}, with probability at least
$1 - r n^{-7}$ over the randomness of the Lewis estimates, the triple
$(x, s, \mu)$ at the start and end of every call is
$\epsilon$-centered and satisfies
$\Psi(x, s, \mu) \le n^{2}$.
\end{lemma}

\begin{proof}
Induction on the round, using Lemma~\ref{lem:computed_step}.
The inductive hypothesis at the start of round $k$ is threefold: the triple is
$\epsilon$-centered with
$\Psi \le n^{2}$, the estimate
$\tautil$ carried from round
$k - 1$ satisfies
$\tautil \approx_{\gamma / 100} \tau(x^{(k - 1)})$, and
$\|\Phi''(x^{(k - 1)})^{1/2} (x^{(k)} - x^{(k - 1)})\|_{\infty} \le 1 / 100$; at
$k = 1$ the last two hold with
$x^{(0)} \defeq x^{(1)}$ by the hypothesis on
$\tautil$.
Given them, Lemma~\ref{lem:solver_form} supplies Items~\ref{itm:lewis:5}
and~\ref{itm:lewis:6} to the call of Lemma~\ref{lem:lewis} in step 2, whose
output is $\tautil \approx_{\gamma / 100} \tau(x^{(k)})$, so
Definition~\ref{def:invariant} holds for the round and
Lemma~\ref{lem:computed_step} applies; Item~\ref{itm:step_bounds:2} of
Lemma~\ref{lem:step_bounds} for the exact step, plus the $\rho^{-8} \gamma$ by
which Lemma~\ref{lem:computed_step} lets the computed step differ from it,
bounds the movement to $x^{(k + 1)}$ by
$O(\gamma) \le 1 / 100$ in the local norm, which is the third hypothesis for
round $k + 1$, and the new
$\tautil$ is the second.
The base case is the initial potential: an
$\epsilon / C_{\mathrm{start}}$-centered point has
$|y_i| \le \epsilon / (10 C)$ for every
$i$, so
\[
\Psi(x, s, \mu_0)
\le
n \cosh\left( \frac{\lambda \epsilon}{10 C} \right)
\le
n \, e^{3 \log n / 10}
\le
n^{2}
\]
by $\lambda \epsilon \le 3 C \log n$ from Equation~\eqref{eq:parameters}.
For the inductive step, if $\Psi \le n^{2}$ at the start of a round then
Lemma~\ref{lem:computed_step} gives, at its end,
$\Psi \le (1 - \kappa / 2) n^{2} + 2 n \le n^{2}$, because
$\kappa = \alpha^{3} / (32 C^{3} d^{1/2}) \ge 4 / n$ for
$n$ larger than a constant depending on
$C$, using
$n \ge d$; the same lemma gives
$\epsilon$-centrality at the end of the round, and step 2 of the next round
restores the hypothesis on $\tautil$.
The count of calls is $\mu_r = (1 - \theta)^{r} \mu_0 \le \mu_f$.
The failure probability is a union bound over the $r$ rounds, each making one
call to Lemma~\ref{lem:lewis}, which fails with probability at most $n^{-7}$
once its solver meets Item~\ref{itm:lewis:6}; that every \textsc{Solve} call,
inside and outside the estimate, does so is the event conditioned on in the
statement, and its probability is the subject of
Theorem~\ref{thm:inv_maintain_stable}, not of this lemma.
\end{proof}

Hence the potential stays at most $n^{2}$ throughout and
$\mu$ decreases from
$\mu_0$ to
$\mu_f$ in
$\O(d^{1/2} \log(\mu_0 / \mu_f))$ steps; Section~\ref{sec:overall} turns this
into a bit complexity.

\subsection{Stability of the weights handed to the data structure}

With $R = I$ the iterates themselves move stably, so unlike \cite[Section
4.4]{LPTWY25:arxiv} and \cite[Lemmas 4.44 and
4.45]{BLLSSSW21} no nearby auxiliary sequence
is needed: the stability below is a statement about the actual iterates, which
are vectors of numbers with $40 \log \rho$ bits by construction.

The following lemma is the stability of the path weights.

\begin{lemma}
\label{lem:stability}
Let $(x^{(k)}, z^{(k)})$,
$1 \le k \le r + 1$, be the iterates of Lemma~\ref{lem:pathfollowing}, on the
event of that lemma, the $k$-th call being made from
$(x^{(k)}, z^{(k)})$, and write
$\tau^{(k)} \defeq \tau(x^{(k)})$,
$T^{(k)} \defeq \diag(\tau^{(k)})$,
$w^{(k)}_i \defeq (\tau^{(k)}_i \phi_i''(x^{(k)}_i))^{-1}$ for the exact path
weights with $W^{(k)} \defeq \diag(w^{(k)})$, and
$\what^{(k)}$ for the vector handed to \textsc{Update} in round
$k$.
Then for every $k \le r$:
\begin{enumerate}
\item \label{itm:stability:1} ($\ell_2$ stability of the iterate)
  $\| \Phi''(x^{(k)})^{1/2} (x^{(k + 1)} - x^{(k)}) \|_{\tau^{(k)} + \infty} \lesssim \gamma$;
\item \label{itm:stability:2} (stability of the barrier)
  $\| \Phi''(x^{(k)})^{1/2} ( \Phi''(x^{(k + 1)})^{-1/2} - \Phi''(x^{(k)})^{-1/2} ) \|_{\tau^{(k)} + \infty} \lesssim \gamma$;
\item \label{itm:stability:3} ($\tau_q$ stability)
  $\| (T^{(k)})^{-1} ( \tau^{(k + 1)} - \tau^{(k)} ) \|_{\tau^{(k)} + \infty} \lesssim \gamma$
  and
  $\| (W^{(k)})^{-1} ( w^{(k + 1)} - w^{(k)} ) \|_{\tau^{(k)} + \infty} \lesssim \gamma$;
\item \label{itm:stability:4} (the fed weights)
  $\what^{(k)} \approx_{\gamma / 50} w^{(k)}$.
\end{enumerate}
\end{lemma}

\begin{proof}
Fix $k$ and write
$x = x^{(k)}$,
$\xhat = x^{(k + 1)}$, and
$\xhat^{\star}$ for the exact step of Equation~\eqref{eq:exact_step} from
$x$.
By Lemma~\ref{lem:computed_step},
$\|\Phi''(x)^{1/2} (\xhat - \xhat^{\star})\|_{\infty} \le \rho^{-8} \gamma$, and
by Equation~\eqref{eq:norm_compare} this is at most
$(1 + \Cnorm (5 d)^{1/2}) \rho^{-8} \gamma \le \gamma$ in
$\|\cdot\|_{\tau^{(k)} + \infty}$; through Equation~\eqref{eq:barrier_local} and
Lemma~\ref{lem:lewis_approx} the same $O(\rho^{-8} \gamma)$ bound holds for the
induced changes of $\Phi''^{-1/2}$ and of
$\tau$ between
$\xhat$ and
$\xhat^{\star}$.
So it suffices to prove Items~\ref{itm:stability:1} to~\ref{itm:stability:3}
with $\xhat^{\star}$ in place of
$\xhat$, and there they are, respectively, Item~\ref{itm:step_bounds:1},
Item~\ref{itm:step_bounds:3}, and the consequence of Lemma~\ref{lem:tau_move};
for the weights, $w = \tau^{-1} c^{2}$ with
$c = \Phi''^{-1/2}$, and with
$u \defeq \tau^{(k+1)} / \tau^{(k)} - 1$ and
$v \defeq c^{(k+1)} / c^{(k)} - 1$, both of
$\ell_\infty$ norm
$O(\gamma) \le 1/10$ by Item~\ref{itm:step_bounds:2} and
Item~\ref{itm:tau_move:1},
\[
\left| \frac{w^{(k + 1)}_i}{w^{(k)}_i} - 1 \right|
=
\left| \frac{(1 + v_i)^{2}}{1 + u_i} - 1 \right|
\le
2 |u_i| + 3 |v_i|,
\]
so the weight bound is $2$ times the
$\tau$ bound plus
$3$ times Item~\ref{itm:stability:2}.
Item~\ref{itm:stability:4} is
$\what^{(k)} \approx_{\rho^{-18}} (\widetilde T \Phi''(x))^{-1}$ from the
truncation in step 3 together with $\tautil \approx_{\gamma / 100} \tau(x)$ from
step 2, and $\rho^{-18} \le \gamma / 100$.
\end{proof}

Item~\ref{itm:stability:3} in $\|\cdot\|_{\tau^{(k)} + \infty}$ contains both an
$\ell_\infty$ bound of
$O(\gamma)$ and a
$\tau$-norm bound of
$O(\gamma / \Cnorm)$, which are the shapes of conditions 2 and 3 of
Definition~\ref{def:tau_stable}, except that the definition weighs condition 3
by the leverage scores $\sigma^{(k)}$ of the reweighted matrix
$(W^{(k)})^{1/2} A$ rather than by the Lewis weights
$\tau^{(k)} = \tau(x^{(k)})$ of
$\Phi''(x^{(k)})^{-1/2} A$.
These are different vectors, but the reweighting is by $\tau^{(k)}$ itself, and
at $q = 1 - \alpha$ that makes them comparable.

The following lemma compares the leverage scores of the reweighted matrix with
the path weights.

\begin{lemma}
\label{lem:sigma_vs_tau}
Let $x$ be a point with
$\tau = \tau(x)$,
$T = \diag(\tau)$,
$W = T^{-1} \Phi''(x)^{-1}$ the exact path weights at
$x$, and
$\sigma \defeq \tau(W^{1/2} A)$ the leverage scores of the reweighted matrix.
Then
\[
v + \sigma \approx_{1/2} \tau,
\qquad\text{in particular}\qquad
\sigma_i \le e^{1/2} \tau_i
\quad\text{for every } i .
\]
\end{lemma}

\begin{proof}
By Definition~\ref{def:central_weights},
$\tau = v + \tau(T^{1/2 - 1/q} \Phi''(x)^{-1/2} A)$, and at
$q = 1 - \alpha$
\[
\frac12 - \frac1q
=
- \frac12 - \frac{\alpha}{1 - \alpha},
\qquad\text{so}\qquad
T^{1/2 - 1/q} \Phi''(x)^{-1/2}
=
T^{-\alpha / (1 - \alpha)} W^{1/2}
=
\left( D^{2} W \right)^{1/2},
\quad
D \defeq T^{-\alpha / (1 - \alpha)} .
\]
The computation of Lemma~\ref{lem:solver_form}, with $\tau$ in place of
$\tautil'$ and
$d / n \le \tau_i \le 5$, gives
$D^{2} = T^{-2 \alpha / (1 - \alpha)} \approx_{1/4} I$, hence
$D^{2} W \approx_{1/4} W$, and leverage scores move by at most twice the
movement of the weights, as recorded after Definition~\ref{def:tau_stable}:
$\tau((D^{2} W)^{1/2} A) \approx_{1/2} \tau(W^{1/2} A) = \sigma$.
The second claim is $\sigma_i \le e^{1/2} (\tau_i - v_i) \le e^{1/2} \tau_i$.
\end{proof}

Lemma~\ref{lem:ipm_is_stable} combines this with Lemma~\ref{lem:stability} to
verify Definition~\ref{def:tau_stable}, and Corollary~\ref{cor:update_counts}
and Theorem~\ref{thm:inv_maintain_stable} then bound the total work spent on
inverse maintenance over the whole path.
Condition 3 of Definition~\ref{def:tau_stable} is proved for the exact $w^{(k)}$
and not for the fed $\what^{(k)}$, whose deviation
$e^{\pm \gamma / 50}$ from
$w^{(k)}$ has leverage score norm up to
$d^{1/2} \gamma / 25$; this is why Section~\ref{sec:dynamic} states its
guarantee for sequences that are stable up to noise, and why the fed sequence
qualifies is the last step of Lemma~\ref{lem:ipm_is_stable}.

\section{Initialization and Overall Guarantees}
\label{sec:overall}

This section collects the three things still missing: that every matrix the
algorithm stores has magnitude $2^{O(B)}$, so that
$O(B)$ bits per entry suffice throughout; that the weights fed to
\textsc{Update} meet Definition~\ref{def:tau_stable}, so that
Theorem~\ref{thm:inv_maintain_stable} and Corollary~\ref{cor:adaptive} apply;
and an initial point from which $\O(d^{1/2})$ path steps suffice.

\subsection{Aspect ratio, magnitudes, and bit lengths}
\label{subsec:magnitudes}

Everything rests on the path aspect ratio, which is where the assumption that
the midpoint is feasible is used.

The following is Lemma 4.46 of
\cite{BLLSSSW21}.

\begin{importedresult}
\label{imp:aspect}
Suppose the midpoint $(l + u) / 2$ is feasible, and let
$W_0$ be the ratio of the largest to the smallest entry of
$\Phi''((l + u) / 2)^{1/2}$.
Then, over all points $x$ encountered by the path following of
Lemma~\ref{lem:pathfollowing}, the ratio $W'$ of the largest to the smallest
entry of $\Phi''(x)^{1/2}$ satisfies
\[
\log W'
=
\O\left(
\log (2 + W_0) + \log(1 / \mu_f) + \log (2 + \|c\|_{\infty})
\right).
\]
\end{importedresult}

The following lemma bounds the magnitudes along the path.

\begin{lemma}
\label{lem:magnitudes}
Run the algorithm on the augmented instance of Lemma~\ref{lem:initialization},
whose data has $O(B)$ bits and whose midpoint is feasible, with
$\mu_f = 2^{-O(B)}$.
Then at every point $x$ of the path, and for the maintained weights
$w_i = (\tau_i(x) \phi_i''(x_i))^{-1}$,
\[
\begin{aligned}
2^{-O(B)}
&\le
\phi_i''(x_i) \le 2^{O(B)},
\\
\frac{d}{n}
&\le
\tau_i(x) \le 1 + v_i \le 5,
\\
2^{-O(B)}
&\le
w_i \le 2^{O(B)},
\end{aligned}
\]
and consequently $A^\top W A$,
$A^\top \Phi''(x)^{-1} A$ and every normal matrix formed by \textsc{Step} has
magnitude at most $\rho = 2^{O(B)}$.
In particular $O(\log \rho) = O(B)$ bits per entry suffice in
Lemma~\ref{lem:inv_maintain_explicit} and Theorem~\ref{thm:inv_maintain_stable}.
\end{lemma}

\begin{proof}
Imported result~\ref{imp:aspect} bounds the ratio between the largest and
smallest $\phi_i''(x_i)^{1/2}$ along the path by
$2^{O(B)}$, since the augmented widths,
$\|c\|_{\infty}$ and
$1 / \mu_f$ are all
$2^{O(B)}$ by Lemma~\ref{lem:initialization}.
An absolute bound follows from centrality.
Let $m \defeq (l + u) / 2$, so
$A^\top m = b$, and let
$(x, s, \mu)$ be
$\epsilon$-centered with
$s = A z + c$, so that
$s_i = \mu \tau_i(x) ( - \phi_i'(x_i) + \nu_i \phi_i''(x_i)^{1/2} )$ with
$\|\nu\|_{\infty} \le \epsilon \le 1/80$ by Item~\ref{itm:centered:1}.
Pairing $s$ with
$x - m$ gives
\[
z^\top \left( A^\top x - b \right)
+ c^\top \left( x - m \right)
+ \mu \sum_i \tau_i(x) g_i
=
0 ,
\]
where
\[
g_i
\defeq
\left( \phi_i'(x_i)
- \nu_i \phi_i''(x_i)^{1/2} \right)
\left( x_i - m_i \right) ,
\]
since $z^\top A^\top (x - m) = z^\top (A^\top x - b)$.
Write $N \defeq A^\top (T(x) \Phi''(x))^{-1} A$.
The first term is at most
$\|z\|_{N} \|A^\top x - b\|_{N^{-1}} \le \|z\|_{N} \, \epsilon \gamma / \Cnorm$
by Item~\ref{itm:centered:3}, and
\[
\begin{aligned}
\|z\|_{N}^{2}
&=
\sum_i \frac{(s_i - c_i)^{2}}{\tau_i(x) \phi_i''(x_i)}
\\
&\le
2 \mu^{2} (1 + \epsilon)^{2} \sum_i \tau_i(x)
+ 2 \|c\|_{\infty}^{2} \, \frac{n}{d} \sum_i \frac{(u_i - l_i)^{2}}{8}
\\
&=
2^{O(B)} ,
\end{aligned}
\]
using $|s_i| \le \mu \tau_i (1 + \epsilon) \phi_i''(x_i)^{1/2}$ from
Lemma~\ref{lem:self_concordant}, $d / n \le \tau_i \le 5$,
$\phi_i''(t) \ge 8 / (u_i - l_i)^{2}$, and
$\mu \le \mu_0 = 2^{O(B)}$.
The second term is at most $\|c\|_{\infty} \sum_i (u_i - l_i) = 2^{O(B)}$.
Each $g_i \ge -1$:
$\phi_i'(x_i)$ has the sign of
$x_i - m_i$, and in the middle half of
$[l_i, u_i]$ one has
$\phi_i''(x_i)^{1/2} |x_i - m_i| \le 2^{1/2}$, so
$|\nu_i| \phi_i''(x_i)^{1/2} |x_i - m_i| \le 1$.
For a coordinate at distance $\delta_i \le (u_i - l_i) / 4$ from its nearer
bound, $|\phi_i'(x_i)| \ge 1 / \delta_i - 1 / (3 \delta_i) \ge 1 / (2 \delta_i)$
and $|\nu_i| \phi_i''(x_i)^{1/2} \le 1.1 / (80 \delta_i) \le 1 / (4 \delta_i)$,
so $g_i \ge |x_i - m_i| / (4 \delta_i) \ge (u_i - l_i) / (16 \delta_i)$.
Isolating such a coordinate, with $g_j \ge -1$ for the others,
\[
\begin{aligned}
\mu \frac{d}{n} \cdot \frac{u_i - l_i}{16 \delta_i}
&\le
\mu \tau_i(x) g_i
\\
&\le
\left| z^\top (A^\top x - b) \right|
+ \left| c^\top (x - m) \right|
+ 5 d \mu
\\
&=
2^{O(B)} ,
\end{aligned}
\]
so $1 / \delta_i \le 2^{O(B)} n / (\mu_f d (u_i - l_i)) = 2^{O(B)}$ and
$\phi_i''(x_i) = \Theta(\delta_i^{-2})$ is
$2^{O(B)}$, while the lower bound on
$\phi_i''$ is
$8 / (u_i - l_i)^{2} = 2^{-O(B)}$.
The bounds on $\tau$ are Definition~\ref{def:central_weights}, and the bound on
$w$ follows.
Finally $A^\top W A \preceq \max_i w_i \cdot A^\top A$ and
$A^\top W A \succeq \min_i w_i \cdot A^\top A$, and
$A^\top A$ has magnitude
$2^{O(B)}$ by assumption, so
$A^\top W A$ has magnitude
$2^{O(B)}$.
\end{proof}

\subsection{The maintained weights are leverage stable}
\label{subsec:ipm_stable}

The following lemma shows that the path weights are leverage stable.

\begin{lemma}
\label{lem:ipm_is_stable}
Let $w^{(k)}_i = (\tau_i(x^{(k)}) \phi_i''(x^{(k)}_i))^{-1}$ be the exact path
weights of Lemma~\ref{lem:stability}, over the
$r = \O(d^{1/2} \log(\mu_0 / \mu_f))$ rounds of Lemma~\ref{lem:pathfollowing}.
Then $w^{(1)}, \ldots, w^{(r + 1)}$ is
$\zeta$-leverage stable in the sense of Definition~\ref{def:tau_stable}, with
\[
\zeta = O(\gamma)
\qquad\text{and}\qquad
\rho = 2^{O(B)},
\]
and condition 3 of that definition holds with the smaller value
$O(\gamma / \Cnorm)$ in place of
$\zeta$.
Consequently the sequence $\what^{(1)}, \ldots, \what^{(r + 1)}$ actually handed
to \textsc{Update} in step 3 of Figure~\ref{fig:step} is $\zeta$-leverage stable
up to noise, in the sense recorded after Definition~\ref{def:tau_stable}.
\end{lemma}

\begin{proof}
Condition 1 is Lemma~\ref{lem:magnitudes}.
For conditions 2 and 3, Item~\ref{itm:stability:3} gives
\[
\left\| (W^{(k)})^{-1}
\left( w^{(k + 1)} - w^{(k)} \right)
\right\|_{\tau^{(k)} + \infty}
\lesssim \gamma,
\]
and by Definition~\ref{def:taunorm} this single bound contains both
$\|\cdot\|_{\infty} \lesssim \gamma$, which is condition 2 since relative and
logarithmic movement agree up to a factor of two once the movement is at most
$0.1$ coordinatewise, and
$\|\cdot\|_{\tau^{(k)}} \lesssim \gamma / \Cnorm$.
Condition 3 at round $k + 1$ weighs the same logarithmic movement by
$\sigma^{(k + 1)} = \tau((W^{(k + 1)})^{1/2} A)$.
By the fact recorded after Definition~\ref{def:tau_stable} and condition 2,
$\sigma^{(k + 1)} \approx_{2 \zeta} \sigma^{(k)}$, and by
Lemma~\ref{lem:sigma_vs_tau} at $x^{(k)}$,
$\sigma^{(k)}_i \le e^{1/2} \tau^{(k)}_i$, so
\[
\begin{aligned}
\sum_i \sigma^{(k + 1)}_i
\log^{2}\left( \frac{w^{(k + 1)}_i}{w^{(k)}_i} \right)
&\le
4 e^{1/2 + 2 \zeta}
\sum_i \tau^{(k)}_i
\left( \frac{w^{(k + 1)}_i - w^{(k)}_i}{w^{(k)}_i} \right)^{2}
\\
&\lesssim
\frac{\gamma^{2}}{\Cnorm^{2}},
\end{aligned}
\]
which is condition 3 with $O(\gamma / \Cnorm)$, and a fortiori with
$\zeta = O(\gamma)$.
Finally $\gamma < 0.1$ by Equation~\eqref{eq:parameters}.
For the fed sequence, Item~\ref{itm:stability:4} gives
$\what^{(k)} \approx_{\gamma / 50} w^{(k)}$, and taking the constant in
$\zeta = O(\gamma)$ to be at least
$1/5$ makes this the
$\approx_{\zeta / 10}$ that stability up to noise asks for.
\end{proof}

So the interior point method supplies the stability conditions with room to
spare: the $\ell_\infty$ condition holds with
$O(\gamma)$ where
$0.1$ is asked, and the leverage score condition, which is the one that drives
the update counts of Corollary~\ref{cor:update_counts}, is $\Cnorm$ times inside
even that.
The weights are computed from the solver's outputs, so it is the adaptive half
of Theorem~\ref{thm:inv_maintain_stable}, that is Corollary~\ref{cor:adaptive},
that is invoked; the number of rounds,
$\O(d^{1/2} \log(\mu_0/\mu_f)) = \O(d^{1/2} B) \le d \log^{O(1)} n \le n^{2}$ by
the standing assumption $B \le d^{1/2}$, is within the hypothesis
$r \le n^{2}$ of that corollary, and the number
$Q$ of solver calls, a polylogarithmic factor more, enters its failure
probability as $2 Q n^{-3}$.

\subsection{Initialization}
\label{subsec:init}

The final point is the projection of Lemma 4.11 of
\cite{BLLSSSW21}, which is proved there from
the matrix norm bounds of its Lemma 4.28 and the self-concordance inequalities
of Nesterov.
We state and prove a version adapted to the bounds of this paper.
Three things change relative to the source.
The projection is taken in the estimate $\tautil$ the algorithm holds, not in
the exact $\tau(x)$, since the exact weights are never available; the movement
bound is the explicit $2 \epsilon \gamma / \Cnorm$, from
Item~\ref{itm:centered:3} through the leverage scores of the reweighted matrix
and Lemma~\ref{lem:sigma_vs_tau}, in place of the source's $O(\epsilon)$, so
that no lower bound on $C$ is needed for the point to stay inside the box; and
the objective gap carries the explicit constant $2$, obtained from the one sided
bound $(-\phi_i'(x_i) + \nu_i \phi_i''(x_i)^{1/2})(x_i - x_i') \le 1 + \epsilon$
on the log barrier, in place of an unspecified absolute constant.

\begin{lemma}
\label{lem:finalpoint}
Let $(x, s, \mu)$ be
$\epsilon$-centered with
$\epsilon \le 1/80$, let
$\tautil \approx_{\epsilon} \tau(x)$, and let
$\overline{W} \defeq (\widetilde T \Phi''(x))^{-1}$.
Then
\[
x^{\mathrm{fin}}
\defeq
x + \overline{W} A
\left( A^\top \overline{W} A \right)^{-1}
\left( b - A^\top x \right)
\]
satisfies $A^\top x^{\mathrm{fin}} = b$,
\[
\left\| \Phi''(x)^{1/2} \left( x^{\mathrm{fin}} - x \right) \right\|_{\infty}
\le
\frac{2 \epsilon \gamma}{\Cnorm} ,
\]
so in particular $l < x^{\mathrm{fin}} < u$, and
\[
c^\top x^{\mathrm{fin}}
- \min\left\{ c^\top x' : A^\top x' = b,\ l \le x' \le u \right\}
\le
2 \mu \, \|\tau(x)\|_1 .
\]
\end{lemma}

\begin{proof}
Write $r \defeq b - A^\top x$,
$\overline{H} \defeq A^\top \overline{W} A$,
$W \defeq (T(x) \Phi''(x))^{-1}$ the exact path weights and
$N \defeq A^\top W A$ the matrix of Item~\ref{itm:centered:3}.
Since $\tautil \approx_{\epsilon} \tau(x)$ we have
$\overline{W} \approx_{\epsilon} W$ and
$\overline{H} \approx_{\epsilon} N$.
Feasibility is
$A^\top x^{\mathrm{fin}} = A^\top x + \overline{H} \, \overline{H}^{-1} r = b$.

\emph{Movement.}
Coordinate $i$ of the correction is
$\overline{w}_i a_i^\top \overline{H}^{-1} r$, and by Cauchy--Schwarz in the
$\overline{H}^{-1}$ inner product
\[
\begin{aligned}
\left| \overline{w}_i a_i^\top \overline{H}^{-1} r \right|
&\le
\overline{w}_i \, \|a_i\|_{\overline{H}^{-1}} \, \|r\|_{\overline{H}^{-1}}
\\
&=
\left( \overline{w}_i \, \overline{\sigma}_i \right)^{1/2}
\|r\|_{\overline{H}^{-1}} ,
\end{aligned}
\]
where
\[
\overline{\sigma}_i
\defeq
\overline{w}_i \, a_i^\top \overline{H}^{-1} a_i
=
\tau_i\left( \overline{W}^{1/2} A \right) .
\]
By the fact recorded after Definition~\ref{def:tau_stable},
$\overline{W} \approx_{\epsilon} W$ gives
$\overline{\sigma} \approx_{2 \epsilon} \tau(W^{1/2} A)$, and
Lemma~\ref{lem:sigma_vs_tau} gives
$\tau_i(W^{1/2} A) \le e^{1/2} \tau_i(x) \le e^{1/2 + \epsilon} \tautil_i$, so
$\overline{\sigma}_i \le e^{1/2 + 3 \epsilon} \tautil_i \le 2 \tautil_i$ and
$\overline{w}_i \, \overline{\sigma}_i \le 2 \overline{w}_i \tautil_i = 2 / \phi_i''(x_i)$.
Hence, with $\beta \defeq 2 \epsilon \gamma / \Cnorm$,
\[
\begin{aligned}
\phi_i''(x_i)^{1/2}
\left| x^{\mathrm{fin}}_i - x_i \right|
&\le
2^{1/2} \, \|r\|_{\overline{H}^{-1}}
\\
&\le
2^{1/2} e^{\epsilon / 2} \, \|r\|_{N^{-1}}
\\
&\le
2^{1/2} e^{\epsilon / 2} \, \frac{\epsilon \gamma}{\Cnorm}
\\
&\le
\frac{2 \epsilon \gamma}{\Cnorm}
\\
&=
\beta ,
\end{aligned}
\]
by $\overline{H} \approx_{\epsilon} N$ and Item~\ref{itm:centered:3}.
Since $\phi_i''(x_i)^{1/2} \ge 1 / \min(x_i - l_i, u_i - x_i)$ the correction
moves coordinate $i$ by at most
$\beta$ times its distance to the nearer bound, and
$\beta < 1$, so
$l < x^{\mathrm{fin}} < u$.

\emph{Objective.}
Let $x'$ be feasible,
$A^\top x' = b$ and
$l \le x' \le u$.
By Item~\ref{itm:centered:2}, $c = s - A z$, and
$A^\top x^{\mathrm{fin}} = A^\top x'$, so
$c^\top (x^{\mathrm{fin}} - x') = s^\top (x^{\mathrm{fin}} - x')$.
By Item~\ref{itm:centered:1},
$s_i = \mu \tau_i(x) ( - \phi_i'(x_i) + \nu_i \phi_i''(x_i)^{1/2} )$ with
$|\nu_i| \le \epsilon$.
For the log barrier, $- \phi_i'(t) = 1 / (t - l_i) - 1 / (u_i - t)$ and
$\phi_i''(t)^{1/2} \le 1 / (t - l_i) + 1 / (u_i - t)$, so for
$x'_i \in [l_i, u_i]$ with
$x'_i \le x_i$, using
$0 \le x_i - x'_i \le x_i - l_i$,
\[
\begin{aligned}
\left( - \phi_i'(x_i) + \nu_i \phi_i''(x_i)^{1/2} \right)
\left( x_i - x'_i \right)
&\le
\left( 1 + \epsilon \right) \frac{x_i - x'_i}{x_i - l_i}
-
\left( 1 - \epsilon \right) \frac{x_i - x'_i}{u_i - x_i}
\\
&\le
1 + \epsilon ,
\end{aligned}
\]
and symmetrically when $x'_i \ge x_i$.
For the remaining piece,
$|{-\phi_i'(x_i)} + \nu_i \phi_i''(x_i)^{1/2}| \le (1 + \epsilon) \phi_i''(x_i)^{1/2}$
by Lemma~\ref{lem:self_concordant}, and
$\phi_i''(x_i)^{1/2} |x^{\mathrm{fin}}_i - x_i| \le \beta$ by the movement
bound.
Adding,
\[
\begin{aligned}
c^\top \left( x^{\mathrm{fin}} - x' \right)
&=
\mu \sum_i \tau_i(x)
\left( - \phi_i'(x_i) + \nu_i \phi_i''(x_i)^{1/2} \right)
\left( (x_i - x'_i) + (x^{\mathrm{fin}}_i - x_i) \right)
\\
&\le
\mu \left( 1 + \epsilon \right) \left( 1 + \beta \right) \|\tau(x)\|_1
\\
&\le
2 \mu \, \|\tau(x)\|_1 ,
\end{aligned}
\]
since $(1 + \epsilon)(1 + \beta) \le (81/80)(1 + 1/40) < 2$.
Taking $x'$ optimal gives the claim.
\end{proof}

The statement is for the exact diagonal $\overline{W}$; the algorithm holds a
truncated copy of it, and the difference is handled by a perturbation bound in
the proof of Lemma~\ref{lem:initialization}, not by weakening the hypothesis.

The augmented instance is the following.

\begin{lemma}
\label{lem:initialization}
Given $A, b, c, l, u$ with
$B$ bits and
$A^\top A$ of magnitude at most
$2^{B}$, one can construct in
$\O(\nnz(A) B)$ bit operations an augmented instance
$(\widetilde{A}, b, \widetilde{c}, \widetilde{l}, \widetilde{u})$ with
\[
\widetilde{A}
=
\begin{bmatrix} A \\ E \end{bmatrix} \in \R^{(n + k) \times d},
\qquad
k \le d,
\qquad
\nnz(\widetilde{A}) \le \nnz(A) + d,
\]
all of whose data has $O(B)$ bits, such that
\begin{enumerate}
\item \label{itm:initialization:1} the midpoint
  $\widetilde{x} = (\widetilde{l} + \widetilde{u}) / 2$ is exactly
  feasible, $\widetilde{A}^\top \widetilde{x} = b$, and
  $\widetilde{A}$ has full column rank
  $d$;
\item \label{itm:initialization:2} with
  $s_{\mathrm{init}} = \widetilde{c}$ and
  $\mu_0 = 2^{O(B)}$, the triple
  $(\widetilde{x}, s_{\mathrm{init}}, \mu_0)$ is
  $\epsilon / C_{\mathrm{start}}$-centered in the sense of
  Definition~\ref{def:centered};
\item \label{itm:initialization:3} if the primal is
  feasible, running the path to
  $\mu_f = 2^{-O(B)}$ and rounding the output inwards to
  $[l, u]$ gives a point
  $x \in [l, u]$ with
  $\|A^\top x - b\|_2 \le 2^{-B}$ and
  $c^\top x \le 2^{-B} + \min\{ c^\top x' : A^\top x' = b,\ l \le x' \le u \}$;
\item \label{itm:initialization:4} $\log(\mu_0 / \mu_f) = O(B)$.
\end{enumerate}
\end{lemma}

\begin{proof}
Let $r_{\mathrm{mid}} = b - A^\top (l + u)/2$ be the residual of the midpoint,
$J = \{ j : (r_{\mathrm{mid}})_j \ne 0 \}$,
$k = |J| \le d$, and
$\Delta$ the least power of two above
$\max_i (u_i - l_i)$. For each
$j \in J$ add an artificial coordinate with box
$[0, 4\Delta]$, midpoint
$2 \Delta$, and row
$E_{j, :} = (r_{\mathrm{mid}})_j (2 \Delta)^{-1} e_j^\top$, so that
$E^\top (2 \Delta \mathbf{1}) = r_{\mathrm{mid}}$ and the augmented midpoint
$\widetilde{x} = ((l + u)/2, 2 \Delta \mathbf{1})$ satisfies
$\widetilde{A}^\top \widetilde{x} = A^\top (l + u)/2 + r_{\mathrm{mid}} = b$, giving
the feasibility of Item~\ref{itm:initialization:1}; adding rows
preserves column rank. Give the artificial coordinates objective
coefficient $\Lambda$, the least power of two above
$2^{B + 2} \|E\|_2 (1 + M)$ with
$M \defeq \sum_i |c_i| (u_i - l_i)$, where
$\|E\|_2 = \max_j |(r_{\mathrm{mid}})_j| / (2 \Delta)$ since the rows of
$E$ are supported on distinct columns. All the new data,
$\Delta$, the entries of $E$, and
$\Lambda$, is formed from the input by $O(1)$ arithmetic operations per
entry on numbers with $O(B)$ bits, using $\log n = O(B)$ as recorded
after Equation~\eqref{eq:parameters}, which is the cost claimed. The
augmented instance has
$n + k \le 2 n$ rows, so Section~\ref{sec:lewis} applies to it with
$n$ replaced by $n + k$, which changes no
$\O(\cdot)$.

For Item~\ref{itm:initialization:2}, at the midpoint every barrier
gradient vanishes, so with $s_{\mathrm{init}} = \widetilde{c}$ and
$z = 0$ Item~\ref{itm:centered:2} of Definition~\ref{def:centered}
holds, Item~\ref{itm:centered:3} holds with left hand side zero by
exact feasibility, and the centrality vector is
$y_i = \widetilde{c}_i / (\mu_0 \tau_i \phi_i''(\widetilde{x}_i)^{1/2})$
with $\phi_i''(\widetilde{x}_i)^{1/2} = 2^{3/2} / (\widetilde u_i - \widetilde l_i)$
and $\tau_i \ge d / (n + k) \ge d / (2 n)$, so
\[
|y_i|
\le
\frac{ 2 n \, \|\widetilde{c}\|_{\infty}
\max_i (\widetilde u_i - \widetilde l_i) }
{ 2^{3/2} d \, \mu_0 }
\le
\frac{\epsilon}{C_{\mathrm{start}}}
\]
for $\mu_0$ any power of two above
$2 C_{\mathrm{start}} \epsilon^{-1} (n / d) \max_i(\widetilde u_i - \widetilde l_i) \|\widetilde c\|_{\infty}$,
which is $2^{O(B)}$ because
$\|\widetilde c\|_{\infty} \le \Lambda = 2^{O(B)}$. We take
$\mu_f \defeq 2^{-B - 7} / d$, also
$2^{-O(B)}$, and Item~\ref{itm:initialization:4} follows.

For Item~\ref{itm:initialization:3}, let
$(\widetilde{x}, \widetilde{s}, \mu_f)$ be the $\epsilon$-centered
triple for the augmented instance at the end of the path, with the
estimate $\tautil$ of Definition~\ref{def:invariant}, and let
$\what$ be the weights of step 3 of Figure~\ref{fig:step} at this point,
$\what_i = \Trunc((\tautil_i \phi_i''(\widetilde x_i))^{-1}, 40 \log \rho)$,
so that \textsc{Solve} is against
$H \defeq \widetilde{A}^\top \widehat{W} \widetilde{A}$. Let
$\overline W \defeq (\widetilde T \Phi''(\widetilde x))^{-1}$ be the exact
diagonal of Lemma~\ref{lem:finalpoint} and
$\overline H \defeq \widetilde A^\top \overline W \widetilde A$; since
$\overline w_i \ge \rho^{-1}$ by Lemma~\ref{lem:magnitudes}, the
truncation error $\rho^{-20}$ per entry of Equation~\eqref{eq:step_roundoff}
gives
$\widehat W \approx_{\rho^{-18}} \overline W$ and
$H \approx_{\rho^{-18}} \overline H$. Lemma~\ref{lem:finalpoint}
gives the exact point
$\widetilde{x}^{\mathrm{fin}} = \widetilde x + \overline{W} \widetilde A \, \overline H^{-1} (b - \widetilde A^\top \widetilde x)$,
and the algorithm forms
$\widetilde{x}' \defeq \Trunc(\widetilde{x} + \widehat{W} \widetilde{A} \, \textsc{Solve}(\widetilde r), 40 \log \rho)$,
$\widetilde r \defeq \Trunc(b - \widetilde A^\top \widetilde x, 20 \log \rho)$,
in $\O((\nnz(A) + d^{2}) B^{2})$ bit operations, the
$B^{2}$ being the cost of the one \textsc{Solve}, the entries being
truncated having magnitude at most
$2 \max(\|\widetilde l\|_{\infty}, \|\widetilde u\|_{\infty}) \le \rho$
by Lemma~\ref{lem:finalpoint}. Writing
$r \defeq b - \widetilde A^\top \widetilde x$, each entry of
$\widetilde r$ is a sum of the second kind on stored numbers cut to
$20 \log \rho$ bits, so
$\|\widetilde r - r\|_2 \le (n + k)^{1/2} (\rho^{-36} + \rho^{-20}) \le \rho^{-19}$
and
$\|H^{-1} \widetilde r - H^{-1} r\|_2 \le \rho^{-18}$. Writing
$\delta \defeq \textsc{Solve}(\widetilde r) - H^{-1} \widetilde r$, the
guarantee of Equation~\eqref{eq:solve_guarantee} is
$\|\delta\|_{H} \le \rho^{-10} \|\widetilde r\|_{H^{-1}} \le 2 \rho^{-10}$, by
Item~\ref{itm:centered:3}, $\gamma \le 1$,
$\widetilde T \approx_{\gamma / 100} T(\widetilde x)$,
$H \approx_{\rho^{-18}} \overline H$ and
$\|\widetilde r - r\|_{H^{-1}} \le \rho^{-18}$, so
$\|\delta\|_2 \le 2 \rho^{-9}$. The truncation of the weights
contributes
\[
\begin{aligned}
&\left\| \widehat W \widetilde A H^{-1} r
- \overline W \widetilde A \, \overline H^{-1} r \right\|_2
\\
&\le
\left\| \left( \widehat W - \overline W \right) \widetilde A H^{-1} r \right\|_2
+ \left\| \overline W \widetilde A
\left( H^{-1} - \overline H^{-1} \right) r \right\|_2
\\
&\le
\rho^{-18} \rho \cdot \rho^{3} \|r\|_2
+ \rho^{3} \cdot 2 \rho^{-18} \rho \cdot \rho \|r\|_2
\\
&\le
\rho^{-11},
\end{aligned}
\]
using $\|\widehat W - \overline W\| \le \rho^{-18} \rho$,
$\|H^{-1} - \overline H^{-1}\| \le \|H^{-1}\| \, \|\overline H - H\| \, \|\overline H^{-1}\|$
with $\|\overline H - H\| \le 2 \rho^{-18} \rho$,
the magnitude bounds
$\|\widehat W\|, \|\overline W\|, \|\widetilde A\|, \|H^{\pm 1}\|, \|\overline H^{\pm 1}\| \le \rho$
of Lemma~\ref{lem:magnitudes}, and
$\|r\|_2 \le \|\overline H\|^{1/2} \|r\|_{\overline H^{-1}} \le \rho$.
Together with the solve error, the truncation of the right hand
side, and the final truncation, which by Equation~\eqref{eq:step_roundoff}
moves each entry by at most
$\rho^{-20}$,
\[
\begin{aligned}
\| \widetilde x' - \widetilde x^{\mathrm{fin}} \|_2
&\le
\| \widehat W \widetilde A \|_2 \left( \|\delta\|_2 + \rho^{-18} \right)
+ \rho^{-11}
+ (n + k)^{1/2} \rho^{-20}
\\
&\le
\rho^{-6},
\end{aligned}
\]
and
\[
\begin{aligned}
\| \widetilde A^\top \widetilde x' - b \|_2
&=
\| \widetilde A^\top ( \widetilde x' - \widetilde x^{\mathrm{fin}} ) \|_2
\\
&\le
\rho^{-5},
\end{aligned}
\]
using $\rho \ge n$ and, below,
$\rho \ge \|\widetilde c\|_2$, both part of the choice of $\rho$ recorded
after Equation~\eqref{eq:parameters}. Finally clip
$\widetilde x'$ into the augmented box coordinatewise; since
$\widetilde x^{\mathrm{fin}}$ is inside the box, clipping moves each
coordinate by at most $|\widetilde x'_i - \widetilde x^{\mathrm{fin}}_i|$
and the two displayed bounds persist with $\rho^{-5}$ replaced by
$\rho^{-4}$. Call the result
$(x, y) \in [l, u] \times [0, 4 \Delta]^{k}$, with
$x$ the returned point.

It remains to compare with the original instance. Let
$x^\star$ be optimal for the primal, so that
$(x^\star, 0)$ is feasible for the augmented instance and the augmented
optimum is at most
$c^\top x^\star$. Lemma~\ref{lem:finalpoint} gives
\[
\begin{aligned}
c^\top x + \Lambda \, \mathbf{1}^\top y
&\le
c^\top x^\star
+ 2 \mu_f \|\tau(\widetilde x)\|_1
+ \|\widetilde c\|_2 \rho^{-4}
\\
&\le
c^\top x^\star + 10 d \mu_f + \rho^{-3},
\end{aligned}
\]
using $\|\tau(\widetilde x)\|_1 \le 5 d$ from
Definition~\ref{def:central_weights}. Since both $x$ and
$x^\star$ lie in $[l, u]$,
$c^\top x - c^\top x^\star \ge -M$, so
$\Lambda \, \mathbf{1}^\top y \le M + 10 d \mu_f + \rho^{-3} \le M + 1$
by the choice of $\mu_f$, and therefore
\[
\begin{aligned}
\| A^\top x - b \|_2
&\le
\| E^\top y \|_2 + \rho^{-4}
\\
&\le
\|E\|_2 \, \mathbf{1}^\top y + \rho^{-4}
\\
&\le
\frac{ \|E\|_2 (1 + M) }{ \Lambda } + \rho^{-4}
\\
&\le
2^{-B - 2} + \rho^{-4}
\\
&\le
2^{-B},
\end{aligned}
\]
the first step because $A^\top x - b = -E^\top y + (\widetilde A^\top (x, y) - b)$,
and the last because $\rho \ge 2^{B}$. For the objective,
$\Lambda \, \mathbf{1}^\top y \ge 0$ gives
$c^\top x \le c^\top x^\star + 10 d \mu_f + \rho^{-3} \le c^\top x^\star + 2^{-B}$
by $\mu_f = 2^{-B - 7} / d$, which is the guarantee of
Item~\ref{itm:initialization:3}. The existence of $x^\star$ is the only
place the feasibility of the primal is used.
\end{proof}

\subsection{Proof of the main theorem}
\label{subsec:mainproof}

\begin{proof}[Proof of Theorem~\ref{thm:main}]
Run \textsc{Step} on the augmented instance of
Lemma~\ref{lem:initialization} from $(\widetilde x, z = 0, \mu_0)$, so
$s = s_{\mathrm{init}}$, until
$\mu \le \mu_f$, with the initial
$\tautil$ from the static estimate of Lemma~\ref{lem:lewis_static} at
accuracy $\gamma / 100$. By Lemma~\ref{lem:pathfollowing} and
Item~\ref{itm:initialization:4}, this is
\[
r
=
\O\left( d^{1/2} \log(\mu_0 / \mu_f) \right)
=
\O\left( d^{1/2} \, B \right)
\]
calls, at the end of which the potential is $n^{O(1)}$ and the point is
$\epsilon$-centered at
$\mu_f$; Item~\ref{itm:initialization:3} converts this, with one more
\textsc{Solve} call, into the claimed feasibility and objective
guarantees for the primal, the feasibility of that
program, which Item~\ref{itm:initialization:3} needs, being witnessed
by $\xini$. It
remains to account for the bit complexity of one call, all matrices
involved having magnitude $\rho = 2^{O(B)}$ by
Lemma~\ref{lem:magnitudes}, so that $\log \rho = O(B)$.
\begin{itemize}
\item The Lewis estimate of step 2 is one call of Lemma~\ref{lem:lewis}
  against \textsc{Solve}, with
  $T_{\mathcal{S}} = \O((\nnz(A) + d^{2}) \log^{2} \rho)$ by
  Item~\ref{itm:inv_maintain_stable:3} of
  Theorem~\ref{thm:inv_maintain_stable} and
  $\varepsilon = \gamma / 100$, so that
  $\varepsilon^{-2} \log^{2}(1 / \varepsilon) = O(\log^{7} n)$ by
  Equation~\eqref{eq:parameters}, at
\[
\O\left( \varepsilon^{-2}
\left( \nnz(A) \log \rho + T_{\mathcal{S}} \right)
\log^{2}(1 / \varepsilon) \right)
=
\O\left( \left( \nnz(A) + d^{2} \right) B^{2} \right)
\]
bit operations per round, hence
$\O(d^{1/2} (\nnz(A) + d^{2}) B^{3})$ over the
$r$ rounds. The one
\textsc{Initialize} is a one-off
$\O((\nnz(A) + \MM(d, d, d)) B)$ by
Item~\ref{itm:inv_maintain_stable:1}, and the static estimate of
Lemma~\ref{lem:lewis_static} that supplies the initial
$\tautil$ is a one-off
\[
\O\left( \left( \varepsilon^{-2} \left( \nnz(A) + d^{2} \right)
+ \MM(d, d, d) \right) \log \rho \log^{2}(1 / \varepsilon) \right)
=
\O\left( \left( \nnz(A) + \MM(d, d, d) \right) B \right),
\]
as $d^{2} \le \MM(d, d, d)$.
\item The weights $\what$ of step 3 are
  $\zeta$-leverage stable up to noise with
  $\zeta = O(\gamma)$ by Lemma~\ref{lem:ipm_is_stable}, so by
  Item~\ref{itm:inv_maintain_stable:2}, applied through
  Corollary~\ref{cor:adaptive} because the weights are computed from the
  solver's own outputs, all
  \textsc{Update} calls over the $r$ rounds cost
\[
\O\left( r \left( \nnz(A) + d^{2} \right) \log^{2} \rho \right)
=
\O\left( d^{1/2} \left( \nnz(A) + d^{2} \right) B^{3} \right)
\]
in total.
\item The two \textsc{Solve} calls of step 5 cost
  $\O((\nnz(A) + d^{2}) \log^{2} \rho) = \O((\nnz(A) + d^{2}) B^{2})$
  per round by Item~\ref{itm:inv_maintain_stable:3}, including the
  $\O(\nnz(A) \log \rho)$ overhead of the noisy solver of Imported
  result~\ref{imp:noisysolver}; \textsc{Direction} in step 4 costs
  $\O(n \log \rho)$ by Lemma~\ref{lem:direction}; and the explicit
  updates of steps 1, 3 and 6 apply
  $A$ and
  $A^\top$ a constant number of times and touch each coordinate
  $O(1)$ times, on numbers with
  $40 \log \rho$ bits, so they cost
  $\O(\nnz(A) B)$ per round. Over the
  $r$ rounds this is
  $\O(d^{1/2} (\nnz(A) + d^{2}) B^{3})$.
\end{itemize}
Summing, and using
$\MM(d, d, d) = \O(d^{1 + \sqrt{2}}) \le \O(d^{1/2} \cdot d^{2})$ from
Equation~\eqref{eq:mm_assumption} to absorb the one-off terms, the
total is
\[
\O\left( d^{1/2} \left( \nnz(A) + d^{2} \right) B^{3} \right)
\]
bit operations, the $\O$ hiding factors polylogarithmic in
$n$ only: one factor
$B$ is the bit length of every number handled, one the iteration count
$O(\log \rho)$ of a \textsc{Solve}, and one the number
$\O(d^{1/2} B)$ of rounds. The failure probability is
$n^{-7}$ from the static estimate, $r n^{-7}$ from
Lemma~\ref{lem:pathfollowing}, on whose event
Lemma~\ref{lem:stability} is deterministic, and
$n^{-5} + 2 Q n^{-3}$ from Theorem~\ref{thm:inv_maintain_stable}, the
number $Q$ of \textsc{Solve} calls being
$\O(r) = \O(d^{1/2} B)$: a constant number per round in
\textsc{Step}, $O(\log n)$ per round inside the structure, and
$\O(\gamma^{-2}) = \O(1)$ per round in the Lewis estimate. The total
failure probability is therefore $\O(d^{1/2} B) \, n^{-3} \le \O(d n^{-3}) \le \O(n^{-2})$,
the hidden factor being polylogarithmic in $n$, which is
the ``high probability'' of the theorem; by the standing assumption
$B \le d^{1/2}$ the number of rounds is
$r = \O(d^{1/2} B) \le d \log^{O(1)} n \le n^{2}$ for
$d$ larger than a constant, which is the hypothesis
$r \le n^{2}$ of Theorem~\ref{thm:inv_maintain_stable} and
Corollary~\ref{cor:adaptive}. No
term of
the form $\MM(n, d, d)$ survives: the only dense operations are on
$d \times d$ matrices and on the
$\O(d)$-row sampled matrix, and the full
$A$ is only ever touched by matrix vector products and by the single
static Lewis estimate of Lemma~\ref{lem:lewis_static}.
\end{proof}


\bibliographystyle{alpha}
\bibliography{ref}

\appendix
\section{Proof of Lemma~\ref{lem:roundoff}}
\label{app:roundoff}

\begin{proof}
Write $L_v$ for the constant
$1$ if operation
$v$ is
$+$ or
$-$,
$2 M$ if it is
$\times$,
$2 M^{3}$ if it is
$\div$, and
$M$ if it is a square root; in every case
$L_v \le 2 M^{3}$. We show by induction along the formula that the error
$e_v$ of the computed value at operation
$v$ satisfies
\[
e_v
\le
2^{-\ell} \sum_{u \preceq v} \prod_{u \prec t \preceq v} L_t ,
\]
where the sum runs over the operations $u$ in the subformula of
$v$, including
$v$, and the product over the operations strictly after
$u$ on the unique path to
$v$. Since at most
$q$ factors of each product exceed
$1$ and each is at most
$2 M^{3}$, the right hand side is at most
$m ( 2 M^{3} )^{q} 2^{-\ell}$, which is the claim for the output and is
at most
$1 / (2 M)$ for every intermediate result by the hypothesis.

Let $v$ have inputs with exact values
$x, y$ and computed values
$\widetilde x, \widetilde y$, with errors
$e_x, e_y \le 1 / (2 M)$ by induction. For
$+$ and
$-$ the exact operation on the computed inputs is off by at most
$e_x + e_y$. For
$\times$,
\[
\begin{aligned}
\left| \widetilde x \widetilde y - x y \right|
&\le
\left| \widetilde x \right| \left| \widetilde y - y \right|
+ \left| y \right| \left| \widetilde x - x \right|
\\
&\le
\left( M + \frac{1}{2 M} \right) e_y + M e_x
\\
&\le
2 M \left( e_x + e_y \right) .
\end{aligned}
\]
For $\div$ the divisor satisfies
$|y| \ge 1 / M$ and hence
$|\widetilde y| \ge 1 / (2 M)$, so
\[
\begin{aligned}
\left| \frac{\widetilde x}{\widetilde y} - \frac{x}{y} \right|
&=
\frac{ \left| \widetilde x y - x \widetilde y \right| }
     { \left| y \widetilde y \right| }
\\
&\le
2 M^{2} \left( \left| y \right| e_x + \left| x \right| e_y \right)
\\
&\le
2 M^{3} \left( e_x + e_y \right) .
\end{aligned}
\]
For the square root the radicand satisfies
$x \ge 1 / M$ and
$\widetilde x \ge 1 / (2 M) > 0$, so
\[
\left| \widetilde x^{1/2} - x^{1/2} \right|
=
\frac{ \left| \widetilde x - x \right| }
     { \widetilde x^{1/2} + x^{1/2} }
\le
M^{1/2} e_x
\le
M e_x .
\]
In every case the exact operation on the computed inputs is within
$L_v ( e_x + e_y )$ of the exact result; its absolute value is at most
$M + 1 / (2 M) < 2^{\ell}$, so it has at most
$\ell$ bits before the decimal point and the truncation adds at most
$2^{-\ell}$, giving
$e_v \le L_v ( e_x + e_y ) + 2^{-\ell}$, which is the inductive
bound.
\end{proof}

\end{document}
